\documentclass[11pt, a4paper]{report}

% --- Paquetes Fundamentales ---
\usepackage[T1]{fontenc} % CRÍTICO: Para usar < y > en el texto sin errores
\usepackage[utf8]{inputenc}
\usepackage[spanish]{babel}
\usepackage{titlesec}
\usepackage{geometry}
\usepackage[hidelinks]{hyperref}
\usepackage{amsthm}
\usepackage{booktabs}

% --- Estética y Tablas ---
\usepackage{eurosym}
\usepackage{multirow}
\usepackage[table,xcdraw]{xcolor}
\usepackage{float}
\usepackage{graphicx}

% --- Matemáticas y Símbolos ---
\usepackage{amsmath}
\usepackage{amsfonts}
\usepackage{amssymb}

% --- Diagramas (TikZ) ---
\usepackage{tikz}
\usetikzlibrary{shapes.geometric, arrows.meta, positioning, fit, calc, babel}

% --- Algoritmos y Código ---
\usepackage{algorithm}
\usepackage{algorithmic}
\usepackage{listings}

% --- Configuración de Listings (Código Python) ---
\lstset{
    language=Python,
    basicstyle=\ttfamily\footnotesize,
    keywordstyle=\color{blue}\bfseries,
    stringstyle=\color{olive},
    commentstyle=\color{gray}\itshape,
    breaklines=true,
    breakatwhitespace=true,
    showstringspaces=false,
    numbers=left,
    numberstyle=\tiny\color{gray},
    frame=single,
    rulecolor=\color{black!20},
    captionpos=b,
    xleftmargin=15pt
}

\theoremstyle{definition}
\newtheorem{ejemplo}{Ejemplo}[chapter]
\renewcommand{\lstlistingname}{Algoritmo} 
\begin{document}

\chapter*{Resumen}

La computación cuántica representa una amenaza inminente para los sistemas de clave pública actuales (RSA, ECC) debido a su capacidad teórica para quebrar los problemas matemáticos mediante algoritmos como el de Shor. Ante la aproximación del denominado \textit{Q-Day}, este Trabajo de Fin de Grado aborda la implementación práctica y funcional del estándar de firma digital poscuántica ML-DSA (FIPS 204) establecido por el NIST en 2024. 

El proyecto abandona la visión puramente teórica para desarrollar un sistema distribuido de extremo a extremo. Por un lado, se ha construido un motor criptográfico algorítmico en Python que ejecuta con fidelidad las operaciones sobre retículos y anillos polinómicos. Por otro lado, se ha diseñado una aplicación web de página única SPA en React que abstrae esta complejidad matemática, ofreciendo un entorno seguro y usable para la gestión de identidades y la firma de documentos PDF. La manipulación de los archivos se resuelve mediante la técnica de inyección binaria \textit{EOF Trick}, permitiendo incrustar firmas poscuánticas de forma nativa sin alterar la validez del archivo original.

Los resultados empíricos demuestran que la adopción temprana de la seguridad poscuántica es plenamente viable sobre arquitecturas informáticas clásicas, logrando un equilibrio óptimo entre el rigor algebraico, la mitigación de vectores de ataque y la usabilidad para el usuario final.

\vspace{0.5cm}
\noindent \textbf{Palabras clave:} Criptografía poscuántica, ML-DSA, FIPS 204, Firma digital, Sistemas distribuidos, Retículos, Seguridad web.

\chapter*{Abstract}

Quantum computing represents an imminent threat to current public-key systems (RSA, ECC) due to its theoretical capability to break mathematical problems using algorithms such as Shor's. In anticipation of the so-called \textit{Q-Day}, this Bachelor's Thesis addresses the practical and functional implementation of the post-quantum digital signature standard ML-DSA (FIPS 204) established by NIST.

The project shifts from a purely theoretical approach to develop an end-to-end distributed system. On the one hand, an algorithmic cryptographic engine has been built in Python that faithfully executes operations over lattices and polynomial rings. On the other hand, a Single Page Application (SPA) has been designed in React that abstracts this mathematical complexity, providing a secure and usable environment for identity management and PDF document signing. File manipulation is resolved through the binary injection technique known as the \textit{EOF Trick}, allowing post-quantum signatures to be natively embedded without altering the validity of the original file.

Empirical results demonstrate that the early adoption of post-quantum security is fully viable on classical computing architectures, achieving an optimal balance between algebraic rigor, mitigation of attack vectors, and end-user usability.

\vspace{0.5cm}
\noindent \textbf{Keywords:} Post-quantum cryptography, ML-DSA, FIPS 204, Digital signature, Distributed systems, Lattices, Web security.

\chapter*{Agradecimientos}
\addcontentsline{toc}{chapter}{Agradecimientos}
Quisiera expresar mi más sincero agradecimiento a todas las personas que me han acompañado desde el inicio de mi etapa universitaria y que han confiado plenamente en mi capacidad para culminar este camino. 

De manera muy especial, quiero dedicar este logro a mi familia. Su incondicional apoyo emocional y su esfuerzo económico han sido los pilares fundamentales para alcanzar esta meta. Una gran parte de este mérito les pertenece, ya que sin su respaldo nada de esto habría sido posible.

A nivel académico, deseo mostrar mi profunda gratitud a mi tutor, Justo Peralta López. Su labor docente en la asignatura de Teoría de Códigos y Criptografía fue el detonante que me introdujo en este apasionante campo, motivándome no solo a desarrollar este Trabajo de Fin de Grado, sino a orientar mi futuro profesional hacia el sector de la ciberseguridad.

\newpage

\tableofcontents
\listoffigures 
\listoftables  
\newpage

\newpage
\chapter*{Lista de Acrónimos y Abreviaturas}
\addcontentsline{toc}{chapter}{Lista de Acrónimos y Abreviaturas}

\begin{description}
    \item[AES] \textit{Advanced Encryption Standard} (Estándar de Cifrado Avanzado)
    \item[API] \textit{Application Programming Interface} (Interfaz de Programación de Aplicaciones)
    \item[ASGI] \textit{Asynchronous Server Gateway Interface} (Interfaz de Pasarela de Servidor Asíncrono)
    \item[BaaS] \textit{Backend-as-a-Service} (Backend como Servicio)
    \item[CA] \textit{Certificate Authority} (Autoridad de Certificación)
    \item[CMS] \textit{Cryptographic Message Syntax} (Sintaxis de Mensaje Criptográfico)
    \item[CORS] \textit{Cross-Origin Resource Sharing} (Intercambio de Recursos de Origen Cruzado)
    \item[CPU] \textit{Central Processing Unit} (Unidad Central de Procesamiento)
    \item[CRQC] \textit{Cryptographically Relevant Quantum Computer} (Ordenador Cuántico Criptográficamente Relevante)
    \item[CRYSTALS] \textit{Cryptographic Suite for Algebraic Lattices} (Suite Criptográfica para Retículos Algebraicos)
    \item[CSS] \textit{Cascading Style Sheets} (Hojas de Estilo en Cascada)
    \item[CVP] \textit{Closest Vector Problem} (Problema del Vector más Cercano)
    \item[CWE] \textit{Common Weakness Enumeration} (Enumeración de Debilidades Comunes)
    \item[DFT] \textit{Discrete Fourier Transform} (Transformada Discreta de Fourier)
    \item[DOM] \textit{Document Object Model} (Modelo de Objetos del Documento)
    \item[DRBG] \textit{Deterministic Random Bit Generator} (Generador Determinista de Bits Aleatorios)
    \item[DSA] \textit{Digital Signature Algorithm} (Algoritmo de Firma Digital)
    \item[ECC] \textit{Elliptic Curve Cryptography} (Criptografía de Curva Elíptica)
    \item[ECDLP] \textit{Elliptic Curve Discrete Logarithm Problem} (Problema del Logaritmo Discreto en Curvas Elípticas)
    \item[ECDSA] \textit{Elliptic Curve Digital Signature Algorithm} (Algoritmo de Firma Digital de Curva Elíptica)
    \item[ECTS] \textit{European Credit Transfer and Accumulation System} (Sistema Europeo de Transferencia y Acumulación de Créditos)
    \item[EOF] \textit{End of File} (Fin de Archivo)
    \item[FAQ] \textit{Frequently Asked Questions} (Preguntas Frecuentes)
    \item[FFT] \textit{Fast Fourier Transform} (Transformada Rápida de Fourier)
    \item[FIPS] \textit{Federal Information Processing Standards} (Estándares Federales de Procesamiento de Información)
    \item[FK] \textit{Foreign Key} (Clave Foránea)
    \item[GNFS] \textit{General Number Field Sieve} (Criba General del Cuerpo de Números)
    \item[HTML] \textit{HyperText Markup Language} (Lenguaje de Marcado de Hipertexto)
    \item[HTTP] \textit{Hypertext Transfer Protocol} (Protocolo de Transferencia de Hipertexto)
    \item[I+D] Investigación y Desarrollo
    \item[ID] \textit{Identifier} (Identificador)
    \item[IFP] \textit{Integer Factorization Problem} (Problema de la Factorización de Enteros)
    \item[$\text{NTT}^{-1}$] \textit{Inverse Number Theoretic Transform} (Transformada Inversa de Teoría de Números)
    \item[IP] \textit{Internet Protocol} (Protocolo de Internet)
    \item[ISO] \textit{International Organization for Standardization} (Organización Internacional de Normalización)
    \item[IVA] Impuesto sobre el Valor Añadido
    \item[JSON] \textit{JavaScript Object Notation} (Notación de Objetos de JavaScript)
    \item[JWT] \textit{JSON Web Token} (Token Web JSON)
    \item[LWE] \textit{Learning With Errors} (Aprendizaje con Errores)
    \item[ML-DSA] \textit{Module-Lattice-Based Digital Signature Algorithm} (Algoritmo de Firma Digital Basado en Retículos sobre Módulos)
    \item[MLWE] \textit{Module Learning With Errors} (Aprendizaje con Errores sobre Módulos)
    \item[MSIS] \textit{Module Short Integer Solution} (Solución de Enteros Cortos sobre Módulos)
    \item[NIST] \textit{National Institute of Standards and Technology} (Instituto Nacional de Estándares y Tecnología de EE. UU.)
    \item[NTT] \textit{Number Theoretic Transform} (Transformada de Teoría de Números)
    \item[OID] \textit{Object Identifier} (Identificador de Objeto)
    \item[PDF] \textit{Portable Document Format} (Formato de Documento Portátil)
    \item[PEM] Presupuesto de Ejecución Material
    \item[PH] \textit{Pre-Hash} (Pre-Resumen Criptográfico)
    \item[PK] \textit{Primary Key} (Clave Primaria)
    \item[PKI] \textit{Public Key Infrastructure} (Infraestructura de Clave Pública)
    \item[PQC] \textit{Post-Quantum Cryptography} (Criptografía Post-Cuántica)
    \item[PRNG] \textit{Pseudo-Random Number Generator} (Generador de Números Pseudoaleatorios)
    \item[QFT] \textit{Quantum Fourier Transform} (Transformada Cuántica de Fourier)
    \item[QR] \textit{Quick Response} (Respuesta Rápida)
    \item[REST] \textit{Representational State Transfer} (Transferencia de Estado Representacional)
    \item[RFC] \textit{Request for Comments} (Petición de Comentarios)
    \item[RLS] \textit{Row Level Security} (Seguridad a Nivel de Fila)
    \item[RSA] \textit{Rivest, Shamir, Adleman} (Algoritmo criptográfico asimétrico)
    \item[SaaS] \textit{Software as a Service} (Software como Servicio)
    \item[SDK] \textit{Software Development Kit} (Kit de Desarrollo de Software)
    \item[SHA] \textit{Secure Hash Algorithm} (Algoritmo de Hash Seguro)
    \item[SHAKE] \textit{Secure Hash Algorithm Keccak} (Algoritmo de Hash Seguro Keccak)
    \item[SP] \textit{Special Publication} (Publicación Especial)
    \item[SPA] \textit{Single Page Application} (Aplicación de Página Única)
    \item[SVP] \textit{Shortest Vector Problem} (Problema del Vector más Corto)
    \item[TDD] \textit{Test-Driven Development} (Desarrollo Guiado por Pruebas)
    \item[TFG] Trabajo de Fin de Grado
    \item[TLS/SSL] \textit{Transport Layer Security / Secure Sockets Layer} (Seguridad de la Capa de Transporte / Capa de Sockets Seguros)
    \item[UI/UX] \textit{User Interface / User Experience} (Interfaz de Usuario / Experiencia de Usuario)
    \item[URL] \textit{Uniform Resource Locator} (Localizador de Recursos Uniforme)
    \item[UUID] \textit{Universally Unique Identifier} (Identificador Único Universal)
    \item[VPS] \textit{Virtual Private Server} (Servidor Privado Virtual)
    \item[XOF] \textit{Extendable-Output Function} (Función de Salida Extensible)
    \item[XOR] \textit{Exclusive OR} (O Exclusivo)
    \item[XSS] \textit{Cross-Site Scripting} (Secuencias de Comandos en Sitios Cruzados)
\end{description}



\newpage

\chapter{Introducción}

\section{Motivación del Proyecto}

La confianza digital sobre la que se sustenta la sociedad moderna, desde transacciones bancarias hasta acuerdos gubernamentales, depende de algoritmos criptográficos asimétricos como RSA y la Criptografía de Curva Elíptica ECC. Sin embargo, el avance acelerado de la computación cuántica amenaza con desmantelar esta infraestructura de seguridad. El algoritmo propuesto por Peter Shor ha demostrado que un ordenador cuántico con suficientes qubits estables podrá resolver en tiempo polinómico los problemas matemáticos que protegen estos estándares, dejando obsoleta la criptografía actual.

Lejos de ser un problema lejano, esta vulnerabilidad representa un riesgo inmediato debido a la táctica cibernética conocida como \textit{Harvest Now, Decrypt Later}, que significa capturar ahora, descifrar después. Ciberdelincuentes ya están almacenando tráfico cifrado masivo a la espera de disponer del hardware cuántico necesario para revelarlo en el futuro. Ante esta inminente ruptura de la seguridad global, el Instituto Nacional de Estándares y Tecnología NIST ha publicado recientemente el estándar FIPS 204 \cite{nist_fips204}, mostrando el esquema ML-DSA, el cual está basado en CRYSTALS-Dilithium, y estableciéndolo como el nuevo estándar internacional para firmas digitales post-cuánticas.

La motivación fundamental de este proyecto nace de la brecha existente entre la publicación de estos estándares y su aplicación en la ingeniería de software real. La industria no solo necesita teoría criptográfica basada en retículos, sino personas que estén capacitadas para implementar estos algoritmos a bajo nivel y conectarlos de manera segura a las arquitecturas actuales. 

Por lo tanto, este Trabajo de Fin de Grado surge de la necesidad técnica de abordar la transición hacia la criptografía \textit{Quantum-Safe}. El propósito es demostrar que es posible construir una plataforma integral que implemente matemáticamente el estándar FIPS 204.

\section{La Amenaza Cuántica y el Algoritmo de Shor}

La computación cuántica representa un cambio respecto a la arquitectura computacional clásica de Von Neumann. Mientras que la computación tradicional está basada en su procesamiento en el bit, apoyado en estados binarios (0 o 1), la computación cuántica opera sobre bits cuánticos o \textit{qubits}. Gracias a los principios de la mecánica cuántica, específicamente el principio de superposición, un qubit puede representar una combinación lineal de ambos estados simultáneamente.

Matemáticamente, el estado de un qubit se define como:

$$|\psi\rangle = \alpha|0\rangle + \beta|1\rangle$$

donde $\alpha$ y $\beta$ son amplitudes de probabilidad complejas que satisfacen la condición de normalización $|\alpha|^2 + |\beta|^2 = 1$. Esta propiedad permite que un registro de $n$ qubits pueda representar y procesar $2^n$ estados de forma concurrente, permitiendo a los sistemas cuánticos un paralelismo masivo inalcanzable para la computación clásica.

La amenaza que representa la computación cuántica no reside en un simple aumento de la velocidad de procesamiento, sino en un cambio de paradigma fundamentado en dos principios de la mecánica cuántica: la superposición y la interferencia \cite{nielsen2010quantum}. 

Mientras que la computación clásica evalúa los estados de forma determinista y secuencial, un ordenador cuántico aprovecha la superposición para procesar una gran cantidad de estados posibles de forma simultánea. Sin embargo, el verdadero motor computacional es la interferencia. Los algoritmos cuánticos manipulan las amplitudes de probabilidad complejas de estos estados mediante puertas lógicas. El diseño de estos algoritmos aplica transformaciones para rotar las fases de las amplitudes de probabilidad complejas. De este modo, fuerza la colisión de fases de signo opuesto para anular los estados incorrectos (interferencia destructiva), mientras que alinea las fases del estado objetivo para maximizar su magnitud final (interferencia constructiva).

\begin{ejemplo}
\label{ex:interferencia_cuantica}
Supongamos un registro cuántico de 2 \textit{qubits} que representa 4 posibles claves de descifrado: $|00\rangle, |01\rangle, |10\rangle$ y $|11\rangle$. Solo el estado $|10\rangle$ es la clave correcta. A diferencia de un ataque clásico que probaría cada una secuencialmente, el algoritmo cuántico procede en bloque:

\begin{enumerate}
    \item \textbf{Superposición Uniforme:} Se inicializa el sistema asignando a cada estado la misma amplitud de probabilidad ($1/2$). La probabilidad real de medir cualquiera de ellos es su módulo al cuadrado ($|1/2|^2 = 0.25$, un $25\%$).
    
    \item \textbf{Marcaje de Fase:} El algoritmo aplica una matriz que evalúa la condición de éxito e invierte únicamente el signo (la fase) del estado correcto. Las amplitudes pasan a ser:
    $$ |00\rangle \to \frac{1}{2}, \quad |01\rangle \to \frac{1}{2}, \quad |10\rangle \to -\frac{1}{2}, \quad |11\rangle \to \frac{1}{2} $$
    En este punto, si midiéramos, la probabilidad seguiría siendo del $25\%$ para todos, ya que $|-1/2|^2 = 0.25$.
    
    \item \textbf{Interferencia Destructiva y Constructiva:} Se aplica el operador, el cual realiza una inversión sobre la media de las amplitudes. La media aritmética de las amplitudes actuales es:
    $$ \text{Media} = \frac{1}{4} \left( \frac{1}{2} + \frac{1}{2} - \frac{1}{2} + \frac{1}{2} \right) = \frac{1}{4} $$
    La fórmula de difusión para cada amplitud es: $\text{Nueva Amplitud} = 2 \cdot \text{Media} - \text{Amplitud Original}$. Al aplicar esta transformación de forma simultánea, ocurre el fenómeno de interferencia:
    \begin{itemize}
        \item Para los estados incorrectos ($|00\rangle, |01\rangle, |11\rangle$): $2 \cdot \left(\frac{1}{4}\right) - \frac{1}{2} = 0$. Sus amplitudes se anulan completamente (interferencia destructiva).
        \item Para el estado correcto ($|10\rangle$): $2 \cdot \left(\frac{1}{4}\right) - \left(-\frac{1}{2}\right) = \frac{1}{2} + \frac{1}{2} = 1$. Su amplitud se suma (interferencia constructiva).
    \end{itemize}
\end{enumerate}

Tras solo un par de operaciones matriciales, las probabilidades finales de observación pasan a ser del $0\%$ para todas las claves erróneas y del $100\%$ ($|1|^2$) para la clave correcta $|10\rangle$. El sistema ha descartado las respuestas incorrectas sin haberlas evaluado una por una.
\end{ejemplo}

Esta capacidad de procesamiento simultáneo y selectivo es exactamente el vector de ataque que compromete los sistemas criptográficos modernos. La amenaza teórica llegó en 1994, cuando Peter Shor propuso un algoritmo cuántico capaz de resolver en tiempo polinómico el problema de la factorización de enteros y el cálculo de logaritmos discretos \cite{shor1994}. Esto implica matemáticamente que, ante la existencia de un ordenador cuántico criptográficamente relevante (CRQC), los estándares actuales de clave pública como RSA y ECC quedarían completamente rotos.

\subsection{Superposición e Interferencia Cuántica}

Aunque el principio de superposición permite representar simultáneamente una cantidad exponencial de estados, existe una limitación física fundamental: el colapso de la función de onda. Según la mecánica cuántica, al realizar una medición sobre un sistema superpuesto, este colapsa irreversiblemente hacia un único estado clásico. La probabilidad de medir un estado específico está dictada por la Regla de Born, siendo proporcional al cuadrado del módulo de su amplitud de probabilidad. En consecuencia, si se midiera un sistema que contiene todas las posibles soluciones superpuestas, se obtendría una respuesta aleatoria, sin valor computacional \cite{nielsen2010}.

La interferencia cuántica se convierte en el mecanismo indispensable para extraer información útil. De esta manera, al igual que las ondas electromagnéticas, las amplitudes de probabilidad de los estados cuánticos pueden ser números complejos con fases distintas. Los algoritmos cuánticos aprovechan esta naturaleza ondulatoria aplicando transformaciones unitarias (puertas lógicas cuánticas) que manipulan dichas fases antes de efectuar la medición, logrando amplificar la probabilidad de medir la solución correcta y cancelando las incorrectas.

El objetivo de estas transformaciones es provocar una interferencia destructiva en las trayectorias que conducen a respuestas incorrectas, causando que sus amplitudes se cancelen, y una interferencia constructiva en los estados que representan la solución correcta. 

En el caso específico del algoritmo de Shor, este proceso de amplificación de probabilidades se logra mediante la aplicación de la Transformada Cuántica de Fourier QFT\cite{nielsen2010}. La QFT actúa sobre un estado computacional básico $|j\rangle$ aplicando la siguiente transformación:

$$|j\rangle \mapsto \frac{1}{\sqrt{N}} \sum_{k=0}^{N-1} e^{2\pi i j k / N} |k\rangle$$

Esta formulación revela la periodicidad oculta de una función mediante el control de las fases exponenciales complejas, un paso matemático que permite reducir la complejidad de la factorización de enteros de un tiempo de ejecución exponencial a uno polinómico \cite{nielsen2010}.

\subsection{El Problema de la Factorización de Enteros y la vulnerabilidad de RSA}

El sistema criptográfico RSA \cite{rivest1978}, pilar fundamental de la infraestructura de clave pública moderna, basa su seguridad en la asimetría computacional de una función matemática. Se basa en la intratabilidad del Problema de la Factorización de Enteros IFP.

El IFP dice así: dados dos números primos aleatorios y suficientemente grandes, $p$ y $q$, es computacionalmente trivial multiplicarlos para obtener su producto $N = p \times q$. Esta operación aritmética puede realizarse de manera eficiente en tiempo polinómico. Sin embargo, la operación inversa, deducir los factores primos originales $p$ y $q$ conociendo únicamente el módulo público $N$, representa un desafío monumental para la arquitectura de computación clásica.

A día de hoy, el algoritmo más eficiente conocido para factorizar enteros grandes es la Criba General del Cuerpo de Números GNFS. La complejidad temporal de este algoritmo escala de forma subexponencial respecto a la longitud en bits de la clave. Esto significa que, para longitudes de clave estándar en la industria actual (como RSA de 2048 o 4096 bits), la factorización requeriría recursos energéticos y temporales inasumibles, garantizando así la invulnerabilidad práctica del esquema.

El problema es que esta robustez es estrictamente dependiente de las limitaciones impuestas por la física clásica. Es precisamente esta barrera de complejidad la que el algoritmo de Shor logra vulnerar. Al transformar el problema de la factorización en un problema de búsqueda de período menor mediante la QFT, Shor reduce la complejidad temporal de la factorización de un orden subexponencial a uno estrictamente polinómico. Esta drástica reducción es la que invalida las garantías de seguridad de RSA frente a la computación cuántica.


\subsection{Vulnerabilidad de la Criptografía de Curva Elíptica ECC}

RSA fue el primer estándar de clave pública ampliamente adoptado, pero la mayor parte de la infraestructura digital contemporánea, incluyendo protocolos como TLS/SSL, mensajería cifrada y criptomonedas, ha migrado hacia la Criptografía de Curva Elíptica ECC \cite{miller1985}. La adopción masiva de ECC se debe a su superioridad en eficiencia: ofrece el mismo nivel de seguridad algorítmica que RSA utilizando tamaños de clave significativamente menores.

La seguridad de este esquema no reside en la factorización, sino en el Problema del Logaritmo Discreto en Curvas Elípticas ECDLP. Dado un punto generador $P$ sobre una curva elíptica definida en un cuerpo finito, y un escalar privado $k$ (que actúa como clave privada), se calcula la clave pública $Q$ mediante la multiplicación escalar:

$$ Q = k \cdot P $$

En la computación clásica, recuperar el valor de $k$ a partir de $Q$ y $P$ requiere algoritmos de tiempo de ejecución exponencial, como el algoritmo rho de Pollard, lo que hace que el problema sea intratable para claves de 256 bits o superiores.

Sin embargo, el alcance del algoritmo cuántico de Shor se extiende más allá de la factorización de enteros. En su formulación más general, el algoritmo resuelve el \textit{Problema del Subgrupo Oculto} para grupos abelianos finitos, categoría matemática a la que pertenece el ECDLP \cite{shor1994}. Mediante una variante de la Transformada Cuántica de Fourier, un ordenador cuántico puede deducir la clave privada $k$ en tiempo polinómico.

Aún más crítico es el hecho de que la principal ventaja de ECC en la computación clásica se convierte en su mayor debilidad en la era cuántica. Investigaciones, como las de Proos y Zalka \cite{proos2003}, han demostrado que un ordenador cuántico requiere significativamente menos \textit{qubits} lógicos para resolver el ECDLP de una clave de 256 bits que para factorizar un módulo RSA de 2048 bits. En consecuencia, se prevé que la criptografía basada en curvas elípticas será comprometida antes que los esquemas RSA tradicionales, evidenciando la urgencia extrema de adoptar el estándar post-cuántico FIPS 204 \cite{nist_fips204}.


\subsection{Criptografía Simétrica y el Algoritmo de Grover}

Aunque la vulnerabilidad arquitectónica más crítica recae sobre la criptografía asimétrica debido a la eficiencia del algoritmo de Shor, es necesario analizar el impacto de la computación cuántica sobre la criptografía simétrica como el \textit{Advanced Encryption Standard} AES \cite{nist2001aes}. El vector de amenaza principal no proviene de Shor, sino del algoritmo de búsqueda propuesto por Lov Grover en 1996 \cite{grover1996fast}.

El algoritmo de Grover está diseñado para la búsqueda óptima en espacios de estados no estructurados. Encontrar una clave criptográfica en un espacio de tamaño $N$ mediante fuerza bruta requiere, en el peor de los casos, un tiempo lineal $\mathcal{O}(N)$. Grover demostró que un sistema cuántico fundamentado en la amplificación de amplitud puede aislar la solución en un tiempo proporcional a $\mathcal{O}(\sqrt{N})$.

Aplicado a la criptografía simétrica, este algoritmo reduce efectivamente los \textit{bits} de seguridad de la clave a la mitad. Para el estándar AES-128, el espacio de búsqueda clásico de $2^{128}$ combinaciones se reduce a $\sqrt{2^{128}} = 2^{64}$ operaciones cuánticas. Dado que un margen de 64 \textit{bits} de seguridad es insuficiente frente a un escenario post-cuántico consolidado, AES-128 queda catalogado como vulnerable ante ataques asistidos por hardware cuántico \cite{nistir8105}.

A diferencia de la criptografía de clave pública, para evitar esta vulnerabilidad no requiere el diseño de algoritmos basados en retículos o códigos. La solución algorítmica y estandarizada para mantener la seguridad simétrica es duplicar la longitud de la clave matriz. Al implementar AES-256, el esfuerzo cuántico requerido por el algoritmo de Grover asciende a $2^{128}$ operaciones, restaurando el margen computacional necesario para garantizar un nivel de seguridad inquebrantable \cite{nistir8105}.

Por lo tanto, la transición hacia arquitecturas \textit{Quantum-Safe} exige una estrategia híbrida: no solo la adopción de estándares como FIPS 204 para blindar las firmas digitales asimétricas, sino la obligatoriedad estricta de establecer AES-256 como el cifrado simétrico mínimo indispensable en la capa de red.


\section{Transición hacia la Criptografía Post-Cuántica}

La amenaza cuántica plantea un desafío temporal crítico. La migración completa de la infraestructura criptográfica global ha requerido décadas. El despliegue de un ordenador cuántico criptográficamente relevante podría ocurrir antes de que los sistemas actuales sean reemplazados. Por ello, el Instituto Nacional de Estándares y Tecnología de los Estados Unidos (NIST) inició en 2016 el Proyecto de Estandarización de Criptografía Post-Cuántica \cite{chen2016report}.

El objetivo de este programa de estandarización impulsado por el NIST no era el desarrollo de hardware cuántico, sino identificar y consolidar algoritmos criptográficos que, ejecutándose sobre arquitecturas clásicas de Von Neumann, ofrecieran seguridad matemática frente a ataques tanto clásicos como cuánticos \cite{nist2016call}. Para lograrlo, la convocatoria solicitó y evaluó propuestas fundamentadas en familias de problemas matemáticos, destacando principalmente la criptografía basada en códigos, en funciones \textit{hash} y en retículos.

Durante más de seis años, los algoritmos candidatos fueron sometidos a la investigación del criptoanálisis mundial a través de múltiples rondas de evaluación. Este proceso de selección demostró la extrema dificultad de diseñar criptografía post-cuántica robusta; candidatos que alcanzaron rondas finales fueron comprometidos en cuestión de horas utilizando hardware clásico, evidenciando debilidades.

Finalmente, en julio de 2022, el NIST anunció los algoritmos ganadores para la estandarización. La criptografía basada en retículos emergió como la familia matemática predominante debido a su excelente equilibrio entre tamaño de clave, velocidad de ejecución y garantías teóricas de seguridad en el peor de los casos. 

Para la firma digital primaria, el esquema seleccionado fue CRYSTALS-Dilithium. Tras un proceso de refinamiento de parámetros para asegurar su robustez técnica y facilitar su implementación en la industria, el NIST publicó oficialmente la especificación final en agosto de 2024 bajo el estándar FIPS 204 \cite{nist_fips204}. Este documento formaliza el algoritmo bajo el acrónimo ML-DSA (\textit{Module-Lattice-Based Digital Signature Algorithm}), estableciéndolo como el nuevo pilar global para la autenticación y no repudio en la era post-cuántica, y convirtiéndose en el objeto central de estudio e implementación del presente Trabajo de Fin de Grado.


\section{Objetivos del Trabajo}

\subsection{Objetivo General}

El objetivo principal de este Trabajo de Fin de Grado es analizar, comprender e implementar de forma integral el estándar de firma digital post-cuántica ML-DSA (FIPS 204), garantizando técnica y operativa mediante la construcción de un entorno distribuido real que opere sobre arquitecturas computacionales clásicas.

\subsection{Objetivos Específicos}

Para alcanzar los objetivos de manera metódica y estructurada, se definen los siguientes hitos técnicos:

\begin{itemize}
    \item \textbf{Asimilación de fundamentos algebraicos:} Estudiar en profundidad las bases matemáticas de la criptografía basada en retículos, comprendiendo la intratabilidad de los problemas \textit{Module Learning With Errors} MLWE y \textit{Module Short Integer Solution} MSIS que garantizan la seguridad teórica del esquema.
    
    \item \textbf{Análisis crítico del estándar FIPS 204:} Desglosar y sintetizar la especificación formal publicada por el NIST, detallando la arquitectura interna, los parámetros de seguridad y las funciones del algoritmo ML-DSA.
    
    \item \textbf{Análisis del código y la firma post-cuántica:} Analizar detalladamente el código y la implementación práctica del esquema de firma.
    \item \textbf{Desarrollo del motor criptográfico:} Implementar desde cero, utilizando el lenguaje de programación Python, un núcleo criptográfico funcional que ejecute con exactitud la generación de claves, firma de documentos y verificación, adaptándose estrictamente a la especificación oficial.
    
    \item \textbf{Diseño e integración del sistema distribuido:} Construir una plataforma integral compuesta por una API RESTful asíncrona en el \textit{backend} y una aplicación de página única (SPA) en el \textit{frontend}. Este ecosistema debe ofrecer un entorno seguro y usable para la gestión de identidades y la delegación de firmas poscuánticas.
    
    \item \textbf{Validación de rendimiento en hardware clásico:} Demostrar empíricamente que la ejecución de algoritmos post-cuánticos avanzados no requiere un cambio de paradigma de \textit{hardware}, comprobando que el motor ML-DSA es eficiente y plenamente viable sobre arquitecturas informáticas clásicas de Von Neumann.
\end{itemize}

\section{Estructura de la Memoria}

Para dar cumplimiento a los objetivos planteados, la memoria se estructura en los siguientes capítulos:

\begin{itemize}
    \item \textbf{Capítulo 1. Introducción:} Contextualiza el problema de la amenaza cuántica, justifica la motivación del proyecto frente a la vulnerabilidad de la criptografía clásica y define los objetivos del Trabajo de Fin de Grado.
    
    \item \textbf{Capítulo 2. Fases de la realización del trabajo y cronogramas:} Detalla la metodología de ingeniería de \textit{software} empleada, la planificación temporal del proyecto y la gestión de los recursos para su ejecución.
    
    \item \textbf{Capítulo 3. Fundamentos Matemáticos:} Aborda la base algebraica necesaria para comprender la criptografía post-cuántica, centrándose en la teoría de retículos, anillos polinómicos y problemas como MLWE.
    
    \item \textbf{Capítulo 4. El Esquema de Firma ML-DSA (FIPS 204):} Analiza el diseño teórico del algoritmo, explicando los mecanismos de generación de claves, el proceso de firma basado en el paradigma \textit{Fiat-Shamir con abortos} y la verificación matemática.\cite{lyubashevsky2009fiat}
    
    \item \textbf{Capítulo 5. Implementación del Motor Criptográfico en Python:} Detalla el desarrollo algorítmico y la traducción estricta del estándar FIPS 204 a código. Documenta la gestión de la entropía de seguridad, la programación del núcleo algebraico polinómico (NTT/$\text{NTT}^{-1}$) y la validación del sistema mediante una suite de pruebas basadas en \textit{Test-Driven Development} (TDD).
    
    \item \textbf{Capítulo 6. Diseño y Desarrollo del Sistema Distribuido:} Presenta la integración de la librería criptográfica en una arquitectura cliente-servidor. Se analiza la infraestructura tecnológica, el diseño de la API REST asíncrona, la arquitectura del cliente web (SPA) junto con su gestión de estado y control de accesos, el motor de procesamiento documental para la inyección de firmas en formato PDF y la fortificación perimetral del sistema contra vectores de ataque web.
    
    \item \textbf{Capítulo 7. Conclusiones y Trabajo Futuro:} Recopila los resultados finales obtenidos, evalúa el grado de cumplimiento de los objetivos iniciales y plantea líneas de investigación para la escalabilidad y transición hacia arquitecturas cero conocimiento.
    
\end{itemize}

\chapter{Fases de la realización del trabajo y cronogramas}

\section{Metodología de Trabajo}

Debido a la gran dificultad matemática de este Trabajo de Fin de Grado, la aplicación de una única metodología de \textit{software} tradicional era insuficiente. El proyecto exigía un análisis profundo de la criptografía basada en retículos antes de poder escribir la primera línea de código. Por ello, se ha optado por aplicar un enfoque metodológico híbrido, estructurado en dos grandes ejes:

\begin{enumerate}
    \item \textbf{Metodología basada en Investigación y Desarrollo (I+D):} Aplicada al núcleo del proyecto, que es el motor criptográfico. En esta fase el estudio se centró en la criptografía basada en retículos y el análisis de la especificación técnica del FIPS 204. Se siguió un proceso iterativo de ingeniería inversa teórica: traducir las ecuaciones matemáticas y las funciones del NIST a diagramas de flujo lógicos y, posteriormente, a algoritmos para poder implementarlos.
    
    \item \textbf{Metodología Iterativa e Incremental (Enfoque Ágil):} Aplicada a la capa de integración de \textit{software} y arquitectura web. Una vez consolidado el núcleo criptográfico, el desarrollo de la API y la interfaz de usuario se gestionó mediante ciclos cortos de desarrollo, prueba y validación. Esto permitió una integración continua (\textit{Bottom-Up}), conectando las rutinas criptográficas  con la plataforma de la red.
\end{enumerate}

Trabajando de esta manera, se permite proporcionar un rigor exhaustivo al estándar de firma digital y, por otro, la agilidad necesaria para desplegar una plataforma funcional en los plazos establecidos.


\section{Fases del Proyecto}

El desarrollo del trabajo se ha estructurado en cinco fases secuenciales e incrementales, para evitar los riesgos técnicos de complejidad de la criptografía post-cuántica:

\begin{itemize}
    \item \textbf{Fase 1: Análisis Teórico y Asimilación de Estándares (Enero - Principios de Febrero).} Estudio exhaustivo y asimilación de los fundamentos matemáticos de los retículos y la especificación original FIPS 204 del NIST.
    
    \item \textbf{Fase 2: Diseño Arquitectónico y Modelado (Febrero).} Mapeo de dependencias de los algoritmos criptográficos y elaboración de diagramas de arquitectura del sistema. En esta fase se estructuró el árbol de directorios y se definieron las interfaces lógicas entre los módulos de generación de clave, firma y verificación.
    
    \item \textbf{Fase 3: Implementación del Core Criptográfico (Finales de Febrero - Principios de Marzo).} Traducción de las matemáticas a código fuente en Python. Se implementaron de forma aislada las primitivas criptográficas y se desarrollaron \textit{scripts} de validación interna para asegurar la fidelidad con el estándar.
    
    \item \textbf{Fase 4: Integración Full-Stack y Despliegue (Marzo).} Construcción de la plataforma. Incluyó el desarrollo de una API REST protegida mediante \textit{JSON Web Tokens} JWT, la persistencia de datos apoyada en Supabase y la creación ágil del \textit{frontend} de usuario interactivo.
    
    \item \textbf{Fase 5: Consolidación, Pruebas y Documentación (Abril - Finales de Mayo).} Elaboración de la presente memoria técnica, integración de los diagramas de diseño generados en fases previas, pruebas de rendimiento finales y preparación del material para la defensa ante el tribunal.
\end{itemize}

\section{Planificación Temporal}

Para mostrar la distribución temporal de este TFG, que hacen una duración total de 5 meses, se presenta el siguiente cronograma. Esta planificación muestra la dedicación inicial a la fase algorítmica y matemática y la posterior aceleración en las fases de despliegue web y documentación.

\vspace{0.5cm}
\begin{table}[h]
\centering
\renewcommand{\arraystretch}{1.5}
\begin{tabular}{|l|c|c|c|c|c|}
\hline
\textbf{Fases del Proyecto} & \textbf{Ene} & \textbf{Feb} & \textbf{Mar} & \textbf{Abr} & \textbf{May} \\ \hline
Fase 1: Análisis Teórico y FIPS 204 & X & X (inicio) & & & \\ \hline
Fase 2: Diseño Arquitectónico & & X  & & & \\ \hline
Fase 3: Core Criptográfico (Python) & & X (final) & X (inicio) & & \\ \hline
Fase 4: Integración API y Web & & & X & & \\ \hline
Fase 5: Documentación y Cierre & & & & X & X \\ \hline
\end{tabular}
\caption{Cronograma de ejecución del proyecto}
\label{tab:cronograma}
\end{table}
\vspace{0.5cm}

\section{Entorno Tecnológico y Herramientas}

La elección del \textit{stack} tecnológico se ha guiado por la necesidad de equilibrar la exigencia matemática de la criptografía post-cuántica, con la agilidad requerida para el despliegue de una arquitectura web moderna. A continuación, se detallan y justifican las herramientas principales utilizadas durante el ciclo de vida del proyecto:

\begin{itemize}
    \item \textbf{Python 3:} Seleccionado como lenguaje principal para el desarrollo del motor criptográfico ML-DSA. Su sintaxis clara y su extenso ecosistema de bibliotecas matemáticas lo convierten en el estándar \textit{de facto} para la investigación y traducción de especificaciones algebraicas complejas a código funcional.
    
    \item \textbf{Draw.io (diagrams.net):} Herramienta empleada en la fase de diseño arquitectónico para el modelado visual de dependencias. Resultó fundamental para desglosar la jerarquía de los más de 50 algoritmos y rutinas que componen el estándar FIPS 204 antes de su implementación.
    
    \item \textbf{Git y GitHub:} Sistema de control de versiones utilizado para la gestión del código fuente, asegurando la trazabilidad de los cambios y protegiendo la integridad del proyecto mediante un repositorio seguro en la nube.
    
    \item \textbf{Supabase:} Plataforma \textit{Backend-as-a-Service} BaaS de código abierto basada en PostgreSQL. Se eligió para la capa de persistencia de datos debido a su facilidad de integración, escalabilidad y soporte nativo para arquitecturas REST.
    
    \item \textbf{JSON Web Tokens JWT:} Estándar abierto (RFC 7519) implementado para securizar la API REST. Proporciona un mecanismo de autenticación \textit{stateless}, ideal para microservicios y plataformas distribuidas de firma digital.
    
    \item \textbf{React y TypeScript (Ecosistema Frontend):} La interfaz de usuario se ha desarrollado como una aplicación de página única SPA utilizando la biblioteca React. Su naturaleza declarativa y basada en componentes permite gestionar de forma eficiente el ciclo de vida de las vistas y el estado de la aplicación. La adopción de TypeScript como superconjunto de JavaScript aporta un tipado estático estricto, previniendo errores en tiempo de compilación y garantizando una integración robusta con las estructuras de datos devueltas por la API.
    
    \item \textbf{LaTeX y Overleaf:} Sistema de composición de textos empleado para la redacción de esta memoria académica. Garantiza la máxima calidad tipográfica, indispensable para la correcta notación de las fórmulas matemáticas, retículos y estados cuánticos descritos en el documento.
\end{itemize}

\section{Estimación de Costes}

Este apartado recoge la parte económica asociada al desarrollo del proyecto. Para evaluar su viabilidad más allá del ámbito puramente académico, se ha diseñado un escenario presupuestario hipotético que muestra el esfuerzo y los recursos utilizados en un entorno empresarial real. 

Los gastos se han estructurado en dos bloques principales: costes directos, haciendo referencia a las horas de ingeniería, y costes indirectos, tales como los costes operativos y de infraestructura tecnológica.

\subsection{Costes de Personal}
El volumen de trabajo del proyecto asciende a 300 horas, correspondientes a la carga de 12 créditos ECTS. Para el cálculo salarial, se ha tomado como referencia el perfil de un Ingeniero de \textit{Software} Junior especializado en Ciberseguridad, aplicando una tarifa estándar de mercado de 25 \euro{}/hora.

\subsection{Costes Indirectos y de Infraestructura}
Se contempla un periodo de desarrollo de cinco meses, por lo tanto se le asocian los siguientes gastos:
\begin{itemize}
    \item \textbf{Equipamiento Informático:} Cuota de amortización de un equipo de desarrollo de altas prestaciones, valor económico de 1.500 \euro{}.
    \item \textbf{Licencias y Servicios \textit{Cloud}:} Estimación del coste comercial de los servicios utilizados, bases de datos en Supabase, entornos \textit{Low-Code} y control de versiones, que, aunque disponen de capas gratuitas, requerirían pago en un entorno de producción.
    \item \textbf{Costes Operativos:} Parte proporcional de gastos derivados del espacio de trabajo, consumo eléctrico y conectividad a Internet.
\end{itemize}

\subsection{Presupuesto Total}
La suma de estos conceptos conforma el Presupuesto de Ejecución Material (PEM), detallado en la Tabla \ref{tab:presupuesto}. Al aplicar el 21\% de IVA vigente, se obtiene el coste total de desarrollo del proyecto.

\vspace{0.5cm}
\begin{table}[h]
\centering
\renewcommand{\arraystretch}{1.5}
\begin{tabular}{|l|l|r|}
\hline
\textbf{Categoría} & \textbf{Concepto} & \textbf{Coste Estimado (\euro{})} \\ \hline
\multirow{1}{*}{Costes Directos} 
 & Ingeniero Junior (300 h $\times$ 25 \euro{}/h) & 7.500,00 \euro{} \\ \hline
\multirow{3}{*}{Costes Indirectos} 
 & Amortización de Equipo (5 meses) & 156,25 \euro{} \\ 
 & Licencias y Servicios Tecnológicos & 120,00 \euro{} \\ 
 & Gastos Operativos (Energía, Internet) & 90,00 \euro{} \\ \hline
\multicolumn{2}{|r|}{\textbf{Subtotal (PEM)}} & \textbf{7.866,25 \euro{}} \\ \hline
\multicolumn{2}{|r|}{\textbf{IVA (21\%)}} & \textbf{1.651,91 \euro{}} \\ \hline
\multicolumn{2}{|r|}{\textbf{PRESUPUESTO TOTAL}} & \textbf{9.518,16 \euro{}} \\ \hline
\end{tabular}
\caption{Presupuesto estimado del proyecto}
\label{tab:presupuesto}
\end{table}
\vspace{0.5cm}

\chapter{Fundamentos Matemáticos}

En este capítulo se describen los pilares algebraicos sobre los que se construye el estándar ML-DSA. La seguridad de la criptografía post-cuántica basada en retículos no depende de la dificultad de factorizar grandes números, sino de la intratabilidad de problemas geométricos en espacios de alta dimensionalidad, estructurados sobre anillos de polinomios.

\section{Estructuras Algebraicas: El Anillo $R_q$}

La unidad mínima de información en Dilithium no es el bit ni el número entero, sino polinomios dentro del anillo cociente $$R_q = \mathbb{Z}_q[x]/(x^{256} + 1)$$. 

\subsection{Aritmética Modular y Reducción Simétrica}
\label{subsec:reduccion_simetrica}

Un concepto fundamental en la implementación de ML-DSA es el tratamiento de los coeficientes mediante la \textbf{reducción modular simétrica} (o centrada). A diferencia de la aritmética modular convencional, donde los resultados se sitúan estrictamente en el rango positivo $[0, q-1]$, en este esquema se busca trabajar en un entorno simétrico respecto al origen, proyectando los valores en el intervalo:

$$ -\lfloor q/2 \rfloor \le r \le \lfloor q/2 \rfloor $$

Este enfoque garantiza que los elementos de error y las claves secretas mantengan una norma pequeña, un factor matemático determinante para dos procesos críticos del algoritmo: la compresión de las firmas digitales y el funcionamiento de los filtros de rechazo basados en la norma infinita ($||\cdot||_\infty$). Mantener los valores oscilando cerca del cero previene el crecimiento descontrolado del ruido durante las complejas multiplicaciones polinómicas de la NTT.

Las Figuras \ref{fig:modulo_estandar} y \ref{fig:modulo_simetrico} ilustran visualmente esta diferencia de representación en el anillo $\mathbb{Z}_q$. En la representación simétrica, los coeficientes se dividen en zonas: los de bajo peso y los de alto peso. Esta división es la que utiliza internamente ML-DSA para calcular los \textit{hint bits} (bits de pista), que se explicarán más adelante, y actúan como indicadores que señalan si un valor original ha caído en la mitad superior o inferior del anillo.

\begin{figure}[H]
    \centering
    \begin{tikzpicture}
        \def\radio{3}
        \def\modulo{17}

        \draw[thick] (0,0) circle (\radio);

        \node at (0,0) {\huge $\mathbb{Z}_{17}$};

        \foreach \i in {0,...,16} {
            \pgfmathsetmacro{\angulo}{90 - \i * (360/\modulo)}
            
            \fill (\angulo:\radio) circle (3pt);
            
            \node at (\angulo:\radio + 0.5) {\large $\i$};
        }
    \end{tikzpicture}
    \caption{Aritmética modular estándar en el anillo $\mathbb{Z}_{17}$.}
    \label{fig:modulo_estandar}
\end{figure}

\begin{figure}[H]
    \centering
    \begin{tikzpicture}
        \def\radio{3}
        \def\modulo{17}

        \draw[thick] (0,0) circle (\radio);

        \node at (0,0) {\huge $\mathbb{Z}_{17}$};

        \foreach \i in {0,...,8} {
            \pgfmathsetmacro{\angulo}{90 - \i * (360/\modulo)}
            
            \fill (\angulo:\radio) circle (3pt);
            \node at (\angulo:\radio + 0.5) {\large $\i$};
        }

        \foreach \i in {9,...,16} {
            \pgfmathsetmacro{\angulo}{90 - \i * (360/\modulo)}
            
            \fill (\angulo:\radio) circle (3pt);
            
            \pgfmathsetmacro{\valorsimetrico}{int(\i - \modulo)}
            
            \node[text=red] at (\angulo:\radio + 0.5) {\large $\valorsimetrico$};
        }
    \end{tikzpicture}
    \caption{Aritmética modular simétrica en el anillo $\mathbb{Z}_{17}$.}
    \label{fig:modulo_simetrico}
\end{figure}

\begin{ejemplo}[Reducción Modular Simétrica vs. Estándar]
\label{ex:reduccion_simetrica}
Para ilustrar el proceso, tomemos el anillo reducido $\mathbb{Z}_{17}$. 

En la aritmética modular estándar, los valores oscilan en el rango $[0, 16]$. Sin embargo, en la representación simétrica empleada por ML-DSA, los valores se proyectan en el intervalo $[-8, 8]$.

Si evaluamos un coeficiente con valor $10$:
\begin{itemize}
    \item \textbf{Reducción estándar:} $10 \pmod{17} = 10$.
    \item \textbf{Reducción simétrica:} Como $10$ excede la cota superior del intervalo simétrico ($8$), da la vuelta al anillo y se representa como su equivalente negativo restando el módulo: $10 - 17 = -7$.
\end{itemize}

Otro ejemplo con un valor que sobrepasa el módulo, como $25$:
\begin{itemize}
    \item \textbf{Reducción estándar:} $25 \pmod{17} = 8$.
    \item \textbf{Reducción simétrica:} Como el equivalente estándar es $8$, el cual se encuentra dentro del límite superior permitido para $q=17$, no requiere ajuste negativo. El resultado es directamente $8$.
\end{itemize}
\end{ejemplo}

En el estándar oficial ML-DSA, este mismo principio algebraico se aplica de forma estricta sobre el módulo del número primo $q = 8380417$. La reducción simétrica mapea todos los coeficientes de los polinomios al intervalo exacto $[-4190208, 4190208]$, garantizando la estabilidad de las operaciones de compresión como la función \texttt{LowBits} empleada durante la fase de firma.

\subsection{El Anillo de Polinomios $R_q$}

Todas las operaciones algebraicas de ML-DSA ocurren dentro del anillo cociente \cite{peikert2016decade}:

$$ R_q = \mathbb{Z}_q[x]/(x^{256} + 1) $$

donde $q = 8380417$ es un número primo. Esto implica que los elementos del criptosistema son polinomios de grado máximo 255. La clave estructural de este anillo es que cualquier operación de multiplicación que genere términos de grado 256 o superior se reduce utilizando la equivalencia polinómica $x^{256} \equiv -1$.

Para comprender la mecánica de esta reducción sin la complejidad de operar con 256 coeficientes, consideremos un anillo de dimensiones reducidas.

\begin{ejemplo}[Reducción Polinómica en un Anillo Cociente]
\label{ex:reduccion_anillo}
Consideremos un anillo con $n=4$, definido como $R_q = \mathbb{Z}_q[x]/(x^4 + 1)$. En este entorno, la regla de reducción modular dicta que $x^4 \equiv -1$.

Supongamos dos polinomios simples $A(x) = x^3 + 2x$ y $B(x) = x^2 + 3$. Su multiplicación es:
\[ A(x) \cdot B(x) = (x^3 + 2x)(x^2 + 3) = x^5 + 3x^3 + 2x^3 + 6x = x^5 + 5x^3 + 6x \]

Al aplicar la reducción del anillo $(x^4 + 1)$, el término de grado 5 excede el límite dimensional del sistema ($n-1 = 3$). Dado que $x^4 \equiv -1 \pmod{x^4 + 1}$, podemos descomponer y reescribir $x^5$ como $x^4 \cdot x$, lo que equivale a $-1 \cdot x = -x$. Sustituyendo este valor en el polinomio resultante obtenemos:
\[-x + 5x^3 + 6x \equiv 5x^3 + 5x \]

El polinomio resultante vuelve a tener un grado estrictamente menor que $n=4$, y el "desbordamiento" se ha reinyectado en los grados inferiores con el signo invertido.
\end{ejemplo}

En el estándar FIPS 204, el motor criptográfico realiza exactamente este mismo proceso matemático, pero operando con $n=256$ y aplicando de forma simultánea la reducción modular simétrica $q$ explicada anteriormente sobre cada uno de los coeficientes.

\section{Operaciones sobre el Anillo y Aceleración Algorítmica}

Dado que los elementos de las claves y las firmas en ML-DSA son matrices y vectores de polinomios en $R_q$, la operación aritmética más crítica y repetitiva del estándar es la multiplicación de polinomios. 

\subsection{Producto de Polinomios y Convolución Negacíclica}

En el álgebra clásica, la multiplicación de dos polinomios $a(x)$ y $b(x)$ de grado $n-1$ da como resultado un polinomio de grado $2n-2$. Sin embargo, al operar dentro del anillo cociente $R_q = \mathbb{Z}_q[x]/(x^{256} + 1)$, cualquier término de grado igual o superior a $n=256$ debe ser reducido. Esta operación se puede realizar teniendo en cuenta que:

$$x^{256} \equiv -1 \pmod{x^{256} + 1}$$

De esta forma, al desbordar el grado máximo, los coeficientes reaparecen en los grados inferiores con el signo invertido. El problema oculto de esta forma de multiplicar polinomios es su complejidad computacional: realizar la multiplicación directa exige que cada coeficiente del primer polinomio se multiplique por todos los coeficientes del segundo, lo que arroja una complejidad asintótica cuadrática de $\mathcal{O}(n^2)$. Para un polinomio de $n=256$, esto supone 65.536 multiplicaciones escalares por cada par de polinomios, un coste inasumible para un sistema criptográfico de alto rendimiento.

\subsection{La Transformada Numérica Teórica NTT}

Para resolver el cuello de botella computacional descrito, el estándar FIPS 204 exige el uso de la Transformada Numérica Teórica NTT \cite{ducas2018crystals}. 

La NTT es una especialización de la Transformada Discreta de Fourier DFT que, en lugar de operar sobre números complejos (los cuales introducirían errores de precisión por el uso de coma flotante), opera estrictamente sobre el cuerpo finito $\mathbb{Z}_q$. Matemáticamente, proyecta los coeficientes del polinomio desde el \textit{dominio estándar} al \textit{dominio transformado} aprovechando las raíces primitivas de la unidad ($2n$-ésimas) que existen dentro de $\mathbb{Z}_q$.

Para ilustrar el funcionamiento y la ventaja computacional de la Transformada Numérica Teórica NTT, es fundamental comprender que el algoritmo no opera sobre vectores abstractos, sino sobre las representaciones vectoriales de los coeficientes de los polinomios.

\begin{ejemplo}[Multiplicación Polinómica vía NTT en un Anillo Reducido]
\label{ex:ntt_juguete}
Supongamos un anillo sencillo con $n=2$ y un módulo primo $q=5$, es decir, $R_5 = \mathbb{Z}_5[x]/(x^2 + 1)$. Definimos dos polinomios y sus respectivos vectores de coeficientes:
\[a(x) = 1 + x \quad \implies \quad a = (1, 1)\]
\[b(x) = 2 + x \quad \implies \quad b = (2, 1)\]

Si realizamos la multiplicación polinómica tradicional en este anillo, obtenemos:
\[c(x) = a(x) \cdot b(x) = (1+x)(2+x) = 2 + 3x + x^2\]
Aplicando la reducción del anillo ($x^2 \equiv -1 \pmod{x^2 + 1}$), obtenemos el resultado final:
\[c(x) = 2 + 3x - 1 = 1 + 3x \quad \implies \quad c = (1, 3)\]
Este método clásico requiere $\mathcal{O}(n^2)$ operaciones. Para $n=2$ es trivial, pero para $n=256$ (el estándar en ML-DSA) resulta prohibitivo.

Para optimizarlo, interviene el \textbf{Teorema de Convolución} en cuerpos finitos, que establece que $c = \text{NTT}^{-1}(\text{NTT}(a) \circ \text{NTT}(b))$. El procedimiento transforma los coeficientes evaluando los polinomios en las raíces de la unidad del anillo (para este caso, las raíces son $\gamma^1 = 2$ y $\gamma^3 = 3$):

\begin{enumerate}
    \item \textbf{Transformada Directa NTT:} Se evalúan $a(x)$ y $b(x)$ en las raíces.
    \begin{itemize}
        \item $\text{NTT}(a) = (a(2), a(3)) = (1+2, 1+3) = (3, 4)$
        \item $\text{NTT}(b) = (b(2), b(3)) = (2+2, 2+3) = (4, 5) \equiv (4, 0) \pmod{5}$
    \end{itemize}
    \item \textbf{Producto de Hadamard $\circ$:} En el dominio transformado, la compleja convolución se reduce a una simple multiplicación elemento a elemento en módulo $5$:
    \[\hat{c} = (3 \times 4 \pmod{5},\; 4 \times 0 \pmod{5}) = (2, 0)\]
    Este paso requiere únicamente $\mathcal{O}(n)$ operaciones.
    \item \textbf{Transformada Inversa $\text{NTT}^{-1}$:} Se recupera el polinomio $c(x)$ que verifica que $c(2)=2$ y $c(3)=0$. Matemáticamente, este proceso se reduce a un problema clásico y sencillo de interpolación de polinomios. Al aplicar el algoritmo inverso, se recupera exactamente:
    \[c = (1, 3) \quad \implies \quad c(x) = 1 + 3x\]
\end{enumerate}
\end{ejemplo}

Para el caso específico del estándar ML-DSA ($n=256$), el hardware clásico ejecutaría la multiplicación en aproximadamente $65.536$ operaciones modulares ($\mathcal{O}(n^2)$). Al mapear los polinomios a vectores y usar la NTT, el coste se reduce a $3 \times (256 \log_2 256) + 256 \approx 6.400$ operaciones. Esta drástica reducción algorítmica es el motor que hace viables las firmas post-cuánticas en entornos de producción.

\section{Álgebra Lineal sobre Retículos}

Mientras que la aceleración algorítmica de Dilithium recae en la NTT, su seguridad fundamental descansa sobre la complejidad de la solución de problemas sobre retículos. En este apartado se define la estructura matemática que da nombre a la familia de algoritmos a la que pertenece ML-DSA.

\subsection{Definición Geométrica y Transición a Módulos}

En su definición formal, un retículo $\mathcal{L}$ es un subgrupo aditivo discreto del espacio euclídeo $\mathbb{R}^n$. Puede visualizarse como una cuadrícula infinita de puntos regulares. Se genera a partir de combinaciones lineales enteras de un conjunto de vectores base linealmente independientes $\{ \mathbf{b}_1, \dots, \mathbf{b}_n \}$ \cite{peikert2016decade}:

$$ \mathcal{L} = \left\{ \sum_{i=1}^n x_i \mathbf{b}_i \;\middle|\; x_i \in \mathbb{Z} \right\} $$

\begin{figure}[H]
    \centering
    \begin{tikzpicture}[scale=0.85]
        \def\voneX{1.5} \def\voneY{1.5}
        \def\vtwoX{1.5} \def\vtwoY{3.5}

        \clip (-5, -4.5) rectangle (5, 4.5);

        \foreach \i in {-10,...,10} {
            \draw[dashed, gray!80] (\i*\vtwoX - 15*\voneX, \i*\vtwoY - 15*\voneY) -- (\i*\vtwoX + 15*\voneX, \i*\vtwoY + 15*\voneY);
            \draw[dashed, gray!80] (\i*\voneX - 15*\vtwoX, \i*\voneY - 15*\vtwoY) -- (\i*\voneX + 15*\vtwoX, \i*\voneY + 15*\vtwoY);
        }

        \draw (-6, 0) -- (6, 0);
        \draw (0, -6) -- (0, 6);

        \foreach \i in {-6,...,6} {
            \foreach \j in {-6,...,6} {
                \fill (\i*\voneX + \j*\vtwoX, \i*\voneY + \j*\vtwoY) circle (2pt);
            }
        }

        \draw[->, thick, black] (0,0) -- (\voneX, \voneY) node[right, yshift=-2pt] {$\mathbf{b}_1$};
        \draw[->, thick, black] (0,0) -- (\vtwoX, \vtwoY) node[right, yshift=2pt] {$\mathbf{b}_2$};
    \end{tikzpicture}
    \caption{Representación bidimensional de un retículo generado por los vectores $\mathbf{b}_1$ y $\mathbf{b}_2$.}
    \label{fig:reticulo_2d}
\end{figure}

Como se ilustra en la Figura \ref{fig:reticulo_2d}, en un espacio bidimensional la estructura del retículo resulta visualmente intuitiva. Con baja dimensionalidad, identificar el vector más corto desde el origen o encontrar el punto del retículo más cercano a una coordenada arbitraria es una tarea sencilla.

La seguridad de la criptografía basada en retículos radica en que, en dimensiones altas (por ejemplo, $n > 500$), encontrar el vector no nulo más corto \textit{Shortest Vector Problem}, SVP o el vector más cercano a un punto dado (\textit{Closest Vector Problem}, CVP) es un problema intratable, incluso para un ordenador cuántico   

Sin embargo, ML-DSA no utiliza retículos estándar, sino \textbf{retículos definidos sobre módulos}, donde un módulo es similar a un espacio vectorial pero definido sobre un anillo en vez de un cuerpo. En lugar de utilizar matrices gigantes de números enteros, lo cual requeriría tamaños de clave inasumibles en la práctica, el estándar opera con matrices y vectores cuyos elementos son los polinomios del anillo $R_q$. 

Si definimos dimensiones $k$ y $l$, que varían según el nivel de seguridad del FIPS 204, la matriz pública $\mathbf{A}$ pertenece a $R_q^{k \times l}$, y los vectores secretos $\mathbf{s}_1$ y $\mathbf{s}_2$ pertenecen a $R_q^l$ y $R_q^k$ respectivamente, siendo vectores de polinomios de longitud $l$ y $k$. Esta transición de escalares enteros a polinomios es la decisión de diseño que permite reducir drásticamente el tamaño en memoria de las claves públicas y privadas, manteniendo la complejidad del problema.

\subsection{La Norma de Infinito ($\| \cdot \|_\infty$)}

En el contexto de la ingeniería de ML-DSA, no basta con generar vectores; es necesario medir su magnitud para garantizar que los valores generados durante la firma sean lo suficientemente pequeños como para no revelar información estadística sobre la clave secreta. 

Para ello, la criptografía basada en retículos no utiliza la norma euclidiana convencional, sino la \textbf{norma de infinito} ($\| \cdot \|_\infty$). 

La norma de infinito de un polinomio $w \in R_q$ se define como el valor absoluto máximo de sus coeficientes, tras aplicar la reducción modular simétrica descrita en la Sección \ref{subsec:reduccion_simetrica}:

$$ \| w \|_\infty = \max_{0 \le i < 256} |w_i| $$

La norma de infinito de un vector de polinomios $\mathbf{w}$ es el máximo de las normas de infinito de todos los polinomios que lo componen. 

\begin{ejemplo}[Cálculo de la Norma Infinito en Polinomios y Vectores]
\label{ex:norma_infinito}
Para ilustrar este cálculo, supongamos que operamos en el anillo polinómico $\mathbb{Z}_{13}[x]$. Definimos un polinomio inicial $w(x) = 3 + 11x + 4x^2$. 

Antes de calcular la norma infinito ($||\cdot||_\infty$), el esquema necesita que los coeficientes estén expresados en su representación modular simétrica, en este caso, proyectados al rango $[-6, 6]$. Aplicando esta reducción al coeficiente $11$ ($11 - 13 = -2$), obtenemos el polinomio equivalente:
$$w(x) = 3 - 2x + 4x^2$$

Sus coeficientes simétricos son $(3, -2, 4)$ y sus valores absolutos correspondientes son $(3, 2, 4)$. Por lo tanto, la norma de infinito del polinomio es el coeficiente absoluto máximo de este conjunto:
$$ \| w \|_\infty = \max(|3|, |-2|, |4|) = 4 $$

Si escalamos esta lógica a un vector polinómico bidimensional $\mathbf{v} = (v_1, v_2)$ definido también sobre $\mathbb{Z}_{13}[x]$, donde $v_1(x) = 2 + 5x$ y $v_2(x) = -8 + x$ (cuya representación simétrica es $5 + x$):
\begin{itemize}
    \item $\| v_1 \|_\infty = \max(|2|, |5|) = 5$
    \item $\| v_2 \|_\infty = \max(|5|, |1|) = 5$
\end{itemize}

La norma de infinito del vector completo será, por definición, el máximo de las normas individuales de sus componentes: 
$$ \| \mathbf{v} \|_\infty = \max(\| v_1 \|_\infty, \| v_2 \|_\infty) = \max(5, 5) = 5 $$
\end{ejemplo}

Desde la perspectiva de la implementación en \textit{software}, la norma de infinito es la métrica de control de flujo fundamental. Permite establecer límites de seguridad estrictos: si durante el proceso de generación, un vector resultante obtiene una norma de infinito que supera un umbral límite predefinido (por ejemplo, si el límite fuera $6$, el vector $\mathbf{v}$ anterior con norma $8$ sería inválido), el algoritmo aborta la operación y reinicia el muestreo. Esta es la base matemática del mecanismo de "muestreo de rechazo", garantizando que las firmas finalmente publicadas, tanto para que la validación de la firma funcione correctamente, como para que nunca se filtren vectores susceptibles de criptoanálisis \cite{nist_fips204}.


\section{Mecanismos de Muestreo y Seguridad}

La seguridad de ML-DSA no reside únicamente en la complejidad de sus estructuras, sino en la imposibilidad computacional de revertir las operaciones que generan las claves y las firmas. Este blindaje se articula mediante el problema MLWE y el uso estratégico del muestreo de rechazo.

\subsection{El problema \textit{Module Learning With Errors} MLWE}

El núcleo de la seguridad de Dilithium es el problema de Aprendizaje con Errores sobre Módulos MLWE. En un esquema de álgebra lineal simple sobre un cuerpo finito, resolver un sistema de ecuaciones de la forma $\mathbf{t} = \mathbf{A}\mathbf{s}$ es trivial mediante la eliminación gaussiana. Sin embargo, MLWE introduce una pequeña perturbación o error que rompe esta linealidad.

La clave pública en ML-DSA se define mediante la ecuación fundamental:
\begin{equation}
    \mathbf{t} = \mathbf{A}\mathbf{s}_1 + \mathbf{s}_2 \pmod q
\end{equation}

Donde:
\begin{itemize}
    \item $\mathbf{A} \in R_q^{k \times l}$ es una matriz pública de polinomios generada de forma pseudoaleatoria.
    \item $\mathbf{s}_1, \mathbf{s}_2$ son los componentes de la clave privada, consistentes en vectores de polinomios con coeficientes muy pequeños.
    \item $\mathbf{t}$ es el vector resultante que conforma la clave pública.
\end{itemize}

Para un atacante, recuperar $\mathbf{s}_1$ y $\mathbf{s}_2$ a partir de $\mathbf{A}$ y $\mathbf{t}$ es equivalente a resolver el problema del vector más cercano CVP en un retículo de alta dimensionalidad, una tarea que se considera intratable incluso para computación cuántica a gran escala \cite{nist_fips204}.

\subsection{Muestreo de Rechazo y el Paradigma \textit{Fiat-Shamir with Aborts}}

Uno de los mayores riesgos en la criptografía de retículos es que la firma digital pueda dibujar involuntariamente la forma de la clave privada. Si el vector de la firma $\mathbf{z}$, cuya función se hablará mas adelante, está demasiado correlacionado con el secreto $\mathbf{s}_1$, un atacante podría acumular suficientes firmas para realizar un análisis estadístico y recuperar la clave.

Para evitar esta fuga de información, ML-DSA implementa el paradigma de \textit{Fiat-Shamir con Abortos}\cite{lyubashevsky2009fiat}. Durante el proceso de firma, el algoritmo genera un candidato a firma y verifica su norma de infinito ($\|\cdot\|_\infty$). Si la firma es demasiado grande o está fuera de un rango de seguridad predefinido, el algoritmo la descarta y reinicia el proceso con nuevos valores aleatorios.

Este mecanismo garantiza que la distribución estadística de las firmas publicadas sea completamente independiente de la clave secreta. Desde el punto de vista de la ingeniería de \textit{software}, esto implica que el tiempo de ejecución del algoritmo de firma es variable, ya que depende del número de iteraciones necesarias hasta encontrar una firma que cumpla los criterios de seguridad del estándar.

\subsection{Problema \textit{Module Short Integer Solution} MSIS}

Mientras que el problema MLWE garantiza la ocultación de la clave privada, la resistencia a la falsificación de firmas en ML-DSA recae sobre un segundo pilar matemático: el problema de la Solución de Enteros Cortos sobre Módulos MSIS.

En el esquema de firma, un atacante que no posea la clave secreta pero desee generar una firma válida $(\mathbf{z}, \mathbf{c})$ para un mensaje fraudulento, se enfrenta al reto de validar la ecuación de verificación del algoritmo. Matemáticamente, esto equivale a encontrar un vector $\mathbf{z}$ no nulo tal que:

$$ \mathbf{A}\mathbf{z} - \mathbf{t}c = \mathbf{w}_1 \pmod q $$

Es decir, que se cumpla con la restricción fundamental de tamaño establecido por el estándar:

$$ \|\mathbf{z}\|_\infty < \gamma_1 - \beta $$

Entonces el problema MSIS se define como: dada una matriz uniforme $\mathbf{A}$, encontrar un vector cuyas componentes sean cortas, es decir, con una norma de infinito muy pequeña, y que resuelva la ecuación matricial en el anillo $R_q$, es un problema asintóticamente equivalente a encontrar el vector más corto en un retículo SVP. 

Por lo tanto, la seguridad de Dilithium es bidireccional: la imposibilidad de resolver MLWE protege el secreto, y la intratabilidad de MSIS impide la suplantación de identidad mediante la falsificación de la firma \cite{peikert2016decade}.


\chapter{El Esquema de Firma ML-DSA (FIPS 204)}

Tras establecer los fundamentos matemáticos que rigen la criptografía basada en retículos, este capítulo desglosa la arquitectura interna del Algoritmo de Firma Digital Basado en Retículos Modulares (ML-DSA), establecido por el NIST en la publicación FIPS 204. Se detallarán sus conjuntos de parámetros y los tres procesos algorítmicos principales que lo componen: generación de claves, firma y verificación.

\section{Introducción al Estándar y Parámetros de Seguridad}

FIPS 204 no es un algoritmo cerrado, se trata de un conjunto de esquemas que se puede ajustar. La idea principal es que se pueda graduar la seguridad según nuestras necesidades, buscando el equilibrio ideal entre protección, velocidad y consumo de recursos.

\subsection{Niveles de Seguridad}

El NIST ha definido tres conjuntos de parámetros oficiales, cada uno asociado a una categoría de seguridad específica, comparable con la dificultad de romper algoritmos clásicos mediante ataques de fuerza bruta:

\begin{itemize}
    \item \textbf{ML-DSA-44 (Categoría 2):} Seguridad equivalente a la dificultad de romper AES-128 mediante fuerza bruta cuántica. Es el perfil más ligero y eficiente.
    \item \textbf{ML-DSA-65 (Categoría 3):} Seguridad equivalente a romper AES-192.
    \item \textbf{ML-DSA-87 (Categoría 5):} Seguridad equivalente a romper AES-256. Es el perfil de máxima resistencia contemplado por el estándar.
\end{itemize}

\subsection{Notaciones del Sistema}

Independientemente del nivel de seguridad elegido, la aritmética fundamental del esquema opera sobre constantes estructurales inmutables \cite{nist_fips204}. El anillo de polinomios $R_q = \mathbb{Z}_q[x]/(x^{256} + 1)$ se define siempre en $n = 256$ y un primo $q =$ $2^{23} - 2^{13} + 1$ = 8380417. Asimismo, el número de bits descartados en las operaciones de compresión de claves públicas es constante ($d = 13$).

El resto de los parámetros varían según el nivel de seguridad, definiendo la entropía de los secretos y los umbrales de rechazo. A continuación se enumera cada uno de los parámetros y su significado:

\begin{itemize}
    \item $k$ y $l$: Dimensiones del retículo subyacente. Determinan el tamaño de la matriz pública $\mathbf{A} \in R_q^{k \times l}$, incrementando la dureza del problema MLWE a medida que crecen.
    \item $\eta$: Rango de extracción para los coeficientes de los polinomios que definen las claves privadas $\mathbf{s}_1$ y $\mathbf{s}_2$. Se seleccionan de forma aleatoria en el intervalo $[-\eta, \eta]$.
    \item $\tau$: Peso de Hamming específico del polinomio de desafío $c$. Define cuántos coeficientes del polinomio (sobre los 256 totales) son estrictamente $+1$ o $-1$, siendo el resto $0$.
    \item $\gamma_1$: Este parámetro define el intervalo $[-\gamma_1 + 1, \gamma_1]$ a partir del cual se define de forma aleatoria la máscara $\mathbf{y}$.
    \item $\gamma_2$: Parámetro de descomposición. Define el tamaño de la ventana utilizada para separar los bits más significativos (\textit{HighBits}) de los menos significativos (\textit{LowBits}) del compromiso.
    \item $\omega$: Límite máximo de bits a 1 permitidos en el vector del \textit{hint} $\mathbf{h}$.
\end{itemize}

En la Tabla \ref{tab:parametros_mldsa} se detalla la configuración exacta de estos parámetros para cada nivel de seguridad, junto con el impacto que tienen en el tamaño final de la clave pública, clave privada y firma digital.

\vspace{0.5cm}
\begin{table}[H]
\centering
\renewcommand{\arraystretch}{1.3}
\begin{tabular}{|l|c|c|c|}
\hline
\rowcolor[HTML]{EFEFEF} 
\textbf{Parámetro} & \textbf{ML-DSA-44} & \textbf{ML-DSA-65} & \textbf{ML-DSA-87} \\ \hline
Categoría NIST & Nivel 2 & Nivel 3 & Nivel 5 \\ \hline
$k \times l$ (Dimensión de $\mathbf{A}$) & $4 \times 4$ & $6 \times 5$ & $8 \times 7$ \\ \hline
$\eta$ (Rango de secreto) & 2 & 4 & 2 \\ \hline
$\tau$ (Peso del desafío) & 39 & 49 & 60 \\ \hline
$\gamma_1$ (Límite de máscara) & $2^{17}$ & $2^{19}$ & $2^{19}$ \\ \hline
$\gamma_2$ (Ventana \textit{HighBits}) & $(q-1)/88$ & $(q-1)/32$ & $(q-1)/32$ \\ \hline
$\omega$ (Peso máximo del \textit{hint}) & 80 & 55 & 75 \\ \hline
Tamaño Clave Pública (\textit{pk}) & 1.312 bytes & 1.952 bytes & 2.592 bytes \\ \hline
Tamaño Clave Privada (\textit{sk}) & 2.560 bytes & 4.032 bytes & 4.896 bytes \\ \hline
Tamaño de la Firma & 2.420 bytes & 3.309 bytes & 4.627 bytes \\ \hline
\end{tabular}
\caption{Conjuntos de parámetros oficiales y tamaños de salida.}
\label{tab:parametros_mldsa}
\end{table}
\vspace{0.5cm}

\section{Generación del Par de Claves}

El proceso de generación de claves en ML-DSA tiene como objetivo producir un par de claves asimétricas: una clave pública ($pk$) utilizada por los verificadores para validar la autenticidad, y una clave secreta ($sk$) utilizada por el firmante. A diferencia de esquemas clásicos como RSA, donde se generan de forma aleatoria dos primos, en Dilithium el proceso consiste en generar de forma aleatoria una matriz y vectores con elementos en el anillo $R_q$.


\subsection{Expansión de semillas criptográficas ($\rho$, $\rho'$, $K$)}

Para minimizar el tamaño de las claves públicas, el estándar FIPS 204 genera las matrices y vectores directamente de un generador de números pseudoaleatorios. Posteriormente se hará pública esa semilla en vez de la matriz $\mathbf{A}$.

El algoritmo comienza generando una semilla aleatoria uniforme de 256 bits (32 bytes), denotada como $\xi$. A continuación, esta semilla raíz se somete a un proceso de expansión utilizando la función de salida extensible (XOF) \texttt{SHAKE-256}, obteniéndose exactamente 128 bytes de salida.

La ecuación formal de este proceso, correspondiente al primer paso del algoritmo, se define como:
\begin{equation}
    (\rho, \rho', K) \leftarrow \text{SHAKE-256}(\xi, 128)
\end{equation}

El flujo de bytes resultante se segmenta estrictamente en tres nuevas semillas con propósitos criptográficos independientes:
\begin{itemize}
    \item $\rho$: Los primeros 32 bytes. Se utiliza como semilla pública para generar de forma aleatoria la matriz $\mathbf{A}$. Al ser derivada de $\xi$, permite que la matriz no tenga que almacenarse completa en memoria, sino que pueda regenerarse cuando sea necesaria.
    \item $\rho'$: Los siguientes 64 bytes. Se emplea como semilla secreta para generar los vectores de error privados $\mathbf{s}_1$ y $\mathbf{s}_2$. Esto quiere decir que $\rho'$ se mantiene en secreto.
    \item $K$: Los últimos 32 bytes. Constituye una semilla secreta adicional que se almacenará en la clave privada $sk$ y se utilizará posteriormente durante el proceso de firma para derivar el vector de enmascaramiento aleatorio $\mathbf{y}$.

\end{itemize}
La Figura \ref{fig:keygen_expansion} muestra este flujo . Como se observa, el uso de la función de salida extensible SHAKE-256 garantiza que, a partir de una única fuente de entropía inicial $\xi$, el sistema extraiga de forma estructurada las tres semillas fundamentales sin comprometer la seguridad del material criptográfico.

\begin{figure}[H]
    \centering
    \begin{tikzpicture}[
        node distance = 1.5cm and 2cm,
        caja/.style = {draw, rectangle, rounded corners, minimum width=3.5cm, minimum height=1cm, align=center, fill=white, thick},
        caja_izq/.style = {draw, rectangle, rounded corners, minimum width=5cm, align=left, fill=white, thick, inner sep=10pt},
        proceso/.style = {draw, ellipse, minimum width=3cm, minimum height=1cm, align=center, fill=white, thick},
        flecha/.style = {-{Stealth[scale=1.2]}, thick},
        linea/.style = {thick}
    ]

    \node[caja] (semilla) {$\xi$ (256 bits aleatorios)};
    \node[proceso, below=of semilla] (shake) {\texttt{SHAKE-256}};
    
    \node[caja_izq, below=of shake] (salida) {
        \textbf{$\rho$} (32 bytes): Para generar la matriz pública $\mathbf{A}$\\
        \textbf{$\rho'$} (64 bytes): Para generar los secretos $\mathbf{s}_1$ y $\mathbf{s}_2$\\
        \textbf{$K$} (32 bytes): Semilla secreta para firmar
    };

    \node[caja, dashed, above left=0.5cm and 1cm of shake] (ecuacion) {$(\rho, \rho', K) \leftarrow \text{SHAKE-256}(\xi, 128)$};

    \draw[flecha] (semilla) -- (shake);
    \draw[flecha] (shake) -- (salida);
    
    \draw[flecha] (shake.west) -| (ecuacion.south);

    \end{tikzpicture}
    \caption{Expansión determinista de la semilla raíz $\xi$.}
    \label{fig:keygen_expansion}
\end{figure}

\subsection{Expansión Matricial $\mathbf{A}$ y Vectores Secretos $ \mathbf{s}_1,\mathbf{s}_2$}

Una vez obtenidas las semillas base, el algoritmo procede a instanciar las estructuras algebraicas sobre el anillo $R_q$. FIPS 204 especifica diferentes métodos y separadores de dominio para garantizar que las expansiones no colisionen.

Para la \textbf{matriz pública $\mathbf{A}$}, se utiliza la semilla $\rho$. Debido al gran volumen de datos necesario para generar una matriz de dimensiones $k \times l$ (donde cada celda es un polinomio de 256 coeficientes), se emplea \texttt{SHAKE-128}, al ser más rápido que su homólogo de 256 bits. La generación de cada coeficiente del polinomio $\mathbf{A}_{i,j}$ se realiza inyectando las coordenadas en el \textit{hash}:
$$ \text{SHAKE-128}(\rho \parallel j \parallel i) $$
El flujo de \textit{bits} resultante se procesa en bloques de 24 \textit{bits}. Si el entero extraído es estrictamente menor que el módulo $q = 8380417$, se acepta como coeficiente válido; de lo contrario, se descarta. 

En paralelo, se generan los \textbf{vectores secretos $\mathbf{s}_1$ y $\mathbf{s}_2$} utilizando la semilla privada $\rho'$ mediante \texttt{SHAKE-256}. En este caso, se utiliza un contador iterativo $k$ para garantizar que cada uno de los polinomios resultantes sea criptográficamente independiente. A diferencia de la matriz $\mathbf{A}$, los coeficientes de los vectores secretos no pueden tomar valores hasta $q$, sino que deben ser cortos. Se aplica un filtro de rechazo estricto para asegurar que cada coeficiente pertenezca al intervalo cerrado $[-\eta, \eta]$, dependiendo este límite del nivel de seguridad configurado.

\subsection{Ecuación MLWE, Compresión y Empaquetado Final}
\label{subsubsec:mlwe_empaquetado}

Con los componentes base, el sistema ejecuta la operación central: el problema de Aprendizaje con Errores en Módulos. Para que la computación sea eficiente y matemáticamente válida, todos los elementos deben operar en el mismo dominio.

La matriz pública $A$, que se expande a partir de la semilla $\rho$, se genera directamente en el dominio de NTT, denotándose como $\hat{A}$. De esta manera se evita el alto coste de transformar la matriz completa tras su generación. Por tanto, el sistema solo necesita aplicar la transformada a los vectores secretos $s_1$ y $s_2$ para realizar la operación:

\[ \hat{t} = \hat{A} \cdot \text{NTT}(s_1) + \text{NTT}(s_2) \]

Donde $\hat{t}$ representa el vector público resultante en el dominio NTT. Para proceder a la fase de compresión, este vector debe regresar al dominio polinómico normal mediante la transformada inversa:

\[ t = \text{NTT}^{-1}(\hat{t}) \]

El vector resultante $t$ contiene la información pública generada, pero su tamaño crudo resulta ineficiente para su transmisión por red y almacenamiento. Para solucionar este cuello de botella, el estándar FIPS 204 aplica el algoritmo de compresión \textit{Power2Round}, que divide matemáticamente cada coeficiente del vector descartando sistemáticamente $d = 13$ bits.

Como resultado de esta operación, el vector $t$ se descompone en dos partes: los bits más significativos $t_1$ y los bits menos significativos $t_0$. Finalmente, se empaqueta exclusivamente el componente de orden alto $t_1$ junto con la semilla $\rho$, conformando así la clave pública definitiva que será distribuida en el sistema.

\begin{ejemplo}
\label{ex:power2round}
A modo de ejemplo, si tras operar el polinomio se obtiene un coeficiente $r = 1.234.567$, el algoritmo lo fragmenta utilizando la potencia $2^d$ (donde $2^{13} = 8192$). Este proceso produce un cociente y un resto que representan los bits de peso alto y bajo respectivamente:
\begin{itemize}
    \item \textbf{Bits más significativos $t_1$:} Representa los bits significativos. En este caso, $t_1 = \lfloor 1.234.567 / 8192 \rfloor = 150$. Constituye la información relevante que se publicará en la clave pública.
    \item \textbf{Bits menos significativos $t_0$:} Representa el residuo exacto ($t_0 = r \pmod{2^{13}} = 5767$). Este valor actúa como el error de redondeo que se mantiene oculto en la clave secreta.
\end{itemize}
Se puede comprobar que $1234567 = 150 \cdot 2^{13} + 5767$.
\end{ejemplo}

Como se observa en el \ref{ex:power2round}, la separación permite reducir drásticamente el peso de la clave pública sin perder la capacidad de reconstrucción necesaria para la verificación.

Finalmente, el sistema codifica estas estructuras en \textit{bytes} (\texttt{pkEncode} y \texttt{skEncode}) y empaqueta las claves:
\begin{enumerate}
    \item \textbf{Clave Pública $pk$:} Se forma mediante la concatenación de la semilla de la matriz y el vector comprimido: $pk = (\rho, \mathbf{t}_1)$.
    \item \textbf{Clave Privada $sk$:} Recopila todos los elementos necesarios para la posterior firma. Incluye las semillas $\rho', K$, los vectores secretos $\mathbf{s}_1, \mathbf{s}_2$, el residuo de compresión $\mathbf{t}_0$ y un elemento de vinculación criptográfica denominado $tr$. Este último se calcula como el \textit{hash} de la clave pública: $tr = \text{SHAKE-256}(\rho \parallel \mathbf{t}_1, 64)$, asegurando una salida constante de 64 \textit{bytes} que ata el secreto a su identidad pública.
\end{enumerate}

\begin{figure}[H]
    \centering
    \resizebox{0.95\textwidth}{!}{
    \begin{tikzpicture}[
        node distance = 1.2cm and 1.5cm,
        caja/.style = {draw, rectangle, rounded corners, minimum width=3.5cm, minimum height=1cm, align=center, fill=white, thick},
        caja_clara/.style = {draw, rectangle, rounded corners, minimum width=3cm, minimum height=0.8cm, align=center, fill=gray!5, thick},
        proceso/.style = {draw, ellipse, minimum width=3cm, minimum height=1cm, align=center, fill=white, thick},
        flecha/.style = {-{Stealth[scale=1.2]}, thick},
        linea/.style = {thick}
    ]

    \node[caja_clara] (rho) {Semilla $\rho$ (32 bytes)};
    \node[caja_clara, right=4cm of rho] (rhoprima) {Semilla $\rho'$ (64 bytes)};

    \node[proceso, below=0.8cm of rho] (expA) {ExpandA (\texttt{SHAKE-128})};
    \node[proceso, below=0.8cm of rhoprima] (expS) {ExpandS (\texttt{SHAKE-256})};

    \node[caja, below=0.8cm of expA] (matA) {Matriz Pública $\mathbf{A} \in R_q^{k \times l}$};
    \node[caja, below=0.8cm of expS] (vecS) {Secretos $\mathbf{s}_1, \mathbf{s}_2 \in S_\eta$};

    \node[caja, dashed, below right=1.2cm and -2cm of matA] (mlwe) {
        \textbf{Cálculo de la clave pública}\\
        $t = \text{NTT}^{-1}\left( \hat{A} \cdot \text{NTT}(\mathbf{s}_1) + \text{NTT}(\mathbf{s}_2) \right)$
    };

    \node[proceso, below=1cm of mlwe] (p2r) {\texttt{Power2Round} ($d=13$)};

    \node[caja_clara, below left=1cm and 0.5cm of p2r] (t1) {Bits más significativos $\mathbf{t}_1$};
    \node[caja_clara, below right=1cm and 0.5cm of p2r] (t0) {Bits menos significativos $\mathbf{t}_0$};

    \node[caja, thick, below=0.8cm of t1] (pk) {\textbf{Clave Pública $pk$}\\$(\rho, \mathbf{t}_1)$};
    \node[proceso, right=0.5cm of pk] (hash_tr) {\texttt{SHAKE-256}};
    \node[caja, thick, below=2cm of t0] (sk) {\textbf{Clave Privada $sk$}\\$(\rho', K, tr, \mathbf{s}_1, \mathbf{s}_2, \mathbf{t}_0)$};

    \draw[flecha] (rho) -- (expA);
    \draw[flecha] (rhoprima) -- (expS);
    
    \draw[flecha] (expA) -- (matA);
    \draw[flecha] (expS) -- (vecS);
    
    \draw[flecha] (matA.south) |- (mlwe.west);
    \draw[flecha] (vecS.south) |- (mlwe.east);
    
    \draw[flecha] (mlwe) -- (p2r);
    
    \draw[flecha] (p2r.south) -| (t1.north);
    \draw[flecha] (p2r.south) -| (t0.north);
    
    \draw[flecha] (t1) -- (pk);
    
    \draw[flecha] (pk) -- (hash_tr);
    \draw[flecha] (hash_tr) |- (sk.west) node[near start, right] {$tr$};
    
    \draw[flecha] (t0) -- (sk);
    
    \end{tikzpicture}
    }
    \caption{Generación de las claves a partir de las semillas.}
    \label{fig:keygen_estructuras}
\end{figure}

En la Figura \ref{fig:keygen_estructuras} es importante analizar la dimensión de las estructuras de datos generadas, ya que afectan directamente la huella en memoria del sistema. 

Aunque conceptualmente $\mathbf{s}_1$ y $\mathbf{s}_2$ se definen como vectores, a nivel de implementación son \textit{arrays} multidimensionales. Cada posición del vector no contiene un número escalar, sino un polinomio completo del anillo $R_q$, es decir, una lista de $n = 256$ coeficientes enteros. El tamaño y el contenido de estos vectores están estrictamente acotados por el nivel de seguridad seleccionado:

\begin{itemize}
    \item \textbf{Longitud del vector:} El vector $\mathbf{s}_1$ tiene siempre longitud $l$, mientras que $\mathbf{s}_2$ tiene longitud $k$.
    \item \textbf{Rango de los coeficientes $\eta$:} Cada uno de los 256 coeficientes de cada polinomio debe pertenecer al intervalo $[-\eta, \eta]$.
\end{itemize}

Para materializar estas restricciones, tómense como ejemplo los extremos del estándar:
\begin{itemize}
    \item \textbf{En el nivel mínimo de seguridad (ML-DSA-44):} Los parámetros son $k=4$, $l=4$ y $\eta=2$. Por tanto, el vector $\mathbf{s}_1$ está compuesto por 4 polinomios, sumando un total de $1.024$ coeficientes. Debido al filtro de rechazo, el valor de todos y cada uno de estos 1.024 coeficientes solo puede pertenecer al intervalo $[-2, 2]$, es decir, los coeficientes de los polinomios en $\mathbf{s}_1$ están acotados por $\eta = 2$.
    \item \textbf{En el nivel intermedio (ML-DSA-65):} Las dimensiones crecen a $k=6$ y $l=5$. Aquí $\mathbf{s}_1$ contiene $1.280$ coeficientes. Al ser $\eta=4$, el rango de valores permitidos se amplía al intervalo $[-4, 4]$, es decir, polinomios con coeficientes acotados por $\eta = 4$.
    \item \textbf{En el nivel máximo de seguridad (ML-DSA-87):} El esquema escala a $k=8$ y $l=7$, generando un $\mathbf{s}_1$ de $1.792$ coeficientes. La cota $\eta$ vuelve a reducirse ($\eta=2$), por lo que el rango vuelve a ser del intervalo $[-2, 2]$. La seguridad en este nivel no reside en el tamaño del secreto, sino en la gran dimensión.
\end{itemize}

Esta variabilidad demuestra por qué el empaquetado final (\texttt{skEncode}) debe ser altamente eficiente, comprimiendo estos miles de coeficientes para que la clave privada pueda ser funcional en un entorno de red actual, ya que de otra manera no es viable.

\subsubsection{Compresión de la Clave Pública: El algoritmo \texttt{Power2Round}}

El vector $t$ resultante de la ecuación central MLWE, definido matemáticamente como:
$$t = \text{NTT}^{-1}(\hat{A} \cdot \text{NTT}(s_1) + \text{NTT}(s_2))$$
es demasiado grande en su formato crudo, lo que representa una huella de memoria ineficiente para su transmisión a través de la red como clave pública. Para solucionar esto, el estándar FIPS 204 (ML-DSA) no almacena $t$ íntegramente, sino que lo comprime utilizando el algoritmo \texttt{Power2Round}.

El objetivo de esta función es dividir cada coeficiente $r$ del vector $\mathbf{t}$ en dos componentes: los \textit{bits} más significativos $r_1$, que conformarán la clave pública $\mathbf{t}_1$; y los \textit{bits} menos significativos $r_0$, que se almacenarán en la clave privada $\mathbf{t}_0$ como residuo necesario para la futura firma.

La partición se realiza en base a una potencia de dos $2^d$, donde $d = 13$ para todos los niveles de seguridad de ML-DSA. Matemáticamente, se busca que $r = r_1 \cdot 2^d + r_0$. 

Para evitar errores y que pueda afectar para más adelante, la extracción de $r_0$ no emplea un truncamiento de \textit{bits} clásico, sino que utiliza una reducción modular simétrica. El algoritmo procesa cada coeficiente de la siguiente manera:

\begin{enumerate}
    \item \textbf{Cálculo del residuo $r_0$:} Se extraen los $d$ \textit{bits} menos significativos proyectándolos en el rango simétrico $[-2^{d-1}, 2^{d-1}]$. Para $d=13$, esto significa que $r_0$ siempre estará acotado en el intervalo $[-4096, 4096]$.
    \item \textbf{Cálculo de la componente más significativa $r_1$:} Una vez obtenido el residuo exacto, se sustrae del valor original y se divide entre $2^d$:
    \begin{equation}
        r_1 = \frac{r - r_0}{2^d}
    \end{equation}
\end{enumerate}

\vspace{0.5cm}
\begin{ejemplo}[Implementación Numérica]
\label{ex:compresion_numerica}
Consideremos un coeficiente con valor $r = 30.000$ y el parámetro de truncamiento $d = 13$ (donde $2^{13} = 8192$). El proceso de división y redondeo se ejecuta de la siguiente manera:
\begin{itemize}
    \item Se calcula la división para aislar los bits más significativos, redondeando al entero más cercano: $30.000 / 8192 \approx 3,66$. El entero más próximo es $4$.
    \item Por tanto, la componente alta asume el valor $r_1 = 4$.
    \item El residuo centrado (la componente baja) se despeja de la ecuación: 
    \[ r_0 = r - (r_1 \cdot 2^d) = 30.000 - (4 \cdot 8192) = 30.000 - 32.768 = -2768 \]
\end{itemize}

De esta forma, el coeficiente original $30.000$ queda comprimido matemáticamente en el par $(r_1=4, r_0=-2768)$. El valor $4$ requiere exponencialmente menos \textit{bits} , logrando una reducción drástica del tamaño final de la clave pública.
\end{ejemplo}

Para mostrar la mecánica descrita en el Ejemplo \ref{ex:compresion_numerica}, la Figura \ref{fig:power2round} esquematiza el flujo de operaciones del algoritmo \texttt{Power2Round}. Como se observa en el diagrama, la arquitectura garantiza una separación estricta tras la sustracción: la componente más significativa $r_1$ fluye hacia la estructura de la clave pública $\mathbf{t}_1$, mientras que el residuo centrado de la división $r_0$ queda aislado, pasando a formar parte del material criptográfico de la clave privada $\mathbf{t}_0$.

\begin{figure}[H]
    \centering
    \begin{tikzpicture}[
        caja/.style = {draw, rectangle, rounded corners, minimum height=0.8cm, inner sep=0.2cm, align=center, fill=white, thick},
        caja_pub/.style = {draw, rectangle, rounded corners, minimum height=0.8cm, inner sep=0.2cm, align=center, fill=blue!5, thick},
        caja_priv/.style = {draw, rectangle, rounded corners, minimum height=0.8cm, inner sep=0.2cm, align=center, fill=red!5, thick},
        flecha/.style = {-{Stealth[scale=1.2]}, thick}
    ]

    \node[caja]      (r)       at (0, 0)    {Coeficiente $r \in \mathbf{t}$};
    \node[caja]      (calc_r0) at (4, -1.5) {$r_0 = r \bmod^{\pm} 2^d$};
    \node[caja]      (calc_r1) at (0, -3.5) {$r_1 = (r - r_0)/2^d$};
    \node[caja_priv] (t0)      at (4, -3.5) {$\mathbf{t}_0$ (Clave Privada)};
    \node[caja_pub]  (t1)      at (0, -5.5) {$\mathbf{t}_1$ (Clave Pública)};

    \draw[flecha] (r.south) -- ++(0,-0.5) -| (calc_r0.north);

    \draw[flecha] (calc_r0.south west) -- (calc_r1.north east) 
        node[midway, left=0.2cm, align=right] {Sustracción};

    \draw[flecha] (calc_r0.south) -- (t0.north) 
        node[midway, left=0.1cm] {$r_0$};

    \draw[flecha] (calc_r1.south) -- (t1.north) 
        node[midway, right=0.1cm] {$r_1$};

    \end{tikzpicture}
    \caption{Descomposición de un coeficiente $r$ mediante \texttt{Power2Round}.}
    \label{fig:power2round}
\end{figure}

\subsection{Aritmética y Dominio NTT}

La traducción de las ecuaciones algebraicas de ML-DSA a código fuente introduce desafíos que deben gestionarse de manera que se  garantiza el rendimiento del esquema. Destacan tres consideraciones fundamentales:

\begin{enumerate}
    \item \textbf{Generación y Muestreo de Rechazo:} Durante la instanciación de vectores pseudoaleatorios, como $\mathbf{A}$, $\mathbf{s}_1$ o $\mathbf{s}_2$, la salida del \textit{hash} XOF rara vez coincide perfectamente con el rango deseado. El sistema implementa un bucle condicional: si el coeficiente generado excede el módulo $q$ (en el caso de $\mathbf{A}$) o el límite $\eta$ (en el caso de los secretos), se rechaza y se extraen nuevos \textit{bits} del XOF. No se aplican operaciones de módulo en esta fase para evitar sesgos estadísticos en la distribución uniforme.
    
    \item \textbf{Aritmética Modular en Operaciones:} Una vez que los polinomios están establecidos, cualquier suma, resta o multiplicación entre coeficientes puede provocar desbordamientos numéricos, fenómeno también conocido como \textit{overflow}. Por ello, toda operación aritmética se envuelve en una reducción modular. Dependiendo de la fase del algoritmo, esta reducción puede ser estándar, resultados entre $0$ y $q-1$, o centrada, que serán resultados acotados en el intervalo absoluto $[-q/2, q/2]$.
    
    \item \textbf{El Estado del Dominio NTT en Memoria:} La multiplicación de dos vectores de polinomios de 256 coeficientes tiene una complejidad cuadrática $\mathcal{O}(n^2)$. Para optimizar el rendimiento, se emplea la NTT, que reduce la complejidad a $\mathcal{O}(n \log n)$. \textbf{Para multiplicar dos polinomios ambos deben transformarse al dominio NTT}. Una vez multiplicados coeficiente a coeficiente, se aplica la transformada inversa $\text{NTT}^{-1}$ para devolver el resultado al dominio original y poder realizar operaciones (\texttt{Power2Round}). 
\end{enumerate}

\section{Generación de la Firma Digital}

El proceso de firma es el núcleo operativo del esquema ML-DSA y, a diferencia de la generación de claves, no es de naturaleza estrictamente lineal. Su seguridad descansa sobre el paradigma \textit{Fiat-Shamir with Aborts}, lo que implica la ejecución de un bucle condicional (muestreo de rechazo) que puede iterar múltiples veces hasta generar una firma que no filtre información sobre los vectores privados.

\subsection{Preparación del Mensaje y Vinculación Criptográfica $\mu$}

Antes de iniciar la aritmética de polinomios y el bucle de enmascaramiento, el algoritmo debe acondicionar la información de entrada. El objetivo de esta fase inicial es vincular el mensaje original con la identidad del firmante para prevenir ataques de sustitución de clave.

La función recibe como parámetros la clave privada $sk$, el mensaje en claro $M$ y, de forma opcional, una cadena de contexto $ctx$. Este contexto es una matriz de \textit{bytes} estrictamente acotada a un máximo de 255 posiciones. Su propósito es actuar como separador de dominio a nivel de aplicación; por ejemplo, impidiendo que una firma válida generada en un entorno de pruebas pueda ser inyectada maliciosamente en un entorno de producción. 

La fase de preparación culmina con el cálculo del \textbf{representante del mensaje}, llamado como $\mu$, un resumen criptográfico de 64 \textit{bytes}. El estándar FIPS 204 especifica que este valor $\mu$ se obtiene procesando mediante la función \texttt{SHAKE-256} la concatenación del \textit{hash} de la clave pública ($tr$, el cual se almacena precalculado dentro de $sk$) y los metadatos del mensaje:

\begin{equation}
    \mu \leftarrow \text{SHAKE-256}(tr \parallel \text{dom} \parallel \text{len}(ctx) \parallel ctx \parallel M, 64)
\end{equation}

El orden de los prefijos en esta ecuación es inmutable:
\begin{itemize}
    \item \textbf{$tr$ Huella de la Clave Pública:} Un vector de 64 \textit{bytes} que contiene el resumen criptográfico  de la clave pública completa. En lugar de recalcular este valor en cada operación de firma, el estándar establece que se precalcule durante la fase de generación y se persista dentro de la estructura de la clave privada $sk$ como una optimización de rendimiento. Al inyectar $tr$ en la derivación del representante del mensaje $\mu$, el sistema mitiga los ataques y de sustitución de clave, garantizando que la firma resultante sea criptográficamente inválida si un atacante intenta verificarla contra una clave pública fraudulenta.
    \item $\text{dom}$: Un único \textit{byte} separador de dominio (fijado a \texttt{0x00} para el esquema ML-DSA puro).
    \item $\text{len } ctx$: Un \textit{byte} que representa la longitud del contexto en formato numérico. Si no hay contexto, este \textit{byte} se fija a \texttt{0x00} y se incluye igualmente en la digestión.
\end{itemize}




\begin{figure}[H]
    \centering
    \resizebox{0.85\textwidth}{!}{
    \begin{tikzpicture}[
        node distance = 1cm and 0.5cm,
        caja/.style = {draw, rectangle, rounded corners, minimum height=0.8cm, align=center, fill=gray!5, thick},
        proceso/.style = {draw, ellipse, minimum width=2.8cm, minimum height=0.8cm, align=center, fill=white, thick},
        resultado/.style = {draw, rectangle, rounded corners, minimum width=3.5cm, minimum height=0.8cm, align=center, thick, font=\bfseries},
        flecha/.style = {-{Stealth[scale=1.2]}, thick}
    ]

    \node[caja] (tr) {$tr$ (64 bytes)};
    \node[caja, right=0.3cm of tr] (dom) {$dom$ (0x00)};
    \node[caja, right=0.3cm of dom] (len) {len($ctx$)};
    \node[caja, right=0.3cm of len] (ctx) {Contexto};
    \node[caja, right=0.3cm of ctx] (M) {Mensaje $M$};

    \node[draw, rectangle, dashed, thick, fit=(tr) (dom) (len) (ctx) (M), inner sep=0.2cm] (concat) {};
    \node[above=0.1cm of concat, font=\small\bfseries] {Concatenación en orden estricto};

    \node[proceso, below=1.2cm of concat] (shake1) {\texttt{SHAKE-256}};
    \node[resultado, fill=blue!5, below=1cm of shake1] (mu) {Representante del Mensaje $\mu$ (64 bytes)};

    \draw[flecha] (concat.south) -- (shake1.north);
    \draw[flecha] (shake1.south) -- (mu.north);

    \end{tikzpicture}
    }
    \caption{Derivación del representante del mensaje $\mu$.}
    \label{fig:sign_preparacion}
\end{figure}


\subsection{El Secreto Temporal y el Compromiso $\mathbf{w}_1$}

Con el mensaje original M está vinculado en el representante $\mu$, el algoritmo entra en el núcleo del esquema de firma: el bucle iterativo fundamentado en el paradigma \textit{Fiat-Shamir with Aborts} \cite{lyubashevsky2009fiat}.

El desafío central de esta fase es la protección de la clave privada. Si el sistema operara el mensaje directamente contra los vectores secretos $\mathbf{s}_1$ y $\mathbf{s}_2$, cada firma generada filtraría información estadística sobre sus coeficientes, exponiendo la clave tras múltiples usos. Para evitar esta vulnerabilidad, el paradigma exige generar un vector de enmascaramiento temporal $\mathbf{y}$ que ciega los valores secretos. La evaluación de esta máscara contra la matriz pública produce lo que se denomina el \textit{compromiso} $\mathbf{w}_1$. Para iniciar este proceso, se requiere una semilla. 

\textit{Nota sobre la nomenclatura:} Aunque el pseudocódigo del estándar FIPS 204 sobreescribe la variable $\rho'$ para derivar esta semilla durante la firma, en esta memoria se utilizará la notación $\rho_{sign}$ para evitar confusiones con la semilla de los secretos generada en la Fase 1.

Esta semilla de firma se calcula como $\rho_{sign} = \texttt{SHAKE-256}(K \parallel \mu, 64)$. Al incluir $K$, que estaba almacenada en la clave privada, el resultado es impredecible para un atacante, pero perfectamente determinista para el firmante. A continuación, se inicializa un contador de intentos $\kappa = 0$ y se arranca el bucle \texttt{while}, ejecutando los siguientes pasos algorítmicos en cada iteración.

La Figura \ref{fig:sign_rho_prima} muestra esta fase de preparación. La inyección de la semilla privada $K$ en la función de salida extensible asegura que la semilla resultante ($\rho_{sign}$) sea criptográficamente segura e impredecible para un observador externo, estableciendo una base determinista robusta antes de iniciar las iteraciones del bucle de firma.

\begin{figure}[H]
    \centering
    \begin{tikzpicture}[
        node distance = 1.5cm and 2cm,
        caja/.style = {draw, rectangle, rounded corners, minimum height=1cm, align=center, fill=gray!5, thick},
        proceso/.style = {draw, ellipse, minimum height=1cm, align=center, fill=white, thick},
        flecha/.style = {-{Stealth[scale=1.2]}, thick}
    ]

    \node[caja] (K) {Semilla secreta $K$ (32 bytes)};
    \node[caja, right=2.5cm of K] (mu) {Representante $\mu$ (64 bytes)};
    
    \node[proceso, below right=1.2cm and -1.2cm of K] (shake) {\texttt{SHAKE-256}};
    \node[caja, fill=green!5, below=1.2cm of shake] (rho) {Semilla de firma $\rho_{sign}$ (64 bytes)};

    \draw[flecha] (K) |- (shake);
    \draw[flecha] (mu) |- (shake);
    \draw[flecha] (shake) -- (rho);

    \end{tikzpicture}
    \caption{Derivación de la semilla $\rho_{sign}$.}
    \label{fig:sign_rho_prima}
\end{figure}

\begin{enumerate}
    \item \textbf{El Secreto Temporal $\mathbf{y}$:} Constituye la máscara criptográfica de la iteración. Se genera mediante \texttt{SHAKE-256} utilizando la semilla $\rho_{sign}$ y el contador $\kappa$ como separador de dominio. Es un vector de $l$ polinomios y, al igual que los secretos de la clave, sus coeficientes son enteros. Sin embargo, su rango es considerablemente mayor, ya que deben pertenecer al intervalo $[-\gamma_1 + 1, \gamma_1]$. Este umbral $\gamma_1$ debe ser lo suficientemente grande para esconder cualquier rastro de la clave privada, pero no tan masivo como para invalidar la verificación matemática posterior. Cada vez que el bucle aborta, $\kappa$ se incrementa en $l$, garantizando que el próximo vector $\mathbf{y}$ generado sea completamente distinto al anterior.
    
    \item \textbf{Proyección sobre el Retículo Público $\mathbf{w}$:} Una vez generado el secreto temporal $\mathbf{y}$, el firmante debe comprometerse a él vinculándolo con su identidad pública (la matriz $\mathbf{A}$). La ecuación teórica es $\mathbf{w} = \mathbf{A}\mathbf{y}$. A nivel de implementación, esta fase aprovecha una optimización arquitectónica masiva del estándar FIPS 204. La matriz $\mathbf{A}$, que se expande a partir de la semilla $\rho$ pública, se genera directamente en el dominio de la NTT $\mathbf{\hat{A}}$. Por tanto, el sistema solo necesita aplicar la transformada al vector temporal para operar:
    \begin{equation}
        \mathbf{\hat{w}} = \mathbf{\hat{A}} \cdot \text{NTT}(\mathbf{y})
    \end{equation}
    El resultado $\mathbf{\hat{w}}$ es un vector de dimensión $k$. Para poder manipular sus \textit{bits} en los pasos siguientes, se extrae del dominio transformado aplicando la inversa $\text{NTT}^{-1}$, obteniendo el vector polinómico en el dominio inicial $\mathbf{w}$.
    
    \item \textbf{Extracción de la Componente Relevante $\mathbf{w}_1$:} El vector $\mathbf{w}$ completo presenta una alta volatilidad en sus \textit{bits} menos significativos. Para optimizar el tamaño de la firma, se extraen únicamente sus \textit{bits} más significativos mediante la función \texttt{HighBits}:
    \begin{equation}
        \mathbf{w}_1 = \text{HighBits}(\mathbf{w}, \gamma_2)
    \end{equation}
\end{enumerate}

El valor $\mathbf{w}_1$ resultante actúa como el compromiso formal. Representa la componente de \textit{bits} más significativos extraída directamente del vector $\mathbf{w} = \mathbf{A}\mathbf{y}$. Al aislar esta información mediante la función \texttt{HighBits}, se obtiene el vector estable que se introducirá en la función de \textit{hash} criptográfica para generar el desafío polinómico contra el que el firmante deberá responder.

La Figura \ref{fig:sign_compromiso_bucle} detalla  la secuencia de operaciones que forman parte de la generación de este compromiso. El vector de enmascaramiento $\mathbf{y}$ se somete a las transformaciones polinómicas necesarias, NTT y su inversa, para operar eficientemente contra la matriz pública $\mathbf{\hat{A}}$, terminando con el uso de la función \texttt{HighBits} para aislar los bits más significativos $\mathbf{w}_1$.

\begin{figure}[H]
    \centering
    \begin{tikzpicture}[
        node distance = 1.2cm and 1.5cm,
        caja/.style = {draw, rectangle, rounded corners, minimum width=4.5cm, minimum height=0.8cm, align=center, fill=gray!5, thick},
        proceso/.style = {draw, ellipse, minimum height=0.8cm, align=center, fill=white, thick},
        flecha/.style = {-{Stealth[scale=1.2]}, thick}
    ]

    \node[caja, fill=green!5] (rho) {Semilla $\rho_{sign}$};
    
    \node[proceso, below=1.2cm of rho] (exp) {\texttt{ExpandMask}($\rho_{sign}, \kappa$)};
    \node[caja, below=0.8cm of exp] (y) {Secreto temporal $\mathbf{y} \in [-\gamma_1+1, \gamma_1]$};
    \node[proceso, below=0.8cm of y] (ntt) {Transformada $\text{NTT}(\mathbf{y})$};
    
    \node[caja, right=1.5cm of ntt, fill=orange!10] (A_hat) {Matriz Pública $\mathbf{\hat{A}}$};
    
    \node[caja, dashed, thick, below=0.8cm of ntt] (mult) {$\mathbf{\hat{w}} = \mathbf{\hat{A}} \cdot \text{NTT}(\mathbf{y})$};
    \node[proceso, below=0.8cm of mult] (invNTT) {Transformada Inversa $\text{NTT}^{-1}$};
    \node[caja, fill=blue!5, below=0.8cm of invNTT] (w1) {Compromiso $\mathbf{w}_1 = \texttt{HighBits}(\mathbf{w}, \gamma_2)$};

    \draw[flecha] (rho) -- (exp) node[midway, right, font=\small] {Inicio Bucle \texttt{while} (contador $\kappa$)};
    \draw[flecha] (exp) -- (y);
    \draw[flecha] (y) -- (ntt);
    \draw[flecha] (ntt) -- (mult);
    \draw[flecha] (A_hat) |- (mult);
    \draw[flecha] (mult) -- (invNTT);
    \draw[flecha] (invNTT) -- (w1);

    \node[draw, dashed, thick, rounded corners, fit=(exp) (y) (ntt) (mult) (invNTT) (w1), inner sep=15pt] (caja_bucle) {};
    \node[above right=0.2cm and -2cm of caja_bucle, font=\bfseries] {Cuerpo del Bucle Iterativo};

    \end{tikzpicture}
    \caption{Cálculo del compromiso $\mathbf{w}_1$ mediante enmascaramiento con $\mathbf{y}$.}
    \label{fig:sign_compromiso_bucle}
\end{figure}


\subsection{El Desafío Polinómico y la Función \texttt{SampleInBall}}

En este punto del algoritmo, el objetivo es crear un desafío que el firmante no pueda manipular. El proceso se resume en los siguientes pasos:

\begin{enumerate}
    \item \textbf{El Compromiso:} El firmante fija su posición inicial mediante el vector $\mathbf{w}_1$. Una vez generado este valor, ya no puede dar marcha atrás ni modificarlo sin invalidar todo el proceso.
    
    \item \textbf{La Simulación del Verificador:} Normalmente, un verificador externo debería elegir un número al azar y lanzarlo como desafío. Para evitar esta interacción y que la firma sea autónoma, aplicamos la \textit{Heurística de Fiat-Shamir} \cite{lyubashevsky2009fiat}, es decir, el firmante genera un número aleatorio usando una función \textit{hash} cuya salida no puede controlar.
    
    \item \textbf{Cálculo del Hash:} En lugar de esperar a un tercero, se crea el desafío pasando el vector $\mathbf{w}_1$ por una función resumen \textit{hash}. El resultado es un valor que actúa como el desafío.
\end{enumerate}

El sistema utiliza el propio punto de partida $\mathbf{w}_1$ para generar las preguntas del desafío. Si se intentara manipular el punto de partida para que el desafío sea fácil, el hash cambiará y dará un desafío distinto, haciendo que sea imposible hacer trampas.

El proceso consta de dos etapas estructurales: la generación de la semilla y su expansión al dominio polinómico.

\subsubsection{Generación de la semilla del desafío $\tilde{c}$}
En primer lugar, el vector de polinomios $\mathbf{w}_1$ se comprime y serializa en una cadena de \textit{bytes} mediante la función \texttt{w1Encode}. A continuación, se concatena con el representante del mensaje $\mu$ y se inyecta en \texttt{SHAKE-256} para obtener la semilla pseudoaleatoria $\tilde{c}$:

\begin{equation}
    \tilde{c} \leftarrow \text{SHAKE-256}(\mu \parallel \text{w1Encode}(\mathbf{w}_1), \lambda / 8)
\end{equation}

La longitud de la salida ($\lambda / 8$) varía según el nivel de seguridad seleccionado: 32 \textit{bytes} para ML-DSA-44, y 64 \textit{bytes} para ML-DSA-65 y ML-DSA-87. Esta cadena $\tilde{c}$ es el valor exacto que se adjuntará en la firma digital final, ya que su reducida huella en memoria optimiza el ancho de banda.

La Figura~\ref{fig:sign_desafio_alto_nivel} muestra la fase de desafío. El representante del mensaje $\mu$ y la versión codificada del compromiso $\mathbf{w}_1$ coinciden en la función de salida extensible para generar la semilla $\tilde{c}$. Esta semilla es usada inmediatamente por el algoritmo \texttt{SampleInBall}, el cual, a partir de la entrada $\tilde{c}$, genera un polinomio en $B_\tau$, es decir, con todos los coeficientes iguales a $0$, salvo $\tau$ coeficientes que toman el valor $1$ o $-1$.


\begin{figure}[H]
    \centering
    \begin{tikzpicture}[
        node distance = 1cm and 1.5cm,
        caja/.style = {draw, rectangle, rounded corners, minimum height=0.8cm, align=center, fill=gray!5, thick},
        caja_clara/.style = {draw, rectangle, rounded corners, dashed, minimum height=0.6cm, align=center, fill=white, thick},
        proceso/.style = {draw, ellipse, minimum width=2.5cm, minimum height=0.8cm, align=center, fill=white, thick},
        resultado/.style = {draw, rectangle, minimum height=0.8cm, align=center, fill=blue!5, thick, font=\bfseries},
        flecha/.style = {-{Stealth[scale=1.2]}, thick}
    ]

    \node[proceso] (hash) {\texttt{SHAKE-256}};

    \node[caja_clara, right=1.5cm of hash] (bytes) {Bytes de $\mathbf{w}_1$};
    \node[proceso, above=0.8cm of bytes] (encode) {\texttt{w1Encode}()};
    \node[caja, above=0.8cm of encode] (w1) {Compromiso $\mathbf{w}_1$};

    \node[caja] at (hash |- w1) (mu) {Representante $\mu$ (64 bytes)};

    \draw[flecha] (w1) -- (encode);
    \draw[flecha] (encode) -- (bytes);
    \draw[flecha] (bytes) -- (hash);
    \draw[flecha] (mu) -- (hash);

    \node[resultado, below=1cm of hash] (c_tilde) {Semilla del desafío $\tilde{c}$ (32/64 bytes)};
    \node[proceso, below=1cm of c_tilde] (sib) {\texttt{SampleInBall}($\tilde{c}, \tau$)};
    \node[resultado, fill=green!5, below=1cm of sib] (c_poly) {Polinomio $c \in R$};

    \draw[flecha] (hash) -- (c_tilde);
    \draw[flecha] (c_tilde) -- (sib);
    \draw[flecha] (sib) -- (c_poly);

    \end{tikzpicture}
    \caption{Derivación del desafío $c$ a partir del compromiso $\mathbf{w}_1$.}
    \label{fig:sign_desafio_alto_nivel}
\end{figure}

\subsubsection{Expansión mediante \texttt{SampleInBall} $c$}
La semilla $\tilde{c}$ es una cadena de \textit{bytes}, incompatible con los vectores del anillo $R_q$ necesarios para las operaciones. Para transformarla, se utiliza la función \texttt{SampleInBall}($\tilde{c}, \tau$), que devuelve un polinomio en $B_\tau$.

Esta función no genera un polinomio aleatorio estándar. Para mantener la seguridad del problema MLWE y evitar que el ruido inyectado en la firma se descontrole, $c$ debe ser un polinomio extremadamente disperso. Exactamente $\tau$ coeficientes tendrán un valor de $+1$ o $-1$, y los $256 - \tau$ coeficientes restantes serán estrictamente $0$. El parámetro $\tau$ equivale a 39, 49 o 60, dependiendo del nivel de seguridad.

\vspace{0.5cm}
\begin{ejemplo}[Implementación Algorítmica de \texttt{SampleInBall} (ML-DSA-65, $\tau = 49$)]
\label{ex:sample_in_ball}
El algoritmo no realiza una búsqueda aleatoria de 49 posiciones libres, lo cual generaría tiempos de ejecución inestables debido a las colisiones. En su lugar, emplea una variante del barajado de Fisher-Yates, haciendo uso del algoritmo \textit{inside-out} \cite{durstenfeld1964algorithm}, para rellenar un polinomio $c$ inicializado a ceros. El proceso se ejecuta en dos fases:

\begin{enumerate}
    \item \textbf{Reserva de Signos:} Primero, el algoritmo expande la semilla a través de un XOF y extrae de inmediato los primeros 8 \textit{bytes}. Esto proporciona 64 \textit{bits} secuenciales en memoria, cantidad suficiente para cubrir el signo de los $\tau = 49$ coeficientes no nulos que se van a generar.
    
    \item \textbf{Inserción Iterativa (Barajado):} A continuación, el algoritmo garantiza que se colocarán exactamente $\tau$ elementos iterando en inversa. En ML-DSA-65, el bucle recorre los índices $i$ desde $207$ ($256 - 49$) hasta $255$. En cada iteración $i$:
    \begin{itemize}
        \item Se extrae un nuevo \textit{byte} del XOF. Se aplica muestreo de rechazo estricto hasta obtener un valor válido $j \le i$.
        \item Para evitar sobrescrituras accidentales y garantizar la dispersión, el algoritmo copia el valor actual de $c[j]$ a la nueva posición segura $c[i]$.
        \item Finalmente, inyecta en $c[j]$ el valor $+1$ o $-1$, dependiendo de si el siguiente \textit{bit} consumido de la Fase 1 es un \texttt{0} o un \texttt{1}.
    \end{itemize}
\end{enumerate}

El resultado final es un polinomio de grado 255 donde exactamente $49$ componentes son $\pm 1$ y el resto son estrictamente nulos. Esta dispersión es lo que permite que la posterior multiplicación polinómica contra el vector secreto $\mathbf{s}_1$ sea altamente eficiente, operación central para generar la respuesta de la firma.
\end{ejemplo}

Para visualizar de forma correcta el uso del algoritmo \textit{inside-out}, la Figura \ref{fig:sign_sampleinball} muestra el flujo de control interno de \texttt{SampleInBall}. Se puede observar cómo la extracción de entropía se segmenta para garantizar un tiempo de ejecución menor, resolviendo primero las asignaciones de signo y la lógica de dispersión en un bucle iterativo.

\begin{figure}[H]
    \centering
    \resizebox{0.95\textwidth}{!}{
    \begin{tikzpicture}[
        node distance = 0.5cm and 1.2cm,
        caja/.style = {draw, rectangle, rounded corners, minimum height=0.8cm, align=center, fill=gray!5, thick},
        proceso/.style = {draw, rectangle, align=center, fill=white, thick, minimum height=0.8cm, text width=3.5cm},
        rombo/.style = {draw, diamond, aspect=1.5, align=center, fill=yellow!10, thick},
        flecha/.style = {-{Stealth[scale=1.2]}, thick}
    ]

    \node[caja, fill=blue!5] (inicio) {Semilla $\tilde{c}$};
    \node[proceso, text width=4.5cm, below=0.5cm of inicio] (xof) {Generar flujo de \textit{bytes} (\texttt{SHAKE})};
    \draw[flecha] (inicio) -- (xof);

    \node[proceso, below=0.8cm of xof] (fase1_byte) {Extraer 1 \textit{byte} ($b$)};
    \node[rombo, below=0.5cm of fase1_byte] (fase1_valido) {$b \le 255$ y\\no reservado?};
    \node[proceso, right=1.5cm of fase1_valido] (fase1_add) {Reservar posición $b$};
    \node[rombo, below=0.5cm of fase1_add] (fase1_fin) {¿Total $\tau$ reservas?};

    \draw[flecha] (xof) -- (fase1_byte);
    \draw[flecha] (fase1_byte) -- (fase1_valido);
    \draw[flecha] (fase1_valido) -- node[above] {Sí} (fase1_add);
    \draw[flecha] (fase1_valido.west) -- ++(-1.5,0) |- node[near start, right] {No (Rechazo)} (fase1_byte.west);
    \draw[flecha] (fase1_add) -- (fase1_fin);
    \draw[flecha] (fase1_fin.east) -- ++(1.4,0) |- node[near start, right] {No} (fase1_byte.east);

    \node[caja, dashed, below=0.8cm of fase1_fin] (lista) {Lista con $\tau$ posiciones ordenadas};
    \draw[flecha] (fase1_fin) -- node[right] {Sí} (lista);

    \node[proceso, below=0.8cm of lista] (fase2_bit) {Para cada posición:\\Extraer 1 \textit{bit} del flujo};
    \node[rombo, below=0.5cm of fase2_bit] (fase2_signo) {¿Valor del \textit{bit}?};
    \node[proceso, below left=0.6cm and 0.3cm of fase2_signo] (pos) {Coeficiente = $+1$};
    \node[proceso, below right=0.6cm and 0.3cm of fase2_signo] (neg) {Coeficiente = $-1$};

    \draw[flecha] (lista) -- (fase2_bit);
    \draw[flecha] (fase2_bit) -- (fase2_signo);
    \draw[flecha] (fase2_signo) -| node[above, pos=0.2] {\texttt{0}} (pos);
    \draw[flecha] (fase2_signo) -| node[above, pos=0.2] {\texttt{1}} (neg);

    \node[caja, fill=green!5, below=1.8cm of fase2_signo] (final) {Polinomio disperso $c$ completado};
    \draw[flecha] (pos) |- (final);
    \draw[flecha] (neg) |- (final);

    \node[draw, dashed, thick, rounded corners, fit=(fase1_byte) (fase1_valido) (fase1_add) (fase1_fin), inner sep=0.4cm] (box1) {};
    \node[fill=white, font=\bfseries, anchor=west, xshift=0.5cm] at (box1.north west) {Fase A: Marcaje (Muestreo de Rechazo)};

    \node[draw, dashed, thick, rounded corners, fit=(fase2_bit) (fase2_signo) (pos) (neg), inner sep=0.4cm] (box2) {};
    \node[fill=white, font=\bfseries, anchor=west, xshift=0.5cm] at (box2.north west) {Fase B: Asignación de Signos};

    \end{tikzpicture}
    }
    \caption{Diagrama de flujo interno de la función \texttt{SampleInBall}.}
    \label{fig:sign_sampleinball}
\end{figure}

\subsection{La Respuesta $\mathbf{z}$ y el Muestreo de Rechazo}
\label{subsec:respuesta_z_rechazo}

Con el secreto temporal $\mathbf{y}$, y el desafío $c$ calculados, el bucle iterativo entra en su fase final. El objetivo es generar una respuesta polinómica $\mathbf{z}$ que, junto con el desafío, demuestre la posesión de la clave privada sin revelar sus coeficientes.

La ecuación fundamental de la respuesta se define como:
\begin{equation}
    \mathbf{z} = \mathbf{y} + c \cdot \mathbf{s}_1
\end{equation}

El vector secreto $\mathbf{s}_1$ se extrae de la clave privada. Para realizar la multiplicación, tanto $c$ como $\mathbf{s}_1$ deben ser proyectados al dominio \texttt{NTT}. Tras la multiplicación coeficiente a coeficiente, el producto se devuelve al dominio del anillo cociente mediante la aplicación $\text{NTT}^{-1}$ para poder sumarse de forma directa con el vector de enmascaramiento $\mathbf{y}$.

\subsubsection{Filtros de Seguridad: Condiciones de Salida del Bucle}

El vector $\mathbf{z}$ resultante es la firma. Sin embargo, en el paradigma \textit{Fiat-Shamir with Aborts}, no se puede publicar $\mathbf{z}$ inmediatamente. Si los coeficientes de $\mathbf{z}$ o del error interno son demasiado grandes, la firma filtraría información estadística sobre $\mathbf{s}_1$ y $\mathbf{s}_2$, comprometiendo la clave privada a largo plazo.

Para evitar esto, FIPS 204 implementa un muestreo de rechazo basado en la norma infinito ($||\cdot||_\infty$). Se debe aplicar una reducción modular simétrica a todos los coeficientes del vector, extraer sus valores absolutos y seleccionar el máximo de todos ellos. Si este valor máximo excede la cota de seguridad predefinida, la firma filtra información, por lo que se rechaza y el proceso de generación se reinicia.

El algoritmo realiza dos comprobaciones:

\begin{enumerate}
    \item \textbf{Filtro de Privacidad (Check sobre $\mathbf{z}$):} Se evalúa que ningún coeficiente del vector respuesta supere el umbral $\gamma_1 - \beta$. Si $|| \mathbf{z} ||_\infty \ge \gamma_1 - \beta$, la firma se considera no útil y se rechaza. Si no se verifica la condición, un oponente podría obtener información de la clave privada a partir de la firma.
    
    \item \textbf{Filtro de Corrección (Check sobre $\mathbf{r}_0$):} Se evalúa el ruido en los \textit{bits} menos significativos. Se calcula el vector interno $\mathbf{r}_0 = \text{LowBits}(\mathbf{w} - c \cdot \mathbf{s}_2, 2\gamma_2)$. Al igual que en el caso anterior, si $|| \mathbf{r}_0 ||_\infty \ge \gamma_2 - \beta$, la firma provocaría un fallo en la fase de verificación, por lo que se rechaza.
\end{enumerate}

La Figura \ref{fig:sign_rechazo} muestra el flujo de este proceso de muestreo de rechazo \textit{Rejection Sampling}. Ambas comprobaciones actúan de manera estricta; el fallo en cualquiera de los umbrales anularía la firma, obligando a la creación de un nuevo secreto temporal $\mathbf{y}$ para evitar fugas estadísticas. Únicamente si ambos filtros se superan, el sistema avanza hacia la creación del vector de pistas y el empaquetado de la firma final.

\begin{figure}[H]
    \centering
    \begin{tikzpicture}[
        node distance = 0.8cm and 1cm,
        caja/.style = {draw, rectangle, rounded corners, minimum height=0.8cm, align=center, fill=gray!5, thick},
        proceso/.style = {draw, rectangle, align=center, fill=white, thick, minimum height=0.8cm},
        rombo/.style = {draw, diamond, aspect=1.6, align=center, fill=yellow!10, thick},
        resultado/.style = {draw, rectangle, rounded corners, minimum height=0.8cm, align=center, fill=green!10, thick, font=\bfseries},
        flecha/.style = {-{Stealth[scale=1.2]}, thick}
    ]

    \node[proceso, fill=blue!5] (gen_y) {Generar vector $\mathbf{y}$};

    \node[proceso, below=0.8cm of gen_y] (calc_z) {Calcular $\mathbf{z} = \mathbf{y} + c \cdot \mathbf{s}_1$};
    \node[proceso, below=0.6cm of calc_z] (calc_r) {Calcular $\mathbf{r}_0$ de $\mathbf{w} - c \cdot \mathbf{s}_2$};
    
    \draw[flecha] (gen_y) -- (calc_z);
    \draw[flecha] (calc_z) -- (calc_r);

    \node[rombo, below=0.8cm of calc_r] (check1) {$||\mathbf{z}||_\infty < \gamma_1 - \beta$ ?};
    \draw[flecha] (calc_r) -- (check1);

    \node[rombo, below=1cm of check1] (check2) {$||\mathbf{r}_0||_\infty < \gamma_2 - \beta$ ?};
    \draw[flecha] (check1) -- node[right] {Sí} (check2);

    \node[proceso, left=2.5cm of check1, fill=red!10] (reject) {Rechazo: $\kappa = \kappa + l$};
    
    \draw[flecha] (check1) -- node[above] {No} (reject);
    \draw[flecha] (check2) -| node[above, near start] {No} (reject);
    
    \draw[flecha, thick, red] (reject.north) |- (gen_y.west);

    \node[resultado, below=1cm of check2] (firma) {Firma Válida $\sigma = (\tilde{c}, \mathbf{z}, h)$ \\[1.5mm] \normalfont\small ($h$: vector de pistas)};
    \draw[flecha] (check2) -- node[right] {Sí} (firma);

    \end{tikzpicture}
    \caption{Muestreo de rechazo en el algoritmo de firma.}
    \label{fig:sign_rechazo}
\end{figure}


Los límites de estos rechazos están codificados en los parámetros del estándar, tal y como se refleja en la Tabla \ref{tab:limites_rechazo}. Por ejemplo, para el nivel ML-DSA-65, un solo coeficiente de $\mathbf{z}$ que alcance el valor $524.092$ invalidará instantáneamente la iteración.

\begin{table}[H]
    \centering
    \renewcommand{\arraystretch}{1.4}
    \resizebox{\textwidth}{!}{
    \begin{tabular}{lccc}
        \toprule
        \textbf{Parámetro} & \textbf{ML-DSA-44} (Nivel 2) & \textbf{ML-DSA-65} (Nivel 3) & \textbf{ML-DSA-87} (Nivel 5) \\
        \midrule
        $\gamma_1$ (Rango de $\mathbf{y}$) & $2^{17} = 131.072$ & $2^{19} = 524.288$ & $2^{19} = 524.288$ \\
        $\gamma_2$ (Rango de $\mathbf{w}_1$) & $(q-1)/88 = 95.232$ & $(q-1)/32 = 261.888$ & $(q-1)/32 = 261.888$ \\
        $\beta$ (Margen de rechazo) & 78 & 196 & 120 \\
        \midrule
        \textbf{Límite Check 1} ($\gamma_1 - \beta$) & $\mathbf{130.994}$ & $\mathbf{524.092}$ & $\mathbf{524.168}$ \\
        \textbf{Límite Check 2} ($\gamma_2 - \beta$) & $\mathbf{95.154}$ & $\mathbf{261.692}$ & $\mathbf{261.768}$ \\
        \bottomrule
    \end{tabular}
    }
    \caption{Parámetros y límites máximos permitidos en los filtros de rechazo del bucle de firma.}
    \label{tab:limites_rechazo}
\end{table}

La probabilidad de que una firma candidata sea apta, es decir, que verifique las dos condiciones de muestreo de rechazo descritas anteriormente sin filtrar información sensible, se aproxima mediante la siguiente expresión:

$$ \Pr[\text{Apta}] \approx e^{-256 \cdot \beta \left( \ell/\gamma_1 + k/\gamma_2 \right)} $$


\subsection{El Vector de Pistas $\mathbf{h}$}

Si la firma supera los estrictos filtros de rechazo de $\mathbf{z}$ y $\mathbf{r}_0$, el algoritmo entra en su fase final. En este punto, Alicia posee una firma criptográficamente segura, pero debe asegurarse de que el receptor (Bob) sea matemáticamente capaz de verificarla. 

El problema radica en el algoritmo de compresión \texttt{Power2Round} ejecutado durante la generación de claves. Bob únicamente posee el vector público $\mathbf{t}_1$ (junto con la semilla $\rho$ para expandir la matriz $\mathbf{A}$). Carece del residuo $\mathbf{t}_0$, que quedó confinado en la clave privada de Alicia. Cuando Bob intente reconstruir el compromiso $\mathbf{w}_1$ utilizando la ecuación del retículo, su resultado carecerá de la precisión que aporta $\mathbf{t}_0$. En la mayoría de los coeficientes, esta pequeña variación no afectará al cálculo de los \textit{bits} altos (\texttt{HighBits}). Sin embargo, si un coeficiente original estaba muy cerca de un límite de truncamiento, la ausencia de $\mathbf{t}_0$ provocará que Bob obtenga un \textit{bit} alto diferente al que calculó Alicia.

Para evitar que la verificación fracase, Alicia ejecuta la función \texttt{MakeHint}, simulando los cálculos del receptor para generar un "vector de pistas" polinómico $\mathbf{h}$ compuesto exclusivamente por ceros y unos:

\begin{enumerate}
    \item \textbf{Cálculo de la perspectiva de Alicia $\mathbf{r}_1$:} Alicia calcula el vector intermedio exacto de Bob $\mathbf{w}'$ sumando el residuo $\mathbf{t}_0$ a la aproximación: $\mathbf{w}' = \mathbf{w} - c \cdot \mathbf{s}_2 + c \cdot \mathbf{t}_0$. Seguidamente, le aplica la función \texttt{HighBits} dividiendo por $2\gamma_2$ para obtener $\mathbf{r}_1$. Como opera con polinomios, los operandos $c$, $\mathbf{s}_2$ y $\mathbf{t}_0$ deben ser transformados dinámicamente mediante \texttt{NTT} antes de multiplicarse, y devueltos al dominio temporal con $\text{NTT}^{-1}$ para realizar las sumas y particiones. El resultado es el vector exacto $\mathbf{r}_1$.
    \item \textbf{Cálculo de la perspectiva de Bob $\mathbf{v}_1$:} Alicia simula la carencia de Bob eliminando $\mathbf{t}_0$ de la ecuación. Calcula la componente alta de $\mathbf{w} - c \cdot \mathbf{s}_2$. El resultado es $\mathbf{v}_1$.
    \item \textbf{Generación de la pista $\mathbf{h}$:} Se comparan ambos vectores coeficiente a coeficiente. Si $\mathbf{r}_1[i][j] \neq \mathbf{v}_1[i][j]$, significa que la carencia de $\mathbf{t}_0$ corromperá el cálculo de Bob en esa posición específica. Alicia marca esa coordenada en el vector $\mathbf{h}$ con un \texttt{1}. Si los coeficientes coinciden, se marca con un \texttt{0}.
\end{enumerate}

\subsection{El Filtro Final y Empaquetado $\sigma$}

El vector $\mathbf{h}$ permite a Bob corregir las desviaciones sumando o restando una unidad a sus \textit{bits} altos allí donde Alicia haya marcado un \texttt{1}. No obstante, enviar demasiadas correcciones facilitaría ataques algebraicos, por lo que el estándar FIPS 204 impone un límite máximo de unos permitidos en la firma, denominado $\omega$ (\textit{omega}).

\begin{table}[H]
    \centering
    \renewcommand{\arraystretch}{1.3}
    \resizebox{\textwidth}{!}{
    \begin{tabular}{lccc}
        \toprule
        \textbf{Parámetro} & \textbf{ML-DSA-44} (Nivel 2) & \textbf{ML-DSA-65} (Nivel 3) & \textbf{ML-DSA-87} (Nivel 5) \\
        \midrule
        $\omega$ (Límite de pistas) & 80 & 55 & 75 \\
        \bottomrule
    \end{tabular}
    }
    \caption{Número máximo de correcciones permitidas en el vector $\mathbf{h}$.}
    \label{tab:limite_omega}
\end{table}

Antes de finalizar, Alicia suma todos los coeficientes del vector $\mathbf{h}$. Tomando como ejemplo el nivel intermedio (ML-DSA-65), si la cantidad de unos generados supera los $55$, Alicia descarta el vector $\mathbf{h}$, desecha la firma temporal, incrementa el contador de iteración ($\kappa = \kappa + l$) y vuelve a iniciar el bucle \texttt{while} desde la generación de una nueva máscara temporal $\mathbf{y}$.

Si el conteo es menor o igual a $\omega$, la firma es definitivamente válida. El motor criptográfico serializa y comprime los datos de salida. El empaquetado final $\sigma$ (\textit{sigma}) queda constituido estrictamente por tres elementos:
\begin{enumerate}
    \item $\tilde{c}$: La semilla original del desafío (32 o 64 \textit{bytes}, dependiendo del nivel).
    \item $\mathbf{z}$: El vector de respuesta polinómica empaquetado.
    \item $\mathbf{h}$: El vector de pistas, altamente comprimido debido a su naturaleza binaria dispersa.
\end{enumerate}

\begin{figure}[H]
    \centering
    \begin{tikzpicture}[
        caja/.style = {draw, rectangle, rounded corners, minimum height=0.8cm, inner sep=0.2cm, align=center, fill=gray!5, thick},
        proceso/.style = {draw, rectangle, align=center, fill=white, thick, minimum height=0.8cm, inner sep=0.2cm},
        decision/.style = {draw, diamond, aspect=2, align=center, fill=yellow!10, thick, inner sep=0.1cm},
        resultado/.style = {draw, rectangle, rounded corners, minimum height=0.8cm, align=center, fill=blue!5, thick, font=\bfseries},
        flecha/.style = {-{Stealth[scale=1.2]}, thick},
        linea/.style = {thick}
    ]

    \node[proceso] (Alicia)    at (-3.5, 0)    {$\mathbf{w} - c\cdot\mathbf{s}_2 + c\cdot\mathbf{t}_0$};
    \node[proceso] (hb_Alicia) at (-3.5, -1.5) {\texttt{HighBits}($\cdot, 2\gamma_2$)};
    \node[caja]    (r1)       at (-3.5, -3)   {Valor Real ($\mathbf{r}_1$)};

    \node[proceso] (bob)      at (3.5, 0)     {$\mathbf{w} - c\cdot\mathbf{s}_2$};
    \node[proceso] (hb_bob)   at (3.5, -1.5)  {\texttt{HighBits}($\cdot, 2\gamma_2$)};
    \node[caja]    (v1)       at (3.5, -3)    {Valor Simulado ($\mathbf{v}_1$)};

    \node[decision] (comp)    at (0, -5)      {¿$\mathbf{r}_1[i] \neq \mathbf{v}_1[i]$?};

    \node[caja, fill=red!10]   (h1) at (-3.5, -7) {$\mathbf{h}[i] = 1$};
    \node[caja, fill=green!10] (h0) at (3.5, -7)  {$\mathbf{h}[i] = 0$};

    \node[resultado] (h_final) at (0, -9) {Vector de Pistas Binario ($\mathbf{h}$)};


    \draw[flecha] (Alicia) -- (hb_Alicia);
    \draw[flecha] (bob) -- (hb_bob);
    \draw[flecha] (hb_Alicia) -- (r1);
    \draw[flecha] (hb_bob) -- (v1);

    \draw[linea] (r1.south) |- (0, -4);
    \draw[linea] (v1.south) |- (0, -4);
    \draw[flecha] (0, -4) -- (comp.north);

    \draw[flecha] (comp.west) -| node[above, near start] {Sí} (h1.north);
    \draw[flecha] (comp.east) -| node[above, near start] {No} (h0.north);

    \draw[linea] (h1.south) |- (0, -8);
    \draw[linea] (h0.south) |- (0, -8);
    \draw[flecha] (0, -8) -- (h_final.north);

    \end{tikzpicture}
    \caption{Simulación en \texttt{MakeHint}.}
    \label{fig:make_hint}
\end{figure}

\begin{figure}[H]
    \centering
    \begin{tikzpicture}[
        caja/.style = {draw, rectangle, rounded corners, minimum height=0.8cm, minimum width=2.8cm, align=center, fill=gray!5, thick},
        decision/.style = {draw, ellipse, minimum height=0.8cm, minimum width=4.5cm, align=center, fill=yellow!10, thick},
        paquete/.style = {draw, rectangle, dashed, rounded corners, minimum height=4.2cm, minimum width=3.8cm, align=center, fill=green!5, thick},
        flecha/.style = {-{Stealth[scale=1.2]}, thick},
        linea/.style = {thick}
    ]

    \node[caja]     (c)     at (0, 3)     {Semilla $\tilde{c}$};
    \node[caja]     (h)     at (0, 1.5)   {Vector $\mathbf{h}$};
    \node[caja]     (z)     at (0, 0)     {Respuesta $\mathbf{z}$};

    \node[decision] (check) at (5, 1.5)   {¿Suma de unos $\le \omega$?};

    \node[paquete]  (sigma) at (10, 1.5) {};
    \node[font=\bfseries]   at (10, 3)   {Firma Digital $\sigma$};
    \node[caja, fill=white] at (10, 1.5) {$(\tilde{c}, \mathbf{z}, \mathbf{h}_1)$};

    \draw[flecha] (h.east) -- (check.west);
    \draw[flecha] (check.east) -- node[above] {Sí} (sigma.west |- check.east);
    
    \draw[flecha] (c.east) -- (sigma.west |- c.east);
    \draw[flecha] (z.east) -- (sigma.west |- z.east);

    \draw[linea, red] (check.south) -- (5, -1.2);
    \draw[flecha, red] (5, -1.2) -- (0.5, -1.2) node[midway, below] {No: Rechazar firma y repetir \texttt{while}};

    \end{tikzpicture}
    \caption{Estructura del paquete de firma digital $\sigma$.}
    \label{fig:empaquetado_sigma}
\end{figure}

\section{Verificación de la Firma}
\label{sec:verificacion}

El proceso de verificación es la fase final del esquema ML-DSA. En esta fase, una entidad receptora (en adelante, Bob) recibe un mensaje $M$, una firma digital empaquetada $\sigma$ y la clave pública del firmante $pk$. El objetivo de Bob es reconstruir el desafío criptográfico utilizando únicamente estos elementos públicos, sin interactuar con el firmante ni conocer los secretos $\mathbf{s}_1$ y $\mathbf{s}_2$.

\subsection{Desempaquetado y Filtros Previos}

El primer paso consiste en extraer los componentes de los paquetes recibidos y derivar el contexto del mensaje. 

Por un lado, Bob desempaqueta la clave pública $pk$ para extraer la semilla de la matriz $\rho$ y el vector polinómico público $\mathbf{t}_1$. Por otro lado, desempaqueta la firma $\sigma$ aislando la semilla del desafío original $\tilde{c}$, el vector de respuesta $\mathbf{z}$ y el vector de pistas binario $\mathbf{h}$. Para transformar $\mathbf{z}$ y $\mathbf{h}$ desde sus cadenas de \textit{bytes} a sus respectivas estructuras de polinomios en $R_q$, se aplican las funciones inversas \texttt{Decode} y \texttt{DecodeHint}.

Antes de proceder con cualquier operación algebraica, el estándar FIPS 204 impone un control de seguridad estricto para evitar ataques de falsificación mediante vectores maliciosos. Bob debe verificar dos condiciones:
\begin{enumerate}
    \item Que la norma infinita del vector de respuesta esté dentro de los límites del nivel de seguridad: $||\mathbf{z}||_\infty < \gamma_1 - \beta$.
    \item Que el vector $\mathbf{h}$ sea válido.
\end{enumerate}
Si alguna de estas validaciones falla, el algoritmo aborta instantáneamente devolviendo un error (\texttt{REJECT}), ahorrando ciclos computacionales al procesador. Tras superar este filtro, Bob reconstruye el ancla del mensaje calculando primero el resumen de la clave pública $tr$ y, seguidamente, el ancla final contextual $\mu$:

\begin{equation}
    tr = \text{SHAKE-256}(\texttt{pkEncode}(pk))
\end{equation}

\begin{equation}
    \mu = \text{SHAKE-256}(tr \parallel M)
\end{equation}

\subsection{Reconstrucción Algebraica del Compromiso $\mathbf{w}'$}

Superados los filtros, Bob debe simular la proyección sobre el retículo que Alicia realizó en la fase de compromiso. La ecuación de reconstrucción se define como:
\begin{equation}
    \mathbf{w}' = \mathbf{A}\mathbf{z} - c \cdot \mathbf{t}_1 \cdot 2^d
\end{equation}

A nivel de arquitectura de \textit{software}, esta operación exige llevar las variables al dominio de la \texttt{NTT} para garantizar la eficiencia. La semilla $\tilde{c}$ se expande en el polinomio disperso $c$ mediante \texttt{SampleInBall}($\tilde{c}, \tau$) y se transforma a $\mathbf{\hat{c}}$. La matriz pública $\mathbf{\hat{A}}$ se genera directamente en \texttt{NTT} a partir de $\rho$. Tras realizar las multiplicaciones polinómicas y las restas correspondientes, el resultado se devuelve al dominio original mediante $\text{NTT}^{-1}$, obteniendo el vector $\mathbf{w}'_{approx}$.

Este vector $\mathbf{w}'_{approx}$ es una aproximación muy cercana a la proyección original $\mathbf{w}$ de Alicia, pero difiere ligeramente debido a que Bob no posee la parte baja de la clave $\mathbf{t}_0$, que es el resultado tras aplicar el algoritmo \texttt{Power2Round}.

\subsection{El Algoritmo de Corrección: \texttt{UseHint}}

Para compensar la carencia de $\mathbf{t}_0$, el esquema inyecta en la ecuación el vector de pistas $\mathbf{h}$ el \textit{hint} generado por Alicia. Bob ejecuta la función \texttt{UseHint}($\mathbf{h}, \mathbf{w}'_{approx}, 2\gamma_2$).

El propósito de esta función es recalcular los \textit{bits} más significativos del vector aproximado. \texttt{UseHint} itera sobre cada coeficiente: si el vector de pistas indica un \texttt{0} en esa posición, Bob extrae los \textit{bits} altos de forma estándar. Si el vector indica un \texttt{1}, Bob sabe que el error de truncamiento hizo que su coeficiente cayera  hacia el lado equivocado . En consecuencia, el algoritmo suma o resta una unidad matemática a la componente alta extraída, restaurando la sincronía exacta con los cálculos de Alicia.

El resultado de esta función es $\mathbf{w}'_1$, que debe ser idéntico al $\mathbf{w}_1$ original sobre el que Alicia se comprometió.

\subsection{Verificación Final}

Finalmente, Bob serializa el vector corregido pasándolo a una cadena de \textit{bytes} mediante \texttt{w1Encode}. Con $\mu$ y el compromiso recuperado $\mathbf{w}'_1$, reconstruye el \textit{hash} simulando el papel de la función de desafío de Fiat-Shamir:
\begin{equation}
    \tilde{c}' = \text{SHAKE-256}(\mu \parallel \texttt{w1Encode}(\mathbf{w}'_1))
\end{equation}

El algoritmo concluye con una comparación binaria en tiempo constante. Si la semilla reconstruida $\tilde{c}'$ es exactamente igual a la semilla original $\tilde{c}$ adjunta en la firma $\sigma$, se demuestra criptográficamente que quien generó $\sigma$ poseía la clave privada $\mathbf{s}_1, \mathbf{s}_2$ asociada a $\mathbf{t}_1$. La firma es aceptada. En cualquier otro caso, es rechazada.

\begin{figure}[H]
    \centering
    \resizebox{\textwidth}{!}{
    \begin{tikzpicture}[
        caja/.style = {draw, rectangle, rounded corners, minimum height=0.8cm, minimum width=3cm, align=center, fill=white, thick},
        proceso/.style = {draw, rectangle, align=center, fill=gray!5, thick, minimum height=0.8cm, minimum width=4cm},
        decision/.style = {draw, diamond, aspect=2, align=center, fill=yellow!10, thick},
        resultado_ok/.style = {draw, rectangle, rounded corners, minimum height=0.8cm, align=center, fill=green!10, thick, font=\bfseries},
        resultado_bad/.style = {draw, rectangle, rounded corners, minimum height=0.8cm, align=center, fill=red!10, thick, font=\bfseries},
        flecha/.style = {-{Stealth[scale=1.2]}, thick},
        linea/.style = {thick}
    ]

    \node[caja] (sigma) at (-3, 0) {Firma $\sigma = (\tilde{c}, \mathbf{z}, \mathbf{h})$};
    \node[caja] (mu)    at (3, 0)  {Mensaje y Ancla ($\mu$)};

    \node[decision] (check_z) at (0, -2) {$||\mathbf{z}||_\infty < \gamma_1 - \beta$};
    \node[resultado_bad] (reject1) at (-5.5, -2) {Rechazar (Falsificación)};

    \node[proceso] (w_approx) at (0, -4.5) {$\mathbf{w}'_{approx} = \mathbf{A}\mathbf{z} - c \cdot \mathbf{t}_1 \cdot 2^d$};

    \node[proceso] (use_hint) at (0, -6) {\texttt{UseHint}($\mathbf{h}, \mathbf{w}'_{approx}, 2\gamma_2$)};
    \node[caja, fill=blue!5] (w1_prima) at (0, -7.5) {Compromiso Recuperado ($\mathbf{w}'_1$)};

    \node[proceso] (hash) at (0, -9) {$\tilde{c}' = \text{SHAKE-256}(\mu \parallel \mathbf{w}'_1)$};

    \node[decision] (comparacion) at (0, -11) {¿ $\tilde{c}' == \tilde{c}$ ?};
    \node[resultado_ok] (accept)  at (5.5, -11) {Firma Válida};
    \node[resultado_bad] (reject2) at (-5.5, -11) {Firma Inválida};

    \draw[flecha] (sigma.south) |- (0, -1) -- (check_z.north);
    \draw[flecha] (check_z.south) -- node[right] {Sí} (w_approx.north);
    \draw[flecha] (w_approx.south) -- (use_hint.north);
    \draw[flecha] (use_hint.south) -- (w1_prima.north);
    \draw[flecha] (w1_prima.south) -- (hash.north);
    \draw[flecha] (hash.south) -- (comparacion.north);

    \draw[flecha, red] (check_z.west) -- node[above] {No} (reject1.east);
    
    \draw[flecha, green!50!black] (comparacion.east) -- node[above] {Sí} (accept.west);
    \draw[flecha, red] (comparacion.west) -- node[above] {No} (reject2.east);

    \draw[flecha, dashed] (sigma.west) -| (-8.5, -4.5) -- node[above, near end] {$\tilde{c} \to c$} (w_approx.west);
    
    \draw[flecha, dashed] (mu.south) |- (hash.east);

    \end{tikzpicture}
    }
    \caption{Arquitectura del proceso de verificación.}
    \label{fig:verificacion}
\end{figure}

\section{Consideraciones del Estándar FIPS 204}
\label{sec:fips204_analisis}

El estándar FIPS 204 \cite{nist_fips204} define una arquitectura completa orientada a la implementación segura y resistente frente a criptoanálisis cuántico. A continuación, se desglosan las decisiones de diseño tomadas por el NIST que fundamentan la implementación.

\subsection{Conjuntos de Parámetros y Niveles de Seguridad}
\label{subsec:parametros_niveles}

Para ofrecer resistencia frente a diferentes requisitos de \textit{hardware} y modelos de amenazas, el NIST ha estandarizado tres conjuntos de parámetros. En ML-DSA, la seguridad no escala aumentando el módulo matemático, como pasaba en RSA, sino expandiendo las dimensiones de la matriz del retículo ($k \times l$) y ajustando los umbrales de error.

Cada conjunto está diseñado para igualar o superar la resistencia teórica de los algoritmos simétricos tradicionales estandarizados en FIPS 197 (AES) \cite{nist_fips197}:

\begin{itemize}
    \item \textbf{ML-DSA-44 (Nivel 2):} Equivalente a la dificultad de romper AES-128 mediante fuerza bruta cuántica. Emplea una matriz de $4 \times 4$.
    \item \textbf{ML-DSA-65 (Nivel 3):} Equivalente a la dificultad de romper AES-192. Emplea una matriz de $6 \times 5$. Es el estándar recomendado por defecto por el NIST.
    \item \textbf{ML-DSA-87 (Nivel 5):} Equivalente a la dificultad de romper AES-256. Emplea una matriz de $8 \times 7$. Reservado para sistemas que exijan máxima confidencialidad a muy largo plazo.
\end{itemize}

La selección de estos parámetros impone una compensación crítica entre seguridad teórica y carga de red. La Tabla \ref{tab:tamanos_mldsa} expone la carga de memoria resultante de la serialización.

\begin{table}[H]
    \centering
    \renewcommand{\arraystretch}{1.3}
    \resizebox{0.95\textwidth}{!}{
    \begin{tabular}{lcccc}
        \toprule
        \textbf{Parámetro} & \textbf{Seguridad Equivalente} & \textbf{Clave Pública ($pk$)} & \textbf{Clave Privada ($sk$)} & \textbf{Firma ($\sigma$)} \\
        \midrule
        \textbf{ML-DSA-44} & AES-128 (Nivel 2) & 1.312 bytes & 2.560 bytes & 2.420 bytes \\
        \textbf{ML-DSA-65} & AES-192 (Nivel 3) & 1.952 bytes & 4.032 bytes & 3.309 bytes \\
        \textbf{ML-DSA-87} & AES-256 (Nivel 5) & 2.592 bytes & 4.896 bytes & 4.627 bytes \\
        \bottomrule
    \end{tabular}
    }
    \caption{Requerimientos de almacenamiento y transmisión según FIPS 204.}
    \label{tab:tamanos_mldsa}
\end{table}

\subsection{Función Esponja y SHA-3 - FIPS 202}
\label{subsec:funcion_esponja}

Una revisión profunda de ML-DSA revela que el motor que impulsa las operaciones no pertenece al álgebra de retículos, sino a la familia de funciones hash SHA-3, pertenecientes al FIPS 202 \cite{nist_fips202}, utilizando las variantes \texttt{SHAKE-128} y \texttt{SHAKE-256}.

A diferencia de funciones de resumen tradicionales, como la conocida SHA-256, que reducen una entrada de tamaño arbitrario a una salida de longitud fija, \texttt{SHAKE} es una Función de Salida Extensible (\textbf{XOF}, \textit{eXtendable-Output Function}). Su arquitectura interna se basa en la construcción de esponja \cite{keccak_sponge}, que opera sobre un estado interno de 1600 \textit{bits} dividido en dos partes lógicas:
\begin{enumerate}
    \item \textbf{Tasa ($r$, \textit{Rate}):} La porción del estado que interactúa directamente con los datos de entrada/salida. Determina el rendimiento del algoritmo.
    \item \textbf{Capacidad ($c$, \textit{Capacity}):} La porción oculta del estado. Los datos nunca interactúan directamente con ella. Define el nivel de seguridad contra ataques de colisión y preimagen ($c = 256$ para \texttt{SHAKE-128}).
\end{enumerate}

La ejecución se divide en dos fases iterativas, ilustradas en la Figura \ref{fig:esponja_shake}:

\begin{figure}[H]
    \centering
    \begin{tikzpicture}[
        bloque/.style = {draw, rectangle, minimum width=2.5cm, minimum height=0.8cm, fill=white, thick},
        func/.style = {draw, rectangle, rounded corners, minimum width=2.8cm, minimum height=1.2cm, fill=blue!5, thick, align=center},
        flecha/.style = {-{Stealth[scale=1.2]}, thick},
        etiqueta/.style = {font=\bfseries\large, align=center}
    ]

    \node[func] (f1) at (0, 0) {Permutación $f$};
    \node[func] (f2) at (0, -3) {Permutación $f$};
    \node[func] (f3) at (0, -7) {Permutación $f$};
    \node[func] (f4) at (0, -10) {Permutación $f$};

    \node[bloque] (m1) at (-4.8, 0) {Bloque Datos $m_1$};
    \node[bloque] (m2) at (-4.8, -3) {Bloque Datos $m_2$};

    \node[bloque, fill=green!5] (z1) at (4.8, -7) {Salida $z_1$};
    \node[bloque, fill=green!5] (z2) at (4.8, -10) {Salida $z_2$};

    \draw[flecha] (0, 1.5) -- node[right] {Estado Inicial (0)} (f1.north);
    \draw[flecha] (f1.south) -- node[right] {Estado 1600b} (f2.north);
    
    \draw[flecha, dashed] (f2.south) -- node[right, fill=white, inner sep=4pt] {Transición (\textit{Padding})} (f3.north);
    
    \draw[flecha] (f3.south) -- node[right] {Estado 1600b} (f4.north);
    \draw[flecha, dashed] (f4.south) -- (0, -11.5) node[right] {\dots};

    \draw[flecha] (m1.east) -- node[midway, above, font=\small] {XOR en $r$} (f1.west);
    \draw[flecha] (m2.east) -- node[midway, above, font=\small] {XOR en $r$} (f2.west);
    
    \draw[flecha] (f3.east) -- node[midway, above, font=\small] {Extraer $r$} (z1.west);
    \draw[flecha] (f4.east) -- node[midway, above, font=\small] {Extraer $r$} (z2.west);

    \draw[thick, loosely dashed, gray] (-6, -5) -- (6, -5);
    
    \node[etiqueta] at (-3, -5) [above=0.4cm] {Fase 1: Absorber (\textit{Absorb})};
    \node[etiqueta] at (3, -5) [below=0.4cm] {Fase 2: Exprimir (\textit{Squeeze})};

    \end{tikzpicture}
    \caption{Construcción de Esponja en FIPS 202.}
    \label{fig:esponja_shake}
\end{figure}

\subsubsection{Ejemplo Práctico en ML-DSA: Muestreo de Rechazo}

La necesidad estricta de utilizar una función XOF en lugar de un \textit{hash} clásico se evidencia en la generación de la matriz pública $\mathbf{\hat{A}}$. 

ML-DSA requiere generar polinomios cuyos coeficientes sean enteros uniformemente distribuidos en el rango $[0, q-1]$, donde $q = 8380417$. Para lograr esto, se utiliza la semilla pública $\rho$ (32 \textit{bytes}). Si se utilizara SHA-256, la semilla produciría únicamente 32 \textit{bytes} fijos, insuficientes para generar todos los coeficientes de los polinomios de la matriz $A$.

Con \texttt{SHAKE-128}, la semilla $\rho$ se \textbf{absorbe} en el estado. A continuación, el algoritmo comienza a \textbf{exprimir} \textit{bytes} aleatorios. Para cada coeficiente, ML-DSA toma 3 \textit{bytes} (24 \textit{bits}) del flujo XOF, anula el \textit{bit} más significativo (dejando 23 \textit{bits}) y evalúa el entero resultante ($x$):
\begin{itemize}
    \item Si $x < q$, el coeficiente es aceptado.
    \item Si $x \ge q$, el coeficiente es \textbf{rechazado} y se descarta.
\end{itemize}
Debido a este muestreo de rechazo, es matemáticamente imposible predecir cuántos \textit{bytes} se necesitarán para completar la matriz $\mathbf{\hat{A}}$. La función XOF resuelve este problema al permitir que el bucle continúe exprimiendo el estado de Keccak indefinidamente hasta que todos los coeficientes válidos hayan sido generados, garantizando la seguridad en la heurística de Fiat-Shamir.

\subsection{Variantes Pre-Hash}
\label{subsec:prehash_oids}

El estándar FIPS 204 define la versión pura del algoritmo, requiriendo que $M$ sea inyectada directamente en la fase de absorción de la función SHAKE durante el proceso de firma. Aunque matemáticamente es el proceso ideal, desde la perspectiva de la ingeniería de \textit{software} presenta un punto crítico debido al agotamiento de memoria.

Si un sistema de producción necesita firmar una imagen de sistema de 2 GB o una base de datos masiva, la API pura de ML-DSA obligaría a cargar todo el flujo de \textit{bytes} en la memoria RAM de forma ininterrumpida. Para solucionar este impedimento y permitir la compatibilidad con protocolos heredados como CMS o X.509, el NIST ha estandarizado la variante \textbf{HashML-DSA} \cite{ietf_mldsa_pki}.

En esta variante, el sistema procesa el mensaje de forma independiente mediante una función de resumen criptográfico tradicional (como SHA-256 o SHA-512) para generar un \textit{hash} previo, denotado como $PH(M)$. ML-DSA firma exclusivamente este bloque de datos reducido.

\subsubsection{Mitigación de Colisiones: El contexto OID}

Delegar el proceso a un algoritmo externo introduce un riesgo de seguridad conocido como sustitución algorítmica. Un atacante podría encontrar una colisión en un algoritmo de \textit{hashing} obsoleto y hacer pasar ese resumen malicioso como legítimo.

Para blindar la variante HashML-DSA, FIPS 204 prohíbe firmar el \textit{hash} desnudo. El estándar obliga a encapsular $PH(M)$ dentro de una estructura de contexto (\texttt{ctx}) que incluye, de forma inmutable, el Identificador de Objeto del algoritmo exacto utilizado.

Si el sistema pre-procesa el mensaje utilizando SHA-512, la función interna de ML-DSA concatenará la secuencia hexadecimal \texttt{06 09 60 86 48 01 65 03 04 02 03} al resumen antes de firmarlo. Esta unión a nivel de \textit{bytes} garantiza que la validación criptográfica fracase si un servidor mal configurado intenta verificar el documento utilizando un algoritmo distinto al empleado originalmente por el emisor.

\section{Ejecución de ML-DSA con Dimensión Reducida}
\label{sec:instanciacion_juguete}

Para consolidar la comprensión del flujo algebraico de ML-DSA sin la sobrecarga computacional que imponen las dimensiones reales del estándar, en esta sección se desarrolla una ejecución numérica completa utilizando un modelo de escala reducida. Los parámetros seleccionados corresponden a los esquemas didácticos propuestos por Menezes \cite{menezes_slides}, los cuales conservan la estructura algebraica del problema \textit{Module Learning With Errors} MLWE pero reducen drásticamente las dimensiones del anillo y los vectores.

\subsection{Definición del Entorno y Parámetros del Sistema}

A diferencia de las instancias de producción de FIPS 204, donde el grado del polinomio se fija estrictamente en $n=256$ y el módulo es el primo $q=8380417$, este modelo simplificado opera sobre un espacio vectorial acotado. Las constantes fundamentales de este ejemplo se resumen en la Tabla~\ref{tab:parametros_juguete}.

\begin{table}[H]
\centering
\caption{Parámetros del modelo de ejemplo de ML-DSA.}
\label{tab:parametros_juguete}
\begin{tabular}{lcc}
\hline
\textbf{Parámetro Criptográfico} & \textbf{Símbolo} & \textbf{Valor Asignado} \\ \hline
Módulo del cuerpo base & $q$ & $16417$ \\
Grado del anillo polinómico & $n$ & $4$ \\
Dimensión de filas de la matriz (MLWE) & $k$ & $3$ \\
Dimensión de columnas de la matriz (MLWE) & $\ell$ & $2$ \\
Cota de los vectores secretos ($\mathbf{s}_1, \mathbf{s}_2$) & $\eta$ & $10$ \\
Parámetro de enmascaramiento local & $\gamma_1$ & $2^{10} = 1024$ \\
Parámetro de enmascaramiento de bajo orden & $\gamma_2$ & $(q-1)/32 = 513$ \\
Peso de dispersión del desafío $c$ & $\tau$ & $4$ \\
\hline
\end{tabular}
\end{table}

El espacio de trabajo queda definido por el anillo cociente de polinomios:

\begin{equation}
    R_q = \mathbb{Z}_q[x] / (x^4 + 1) = \mathbb{Z}_{16417}[x] / (x^4 + 1)
\end{equation}

Cualquier elemento elemento del anillo es un polinomio de la forma $a_0 + a_1x + a_2x^2 + a_3x^3$, donde los coeficientes pertenecen al cuerpo finito $\mathbb{X}_{16417}$. Es importante mostrar que dado que $n = \tau = 4$, el algoritmo \texttt{SampleInBall} generará un polinomio de desafío $c$ donde absolutamente todos sus coeficientes serán no nulos ($\pm 1$), eliminando los coeficientes cero habituales de las instancias reales.

\subsection{Cotas del Muestreo de Rechazo}

Como se analizó en la Sección \ref{subsec:respuesta_z_rechazo}, el éxito del protocolo \textit{Fiat-Shamir con abortos}\cite{lyubashevsky2009fiat} depende de que las componentes de la firma no filtren información sobre los secretos. Utilizando los parámetros:

\begin{equation}
    \beta = \tau \cdot \eta = 4 \cdot 10 = 40
\end{equation}

Los umbrales críticos de aceptación que controlarán los filtros de rechazo del algoritmo se derivan de la siguiente manera:

\begin{itemize}
    \item \textbf{Filtro sobre el vector de respuesta $\mathbf{z}$:} El vector no debe aproximarse a los límites de la cota de enmascaramiento:
    \begin{equation}
        ||\mathbf{z}||_\infty < \gamma_1 - \beta \implies 1024 - 40 = 984
    \end{equation}
    
    \item \textbf{Filtro sobre los bits de bajo orden $\mathbf{r}_0$:} El cambio de valores de la reducción no deben alterar el bit de signo del compromiso:
    \begin{equation}
        ||\mathbf{r}_0||_\infty < \gamma_2 - \beta \implies 513 - 40 = 473
    \end{equation}
\end{itemize}

Si durante la ejecución del algoritmo cualquiera de las componentes calculadas extrae una norma infinito que iguale o supere los límites $[984, 473]$, la firma apta se descartará inmediatamente reiniciando el bucle.

\subsection{Fase 1: Generación de Claves}

El primer paso del consiste en la instanciación de las variables del sistema: la matriz pública $\mathbf{A}$ y los vectores secretos $\mathbf{s}_1$ y $\mathbf{s}_2$.

La matriz $\mathbf{A} \in R_q^{k \times \ell}$ (con dimensiones $3 \times 2$) se genera utilizando la función $\text{Expand}(\mathbf{A})$, la cual construye la matriz sobre el anillo $\mathbb{X}_{16417}$ con los parámetros del sistema:
\begin{equation}
    \mathbf{A} = \text{genA}(R_q, k, \ell)
\end{equation}

La ejecución de esta instrucción produce el siguiente resultado numérico:
\begin{equation}
    \mathbf{A} = \begin{pmatrix}
    7120x^3 + 14306x^2 + 6781x + 6110 & 1568x^3 + 6180x^2 + 2401x + 1490 \\
    4392x^3 + 15280x^2 + 2900x + 13812 & 15082x^3 + 2013x^2 + 3956x + 3554 \\
    5174x^3 + 6749x^2 + 1040x + 6670 & 11871x^3 + 7903x^2 + 8679x + 1088
    \end{pmatrix}
\end{equation}

Por su parte, los vectores secretos $\mathbf{s}_1$ y $\mathbf{s}_2$ se instancian utilizando la función generadora de secretos del algoritmo. Esta función recibe como argumento la cota $n_1 = 10$ (equivalente al parámetro $\eta$ de la especificación oficial) para garantizar el enmascaramiento:
\begin{equation}
    \mathbf{s}_1 = \text{ExpandS}(\rho', \ell)
\end{equation}
\begin{equation}
    \mathbf{s}_2 = \text{ExpandS}(\rho', k)
\end{equation}

Los vectores secretos $\mathbf{s}_1$ y $\mathbf{s}_2$ se derivan de la expansión de la semilla $\rho'$ mediante la función \texttt{ExpandS}. En la implementación de este modelo, dicha operación se instancia a través de la función generadora \texttt{genSnk}. Esta rutina recibe como argumento la cota $n_1 = 10$ (equivalente al parámetro $\eta$ de la especificación oficial) para garantizar el enmascaramiento:
\begin{equation}
    \mathbf{s}_1 = \begin{pmatrix}
    7x^3 - 7x^2 - 9x - 6 \\
    -5x^3 + 10x^2 - 3x - 3
    \end{pmatrix}
\end{equation}
\begin{equation}
    \mathbf{s}_2 = \begin{pmatrix}
    -2x^3 - 8x^2 + 6x - 5 \\
    -5x^3 + 8x^2 + 3x + 5 \\
    -9x^3 - 9x - 5
    \end{pmatrix}
\end{equation}

Para terminar la fase de generación de claves, se calcula el vector de la clave pública $\mathbf{t} \in R_q^k$. Esta operación se ejecuta obligatoriamente en el dominio NTT para optimizar el rendimiento de las multiplicaciones polinómicas, aplicando posteriormente una compresión de bits. Sin embargo, este modelo permite realizar la convolución directamente en el anillo $R_q$ utilizando la ecuación del problema MLWE:

\begin{equation}
    \mathbf{t} = \mathbf{A} \cdot \mathbf{s}_1 + \mathbf{s}_2 \pmod{q}
\end{equation}

Considerando las dimensiones acotadas del sistema ($k=3, \ell=2$), la operación se expande en el siguiente sistema:

\begin{equation}
    \begin{pmatrix} t_1 \\ t_2 \\ t_3 \end{pmatrix} = 
    \begin{pmatrix} 
    A_{11} & A_{12} \\ 
    A_{21} & A_{22} \\ 
    A_{31} & A_{32} 
    \end{pmatrix} 
    \begin{pmatrix} s_{1,1} \\ s_{1,2} \end{pmatrix} + 
    \begin{pmatrix} s_{2,1} \\ s_{2,2} \\ s_{2,3} \end{pmatrix} \pmod{16417}
\end{equation}

La ejecución de esta instrucción en el entorno (\texttt{t = A*s1 + s2}) resuelve el producto escalar de las filas de la matriz pública por el vector secreto, aplicando la reducción de ($x^4 \equiv -1$) y el módulo $q = 16417$. 

A diferencia de los vectores que representamos de forma simétrica para comprobar que respetan la cota $\eta$, el vector $\mathbf{t}$ es público. Al no requerir esta comprobación de seguridad, se extrae y se empaqueta directamente con los valores positivos del módulo $16417$:

\begin{equation}
    \mathbf{t} = \begin{pmatrix}
    14147x^3 + 3769x^2 + 5328x + 14112 \\
    14348x^3 + 12013x^2 + 12802x + 11887 \\
    2478x^3 + 2646x^2 + 11204x + 10384
    \end{pmatrix}
\end{equation}

Con la obtención de $\mathbf{A}$ y $\mathbf{t}$ queda formada la clave pública del sistema $pk$, mientras que los vectores $\mathbf{s}_1$ y $\mathbf{s}_2$ se custodian como la clave privada $sk$.


Las variables instanciadas se empaquetan en listas (\texttt{pk = [A, t]} y \texttt{sk = [s1, s2]}) para ser transferidas a las siguientes fases. El par de claves generado se representa de la siguiente manera:

\begin{equation}
    pk = (\mathbf{A}, \mathbf{t}) \in R_q^{k \times \ell} \times R_q^k
\end{equation}

\begin{equation}
    sk = (\mathbf{s}_1, \mathbf{s}_2) \in S_\eta^\ell \times S_\eta^k
\end{equation}

Con este empaquetado concluye la fase de generación de claves.


\subsection{Fase 2: Generación de la Firma}

\subsubsection*{Paso 0: Preparación del Mensaje}
El proceso de firma comienza con la preparación del mensaje $M$. En la especificación FIPS 204, los mensajes son procesados inicialmente mediante la función hash SHAKE-256 para obtener una cadena de longitud fija. Pero para simplificar la ejecución en este ejemplo, se aplica un mapeo directo donde los bits del mensaje estructuran los coeficientes de un polinomio en $R_q$. 

Dada la dimensión del anillo $n=4$, se define un mensaje binario de 4 bits ($M = \texttt{'1001'}$), el cual se codifica en el sistema de la siguiente manera:
\begin{equation}
    M(x) = 1 \cdot x^3 + 0 \cdot x^2 + 0 \cdot x^1 + 1 \cdot x^0 = x^3 + 1
\end{equation}

\subsubsection*{Paso 1: Muestreo del Vector de Enmascaramiento $\mathbf{y}$}
Una vez estructurado el mensaje, el protocolo inicia el bucle de firma con el primer paso: la instanciación del vector de enmascaramiento $\mathbf{y} \in R_q^\ell$. Este vector debe muestrearse asegurando que todos sus coeficientes estén estrictamente acotados por el parámetro $\gamma_1 = 1024$.

Para ello, el estándar oficial utiliza la función \texttt{ExpandMask}, la cual deriva el vector a partir de la semilla privada $\rho'$ y un contador de iteraciones $\kappa$. Para mantener la coherencia con la ejecución de este modelo y evidenciar su dimensión estructural, la llamada se formula de la siguiente manera, asegurando la cota $\gamma_1$:
\begin{equation}
    \mathbf{y} = \text{ExpandMask}(\rho', \kappa, \ell)
\end{equation}

El sistema algebraico devuelve el vector utilizando los representantes positivos por defecto de $\mathbb{Z}_{16417}$:
\begin{equation}
    \mathbf{y} = \begin{pmatrix}
    15572x^3 + 15677x^2 + 16198x + 672 \\
    15973x^3 + 15952x^2 + 219x + 575
    \end{pmatrix}
\end{equation}

Para verificar que la función ha respetado el límite criptográfico $\gamma_1$, se extrae la representación simétrica del vector. Se comprueba que, efectivamente, todos los coeficientes pertenecen al intervalo $[-1024, 1024]$:
\begin{equation}
    \mathbf{y} = \begin{pmatrix}
    -845x^3 - 740x^2 - 219x + 672 \\
    -444x^3 - 465x^2 + 219x + 575
    \end{pmatrix}
\end{equation}

\subsubsection*{Paso 2: Cálculo del Compromiso Comprimido $\mathbf{w}_1$}
El segundo paso del bucle de firma consiste en calcular el vector intermedio de compromiso $\mathbf{w} \in R_q^k$. Esta operación constituye una multiplicación matricial sobre el anillo:
\begin{equation}
    \mathbf{w} = \mathbf{A} \cdot \mathbf{y} \pmod{q}
\end{equation}

La ejecución de este producto da como resultado el siguiente vector, evaluado con los coeficientes positivos del módulo $16417$:
\begin{equation}
    \mathbf{w} = \begin{pmatrix}
    8478x^3 + 14476x^2 + 14181x + 3317 \\
    3868x^3 + 12517x^2 + 8218x + 4380 \\
    8513x^3 + 14058x^2 + 10825x + 1782
    \end{pmatrix}
\end{equation}

Para evitar que la firma final filtre información sensible sobre la clave privada, y para reducir la longitud de los datos transmitidos, el estándar no utiliza el vector $\mathbf{w}$ completo. En su lugar, el algoritmo descompone cada coeficiente en sus partes de mayor y menor peso mediante la función \texttt{HighBits}.

La partición del espacio se rige por el parámetro de $\gamma_2 = 513$, lo que define un tamaño de bloque de $\alpha = 2\gamma_2 = 1026$. Aplicando la extracción de los bits más significativos a cada coeficiente del vector, se obtiene el compromiso comprimido $\mathbf{w}_1$:
\begin{equation}
    \mathbf{w}_1 = \begin{pmatrix}
    8x^3 + 14x^2 + 14x + 3 \\
    4x^3 + 12x^2 + 8x + 4 \\
    8x^3 + 14x^2 + 11x + 2
    \end{pmatrix}
\end{equation}

\subsubsection*{Paso 3: Generación del Desafío $c$}
El tercer paso del bucle es la generación del polinomio de desafío $c \in B_\tau$. Este elemento vincula el mensaje original con el compromiso mediante el paradigma de Fiat-Shamir.

En la especificación oficial, esta vinculación se realiza instanciando la función \textit{Sponge} SHAKE-256 para absorber el mensaje y la compresión $\mathbf{w}_1$, exprimiendo $\tilde{c}$ de longitud fija. Para este entorno de validación, la capa de \textit{hashing} se ha instanciado utilizando la función estándar SHA-256. El sistema concatena el mensaje binario con la representación en cadena del vector $\mathbf{w}_1$, obteniendo el siguiente digesto hexadecimal:

\begin{equation}
    \tilde{c} = \text{SHA-256}(M \parallel \mathbf{w}_1) = \texttt{bf022358b6e3...394ec73f6493806bf3}
\end{equation}

Para transformarlo en un elemento del anillo válido, se interpreta la cadena hexadecimal como un número entero que actúa como semilla para el algoritmo \texttt{SampleInBall}. Esta función es la responsable de generar un polinomio en $R_q$ que pertenezca a $B_\tau$, es decir, un polinomio que contenga exactamente $\tau$ coeficientes no nulos (evaluados como $\pm 1$) y el resto a cero.
\begin{equation}
    c = \texttt{SampleInBall}(\tilde{c}, n, \tau)
\end{equation}

La ejecución en el sistema devuelve inicialmente los coeficientes evaluados como enteros positivos módulo $q$:
\begin{equation}
    c = x^3 + 16416x^2 + x + 16416
\end{equation}

Para verificar el cumplimiento del estándar, se reescribe el polinomio utilizando su representación simétrica en el origen ($16416 \equiv -1 \pmod{16417}$). Este modelo de prueba donde el grado del anillo ($n=4$) coincide con la constante de dispersión ($\tau=4$), obliga a que todos los coeficientes sean no nulos, confirmándose el siguiente resultado:
\begin{equation}
    c = x^3 - x^2 + x - 1
\end{equation}

\subsubsection*{Paso 4: Cálculo de la Firma Apta y Muestreo de Rechazo}

En este paso se calcula el vector de firma apta $\mathbf{z} \in R_q^\ell$ combinando el vector de enmascaramiento aleatorio con el secreto $\mathbf{s}_1$ por el desafío:
\begin{equation}
    \mathbf{z} = \mathbf{y} + c \cdot \mathbf{s}_1 \pmod{q}
\end{equation}

La ejecución de estas operaciones \texttt{z = y + c*s1} genera el siguiente resultado en base a los coeficientes positivos del módulo $16417$:
\begin{equation}
    \mathbf{z} = \begin{pmatrix}
    15561x^3 + 15674x^2 + 16215x + 673 \\
    15988x^3 + 15947x^2 + 204x + 596
    \end{pmatrix}
\end{equation}

Para evaluar si este candidato compromete la seguridad o si se ha desviado de los límites permitidos, se reduce a su representación simétrica:
\begin{equation}
    \mathbf{z} = \begin{pmatrix}
    -856x^3 - 743x^2 - 202x + 673 \\
    -429x^3 - 470x^2 + 204x + 596
    \end{pmatrix}
\end{equation}

\textbf{Verificación de Condiciones de Rechazo}

Para garantizar que la firma no filtre información sobre la clave privada y que la verificación posterior sea consistente, el estándar aplica un muestreo de rechazo. Si el candidato no supera alguna de las siguientes cotas de seguridad, la firma se descarta automáticamente, se aborta el bucle y se vuelve al Paso 1 para generar un nuevo vector $\mathbf{y}$.

\textbf{Check 1: Cota del vector de firma}\\
Se comprueba si la norma infinito del vector supera o iguala el umbral seguro establecido por $\gamma_1 - \beta$, donde para este modelo el margen de rechazo es $\beta = \tau \cdot n_1 = 4 \cdot 10 = 40$:
\begin{equation}
    ||\mathbf{z}||_\infty \ge \gamma_1 - \beta \implies ||\mathbf{z}||_\infty \ge 1024 - 40 = 984
\end{equation}
El coeficiente máximo absoluto de $\mathbf{z}$ es $856$. Evaluando la condición en el sistema (\texttt{size2(z) >= g\_1 - beta}), el resultado es \texttt{False}, superando con éxito el primer filtro.

\textbf{Check 2: Cota de los bits de bajo orden}\\
Se evalúa si la parte menos significativa del error introducido en el compromiso supera el umbral para evitar errores. Para ello, se calcula el vector de bits de bajo orden o menos significativos $\mathbf{w}_0 = \text{LowBits}(\mathbf{A}\mathbf{y} - c\mathbf{s}_2, 2\gamma_2)$:
\begin{equation}
    ||\mathbf{w}_0||_\infty \ge \gamma_2 - \beta \implies ||\mathbf{w}_0||_\infty \ge 513 - 40 = 473
\end{equation}

La extracción de estos residuos en el entorno devuelve el siguiente vector de polinomios:
\begin{equation}
    \mathbf{w}_0 = \begin{pmatrix}
    287x^3 + 91x^2 + 16239x + 246 \\
    16166x^3 + 210x^2 + 21x + 271 \\
    292x^3 + 16106x^2 + 15961x + 16124
    \end{pmatrix}
\end{equation}

Llevando este vector a su forma simétrica para medir su magnitud real, se obtiene:
\begin{equation}
    \mathbf{w}_0 = \begin{pmatrix}
    287x^3 + 91x^2 - 178x + 246 \\
    -251x^3 + 210x^2 + 21x + 271 \\
    292x^3 - 311x^2 - 456x - 293
    \end{pmatrix}
\end{equation}

El coeficiente absoluto máximo de este vector es $456$. Al evaluar la condición matemática (\texttt{size2(lb) >= g\_2 - beta}), el sistema confirma que $456 \ge 473$ es \texttt{False}.

Dado que ambas condiciones han resultado negativas, la firma potencial no se aborta y se considera geométricamente segura para su empaquetado final.


\subsubsection*{Paso 5: Empaquetado de la Firma}

Una vez superados los filtros de seguridad del muestreo de rechazo, el vector $\mathbf{z}$ se asume como seguro y definitivo. Para concluir la fase de generación, el algoritmo empaqueta este vector junto con el polinomio de desafío $c$. 

La firma digital resultante $\sigma$ para el mensaje $M$ queda definida de siguiente forma:
\begin{equation}
    \sigma = (c, \mathbf{z})
\end{equation}

Con la emisión de $\sigma$, el bucle finaliza con éxito y la firma queda lista para ser transmitida junto al mensaje original.


\subsection{Fase 3: Verificación de la Firma}

El proceso de verificación permite a cualquier entidad que posea la clave pública $pk = (\mathbf{A}, \mathbf{t})$ comprobar la autenticidad del mensaje $M$ y la integridad de la firma $\sigma = (c, \mathbf{z})$, sin necesidad de conocer los vectores secretos. El algoritmo se considera exitoso si, y solo si, se cumplen dos condiciones simultáneas.

\subsubsection*{Paso 1: Reconstrucción del Compromiso $\mathbf{w}_1'$}
Se utiliza la clave pública y la firma proporcionada para intentar reconstruir la parte de bits más significativos del vector de compromiso original. Se calcula el vector candidato $\mathbf{w}_1'$ mediante la siguiente operación:
\begin{equation}
    \mathbf{w}_1' = \text{HighBits}(\mathbf{A} \cdot \mathbf{z} - c \cdot \mathbf{t} \pmod{q}, 2\gamma_2)
\end{equation}

La ejecución de \texttt{w2 = HighBitsV(...)} devuelve el siguiente vector de polinomios:
\begin{equation}
    \mathbf{w}_1' = \begin{pmatrix}
    8x^3 + 14x^2 + 14x + 3 \\
    4x^3 + 12x^2 + 8x + 4 \\
    8x^3 + 14x^2 + 11x + 2
    \end{pmatrix}
\end{equation}
Como se puede observar, este resultado coincide de forma exacta con el vector comprimido $\mathbf{w}_1$ generado por el firmante durante el Paso 2 de la fase anterior.

\subsubsection*{Paso 2: Comprobación del Límite de Seguridad}
Bob no confía en la firma recibida. Antes de aceptar el desafío, debe asegurar que el vector $\mathbf{z}$ no ha sido manipulado y cumple con los requisitos de enmascaramiento. Se evalúa la siguiente condición:
\begin{equation}
    ||\mathbf{z}||_\infty < \gamma_1 - \beta
\end{equation}
Considerando que el coeficiente absoluto máximo de $\mathbf{z}$ en su forma simétria es $856$, y que el umbral máximo permitido es $984$, el sistema valida la inecuación (\texttt{True}). La firma es segura.

\subsubsection*{Paso 3: Validación del Desafío}
La segunda condición es comprobar si el desafío adjunto en la firma $c$ se corresponde realmente con el mensaje $M$ y el compromiso reconstruido $\mathbf{w}_1'$. Para ello, Bob repite el proceso de empaquetado y \textit{hashing}:
\begin{equation}
    c' = \texttt{SampleInBall}(\text{SHA-256}(M \parallel \mathbf{w}_1'), n, \tau)
\end{equation}
El cálculo da como resultado el siguiente polinomio, evaluado con los representantes positivos en $\mathbb{Z}_{16417}$:
\begin{equation}
    c' = x^3 + 16416x^2 + x + 16416
\end{equation}
Puesto que $c' = c$ y $||\mathbf{z}||_\infty < \gamma_1 - \beta$, ambas condiciones se satisfacen. El algoritmo concluye que \textbf{la firma es completamente válida}.

\subsubsection*{Modificación del Mensaje}
Para demostrar la robustez del esquema, se simula un escenario donde el mensaje es alterado. Se modifica el mensaje original $M = \texttt{'1001'}$ por un mensaje fraudulento $M^* = \texttt{'1111'}$.

Al intentar verificar la firma original $\sigma$ contra el mensaje manipulado $M^*$, el efecto avalancha de la función \textit{hash} provoca que la concatenación $M^* \parallel \mathbf{w}_1'$ produzca un elemento completamente distinto. Como resultado, la función \texttt{SampleInBall} genera un nuevo desafío $c^*$:
\begin{equation}
    c^* = x^3 + x^2 + 16416x + 1
\end{equation}

Los coeficientes muestran la diferencia ($16416 \equiv -1 \pmod{q}$), se observa:
\begin{equation}
    c^* = x^3 + x^2 - x + 1 \neq c
\end{equation}

Al evaluar $c^* == c$, el sistema devuelve \texttt{False}. El algoritmo detecta la diferencia entre el mensaje y el compromiso, \textbf{rechazando la firma automáticamente}. Esto muestra la resistencia del protocolo ML-DSA frente a manipulaciones externas.














\chapter{Implementación del Motor Criptográfico en Python}
\label{ch:implementacion_motor}

La seguridad en un esquema criptográfico post-cuántico no reside únicamente en la matemática teórica, sino en la precisión y fidelidad de su implementación en \textit{software}. Un algoritmo basado en retículos puede colapsar en producción si su código fuente introduce errores de información, desbordamientos de enteros o una mala gestión de la entropía del sistema.

En este capítulo se detalla la arquitectura, el diseño y la validación de la Prueba de Concepto PoC desarrollada para este proyecto: un motor criptográfico ML-DSA funcional, modular y auditable.

\section{Entorno de Desarrollo y Metodología}
\label{sec:entorno_metodologia}

El desarrollo de software criptográfico sigue reglas distintas a las del software convencional. En este campo, el principio fundamental es evitar la creación de algoritmos propios, ya que la seguridad reside en la robustez de estándares probados. 

Por esta razón, el objetivo de esta implementación no ha sido la innovación en el diseño del algoritmo, sino garantizar una \textbf{transcripción fiel y segura del estándar} descrito por el NIST en el FIPS 204\cite{nist_fips204}.

\subsection{Justificación del Lenguaje: Python vs. Rendimiento}
\label{subsec:justificacion_python}

El estándar ML-DSA requiere cálculos matemáticos complejos, especialmente durante las multiplicaciones polinómicas en el dominio de la Transformada Numérica Teórica NTT y el muestreo de rechazo. Este tipo de operaciones se suelen implementar en lenguajes de bajo nivel muy cercanos al \textit{hardware}, como C o Rust, para optimizar el rendimiento de la CPU y el consumo de memoria.

Sin embargo, para el desarrollo de esta Prueba de Concepto, se seleccionó \textbf{Python} como lenguaje principal, asumiendo un sacrificio consciente en los tiempos de ejecución a favor de tres ventajas para el alcance de este trabajo:

\begin{enumerate}
    \item \textbf{Mapeo Semántico 1:1:} Python prioriza la legibilidad del código. Esto permitió que la sintaxis de la implementación fuera un reflejo casi literal del pseudocódigo en el FIPS 204. Las variables, bucles y operaciones polinómicas en el código fuente pueden ser auditadas visualmente con las especificaciones del NIST sin la preocupación por la gestión manual de punteros o la reserva explícita de memoria como pasa en el lenguaje C.
    \item \textbf{Aritmética de Precisión Arbitraria Nativa:} A diferencia de C, donde el desbordamiento de enteros  en la aritmética modular del anillo $R_q$ requiere controles exhaustivos, el modelo de enteros de Python maneja precisión de forma nativa, reduciendo una de las superficies de error más comunes en criptografía.
    \item \textbf{Ecosistema de Validación:} Python dispone de herramientas para la rápida implementación de pruebas unitarias, facilitando el control de calidad continuo del motor criptográfico.
\end{enumerate}

\subsection{Metodología: Traducción FIPS y TDD en Capas}
\label{subsec:metodologia_tdd}

Para garantizar que la implementación no introdujera vulnerabilidades lógicas y superara las pruebas oficiales, el ciclo de vida del \textit{software} se rigió por dos metodologías:

\begin{itemize}
    \item \textbf{Traducción Basada en el Estándar:} Cada función del motor se programó con el documento FIPS 204 como única fuente. No se aplicaron modificaciones que pudieran crear un flujo de ejecución distinto al estandarizado por el NIST, garantizando así una trazabilidad absoluta entre la teoría y la práctica.
    \item \textbf{Desarrollo Guiado por Pruebas TDD Ascendente:} Dada la complejidad matemática de ML-DSA, donde un error en el redondeo del bit menos significativo invalida toda la firma digital, se aplicó un enfoque TDD, es decir, un \textit{Test-Driven Development}. Se creó el grafo de dependencias de todos los subalgoritmos y se estructuró el desarrollo en un modelo ascendente \textit{bottom-up}, clasificado desde el Nivel 0 (primitivas aritméticas) hasta el Nivel 10 (API de firma).
\end{itemize}

Bajo esta metodología, ningún algoritmo de un nivel superior fue implementado hasta que todas sus dependencias inferiores estuvieran programadas y validadas mediante \textit{tests} unitarios aislados. Por ejemplo, el módulo de generación de claves (\texttt{keygen}) no se inició hasta certificar la exactitud matemática de la Transformada Numérica Teórica NTT


\section{Arquitectura del Repositorio y Jerarquía de Dependencias}
\label{sec:arquitectura_repo}

El código fuente ha sido diseñado aplicando el Principio de Responsabilidad Única, para lograr un alto nivel de desacoplamiento. Esto permite aislar de forma estricta la lógica matemática pura de la lógica de integración, facilitando las auditorías de seguridad y las pruebas unitarias.

\subsection{El Paquete Principal y la Suite de Pruebas}

El repositorio se divide en dos grandes bloques que actúan de manera coordinada: el paquete principal (\texttt{mldsa/}) y la suite de validación (\texttt{tests/}). 

El directorio \texttt{mldsa/} constituye el núcleo del motor criptográfico y se subdivide en los siguientes submódulos especializados:

\begin{itemize}
    \item \textbf{\texttt{mldsa.py} (\textit{API Pública}):} Define la interfaz de alto nivel que consume el usuario final. Contiene las funciones \texttt{keygen}, \texttt{sign}, \texttt{verify}, y sus variantes asíncronas para documentos masivos (\texttt{hash\_sign} y \texttt{hash\_verify}).
    \item \textbf{\texttt{core/}:} Contiene los métodos principales del protocolo. Implementa los algoritmos de generación, firma y verificación siguiendo estrictamente los flujos de control del estándar FIPS 204.
    \item \textbf{\texttt{ntt/}:} Módulo de aceleración algebraica. Implementa la Transformada Numérica Teórica NTT para garantizar multiplicaciones polinómicas eficientes.
    \item \textbf{\texttt{crypto/} y \texttt{sampling/}:} Centralizan la generación de entropía. El primero consume las primitivas \textit{hash} (SHAKE-128/256), mientras que el segundo aplica el muestreo por rechazo para generar distribuciones estadísticas seguras.
    \item \textbf{\texttt{encoding/} y \texttt{decomposition/}:} Responsables de la serialización (\textit{bitpacking}) y la gestión del error, incluyendo algoritmos como \texttt{Power2Round} y la división modular de componentes altos y bajos.
    \item \textbf{\texttt{parameters/} y \texttt{primitives/}:} Definen las constantes inmutables para los niveles de seguridad y las operaciones aritméticas elementales sobre el anillo $R_q$.
\end{itemize}

Para garantizar la seguridad, el directorio \texttt{tests/} replica esta estructura, como por ejemplo \texttt{test\_core/}, \texttt{test\_ntt/}, etc., aislando las pruebas de integración de las pruebas matemáticas.

\subsection{Grafo de Dependencias}

Bajo la metodología TDD ascendente descrita anteriormente, el código no fue programado de forma lineal, sino categorizado en una estricta jerarquía de dependencias. Se establecieron 11 niveles lógicos de implementación, del Nivel 0 al Nivel 10. 

Ningún módulo de una capa superior puede acceder a los datos de una capa inferior sin pasar por las interfaces definidas, evitando referencias circulares. La Figura \ref{fig:arquitectura_capas} ilustra este diseño en cascada inversa \textit{bottom-up}:

\begin{figure}[H]
    \centering
    \resizebox{0.9\textwidth}{!}{
    \begin{tikzpicture}[
        capa_base/.style = {draw, rectangle, rounded corners, minimum width=12cm, minimum height=1.2cm, fill=gray!10, thick, align=center},
        capa_media/.style = {draw, rectangle, rounded corners, minimum width=2.8cm, minimum height=1.2cm, fill=blue!5, thick, align=center},
        capa_alta/.style = {draw, rectangle, rounded corners, minimum width=12cm, minimum height=1.2cm, fill=green!5, thick, align=center},
        capa_api/.style = {draw, rectangle, rounded corners, minimum width=12cm, minimum height=1.2cm, fill=yellow!10, thick, align=center, font=\bfseries},
        flecha/.style = {-{Stealth[scale=1.2]}, thick},
        texto_nivel/.style = {font=\bfseries, text=black!70}
    ]

    \node[capa_base] (base) at (0, 0) {\textbf{Capa Base (Niveles 0 - 2)}\\ \texttt{parameters/} \quad | \quad \texttt{primitives/} \quad | \quad \texttt{crypto/} (SHAKE)};
    \node[texto_nivel, left] at (-6.2, 0) {Fundamentos};

    \node[capa_media] (ntt) at (-4.5, 2.5) {\texttt{ntt/}\\(Aritmética)};
    \node[capa_media] (samp) at (-1.5, 2.5) {\texttt{sampling/}\\(Muestreo)};
    \node[capa_media] (dec) at (1.5, 2.5) {\texttt{decomposition/}\\(Redondeo)};
    \node[capa_media] (enc) at (4.5, 2.5) {\texttt{encoding/}\\(Bitpacking)};
    
    \node[texto_nivel, left] at (-6.2, 2.5) {Niveles 3 - 6};

    \node[capa_alta] (core) at (0, 5) {\textbf{Capa de Protocolo (\texttt{core/}) - Niveles 7 - 9}\\ \texttt{mldsa\_internal.py} (Lógica de Protocolo y Orquestación)};
    \node[texto_nivel, left] at (-6.2, 5) {Protocolo};

    \node[capa_api] (api) at (0, 7.5) {Nivel 10: \texttt{mldsa.py} (API Pública - \texttt{keygen}, \texttt{sign}, \texttt{verify})};
    \node[texto_nivel, left] at (-6.2, 7.5) {Capa de Usuario};

    \draw[flecha] (base.north -| ntt.south) -- (ntt.south);
    \draw[flecha] (base.north -| samp.south) -- (samp.south);
    \draw[flecha] (base.north -| dec.south) -- (dec.south);
    \draw[flecha] (base.north -| enc.south) -- (enc.south);

    \draw[flecha] (ntt.north) -- (core.south -| ntt.north);
    \draw[flecha] (samp.north) -- (core.south -| samp.north);
    \draw[flecha] (dec.north) -- (core.south -| dec.north);
    \draw[flecha] (enc.north) -- (core.south -| enc.north);

    \draw[flecha] (core.north) -- (api.south);

    \draw[thick, loosely dashed, gray] (-8, -1.5) rectangle (6.5, 8.8);
    \node[anchor=south west, text=gray] at (-8, 8.8) {Arquitectura del Repositorio \texttt{mldsa/}};

    \end{tikzpicture}
    }
    \caption{Jerarquía de dependencias del motor ML-DSA.}
    \label{fig:arquitectura_capas}
\end{figure}

Esta arquitectura en capas garantiza que cualquier modificación o parche de seguridad no impacte directamente en la interfaz pública, mitigando el riesgo de introducir vulnerabilidades durante el mantenimiento del código.

\section{Gestión de la Entropía}
\label{sec:entropia}

La seguridad de ML-DSA, al igual que cualquier criptosistema moderno, colapsa si la generación de números aleatorios es predecible. En la fase de generación de claves (\texttt{keygen}) y en la creación de máscaras para la firma, el sistema requiere una semilla de alta entropía. Un error común en implementacioneses el uso de generadores pseudoaleatorios estándar PRNG, los cuales son deterministas y vulnerables.

\subsection{Uso de \texttt{os.urandom} frente a \texttt{random.py}}

Para blindar el motor criptográfico, se ha evitado estrictamente el uso de la librería estándar \texttt{random} de Python. Dicha librería emplea el algoritmo \textit{Mersenne Twister}, el cual posee excelentes propiedades estadísticas para simulaciones, pero no es seguro. La propia especificación del lenguaje advierte explícitamente contra su uso para fines de seguridad \cite{python_docs_random}. Un atacante que observe una secuencia suficiente de bits podría deducir el estado interno del generador y predecir los siguientes valores, comprometiendo la clave privada.

En su lugar, siguiendo las directrices del estándar RFC 4086 \cite{rfc4086}, la implementación delega la obtención de aleatoriedad al sistema operativo mediante la llamada \texttt{os.urandom()}. Esta función extrae entropía de fuentes como interrupciones del kernel o tiempos de latencia, garantizando que las semillas $\rho$, $\rho'$ y $K$ posean la impredecibilidad criptográfica absoluta exigida por el FIPS 204.

\subsection{Determinismo Interno y Expansión}

Es fundamental diferenciar arquitectónicamente entre la \textbf{entropía externa}, la proporcionada por el sistema operativo, y la \textbf{expansión determinista} interna. 

En el motor desarrollado, la aleatoriedad real solo se emplea para generar las semillas iniciales de 32 bytes. A partir de ese punto, el sistema aísla la ejecución y utiliza las funciones XOF descritas en el Capítulo 4, basadas en SHAKE para expandir esa entropía, cumpliendo con las recomendaciones del NIST para Generadores de Bits Aleatorios Deterministas DRBG estandarizados en SP 800-90A \cite{nist_sp800_90a}. Esta separación garantiza que el proceso de firma sea reproducible para la verificación posterior, pero invisible para un atacante sin acceso a la semilla original.


\section{Implementación de la NTT}
\label{sec:implementacion_ntt}

La eficiencia del esquema ML-DSA reside en su capacidad de transformar multiplicaciones polinómicas costosas de complejidad $\mathcal{O}(n^2)$ en operaciones de punto a punto de complejidad $\mathcal{O}(n \log n)$. Esto se logra mediante la Transformada Numérica Teórica NTT. En esta sección se expone la transición desde el estándar FIPS 204 hacia una implementación funcional, siendo demostranda empíricamente

\subsection{Algoritmo NTT al Código Python}

La implementación se ha diseñado para ser una traducción idéntica del estándar. Para comprender este enfoque, primero observamos el pseudocódigo matemático formal dictado por el NIST en el Algoritmo \ref{alg:ntt_fips}.

\begin{algorithm}[H]
\caption{NTT($w$) - Algoritmo 41 del estándar FIPS 204}
\label{alg:ntt_fips}
\textbf{Input:} Polinomio $w(X) = \sum_{j=0}^{255} w_j X^j \in R_q$. \\
\textbf{Output:} $\hat{w} = (\hat{w}[0], \dots, \hat{w}[255]) \in T_q$.
\begin{algorithmic}[1]
\FOR{$j=0$ \TO $255$}
    \STATE $\hat{w}[j] \gets w_j$
\ENDFOR
\STATE $m \gets 0$
\STATE $len \gets 128$
\WHILE{$len \geq 1$}
    \STATE $start \gets 0$
    \WHILE{$start < 256$}
        \STATE $m \gets m + 1$
        \STATE $z \gets \text{ZETAS}[m]$ \COMMENT{$z \gets \zeta^{\text{BitRev8}(m)} \bmod q$}
        \FOR{$j=start$ \TO $start + len - 1$}
            \STATE $t \gets (z \cdot \hat{w}[j + len]) \bmod q$
            \STATE $\hat{w}[j + len] \gets (\hat{w}[j] - t) \bmod q$
            \STATE $\hat{w}[j] \gets (\hat{w}[j] + t) \bmod q$
        \ENDFOR
        \STATE $start \gets start + 2 \cdot len$
    \ENDWHILE
    \STATE $len \gets \lfloor len / 2 \rfloor$
\ENDWHILE
\RETURN $\hat{w}$
\end{algorithmic}
\end{algorithm}

A continuación, el Código \ref{lst:ntt_python} presenta la función desarrollada en Python. Se puede apreciar cómo las variables estructurales (\texttt{m}, \texttt{len}, \texttt{start}) y el triple bucle anidado replican línea a línea el documento oficial.

\begin{minipage}{\linewidth}
\begin{lstlisting}[language=Python, caption={Implementación de la NTT reflejando el Algoritmo 41 de FIPS 204.}, label={lst:ntt_python}]
def ntt(w: list) -> list:
    # 1-3: Copia para preservar la inmutabilidad del polinomio
    w_hat = w.copy()
    m = 0
    length = 128
    
    # 6: Iteracion por etapas (log n etapas)
    while length >= 1:
        start = 0
        # 8: Procesamiento de bloques por etapa
        while start < 256:
            m = m + 1
            z = ZETAS[m] # Raiz de la unidad precomputada
            
            # 11: Bucle interno y operacion Mariposa
            for j in range(start, start + length):
                t = (z * w_hat[j + length]) % Q
                w_hat[j + length] = (w_hat[j] - t) % Q
                w_hat[j] = (w_hat[j] + t) % Q
                
            start = start + 2 * length
        length = length // 2
        
    return w_hat
\end{lstlisting}
\end{minipage}

El núcleo del rendimiento es la \textbf{operación de mariposa} que corresponden a las líneas 16-19 del código en Python. En cada iteración, se toma un par de coeficientes y se combinan utilizando una potencia de la raíz de la unidad $\zeta$. Esta estructura permite que el algoritmo opere sobre el mismo espacio de memoria. De esta manera, la reducción modular explícita módulo $q$ tras cada suma y multiplicación previene el desbordamiento que ocurriría en una convolución polinómica.

\begin{ejemplo}[Multiplicación en un Anillo Reducido]
La verdadera potencia de la NTT se manifiesta al realizar la multiplicación de dos polinomios. En el dominio estándar, multiplicar dos polinomios de grado $n$ requiere una convolución clásica con una complejidad computacional de $\mathcal{O}(n^2)$. 

Al aplicar la NTT, el proceso global reduce esta complejidad a $\mathcal{O}(n \log n)$. Esto se logra transformando ambos polinomios a elementos de $\mathbb{Z}_q^n$ (con un coste de $\mathcal{O}(n \log n)$), donde la pesada convolución se convierte en una simple multiplicación escalar coeficiente a coeficiente de complejidad $\mathcal{O}(n)$, para finalmente devolver el resultado al dominio estándar mediante la transformada inversa.
\vspace{0.3cm}
\textbf{Cálculo de las Raíces de Evaluación}

En la NTT, las raíces de evaluación no se eligen de forma aleatoria; corresponden exactamente a las potencias impares de $\zeta \in \mathbb{Z}_q$ calculadas bajo el módulo $q$. Para este anillo de prueba ($n=4$), buscamos una raíz octava de la unidad. Es decir, un elemento $\zeta = 2$, tal que $2^8 \equiv 1 \pmod{17}$. 

Las cuatro raíces necesarias ($r_i = \zeta^{2i+1} \pmod q$) se derivan de la siguiente manera:
\begin{align*}
    r_0 &= \zeta^1 = 2^1 \equiv 2 \\
    r_1 &= \zeta^3 = 2^3 \equiv 8 \\
    r_2 &= \zeta^5 = 32  \equiv 15 \pmod{17} \\
    r_3 &= \zeta^7 = 128 \equiv 9 \pmod{17}
\end{align*}

\textbf{Paso 1: Transformación al Dominio NTT}

Evaluamos los polinomios de entrada $\mathbf{w}$ y $\mathbf{v}$ sustituyendo la variable $x$ por la secuencia de raíces obtenida $\{2, 8, 15, 9\}$. Para preservar la trazabilidad, se detalla la reducción modular de cada coeficiente:

\vspace{0.3cm}


\textit{Polinomio $\mathbf{w}(x) = 1 + 2x + 3x^2 + 4x^3$:}
\begin{align*}
    \hat{w}_0 &= \mathbf{w}(2)  = 1 + 2(2) + 3(4) + 4(8)   = 49  \equiv \mathbf{15} \pmod{17} \\
    \hat{w}_1 &= \mathbf{w}(8)  = 1 + 2(8) + 3(13) + 4(2)  = 64  \equiv \mathbf{13} \pmod{17} \\
    \hat{w}_2 &= \mathbf{w}(15) = 1 + 2(15) + 3(4) + 4(9)  = 79  \equiv \mathbf{11} \pmod{17} \\
    \hat{w}_3 &= \mathbf{w}(9)  = 1 + 2(9) + 3(13) + 4(15) = 118 \equiv \mathbf{16} \pmod{17}
\end{align*}
El vector transformado es $\hat{\mathbf{w}} = (15, 13, 11, 16)$.
\vspace{0.3cm}
\textit{Polinomio $\mathbf{v}(x) = 5 + 0x + 1x^2 + 1x^3$:}
\begin{align*}
    \hat{v}_0 &= \mathbf{v}(2)  = 5 + 0 + 4 + 8    = 17 \equiv \mathbf{0} \pmod{17} \\
    \hat{v}_1 &= \mathbf{v}(8)  = 5 + 0 + 13 + 2   = 20 \equiv \mathbf{3} \pmod{17} \\
    \hat{v}_2 &= \mathbf{v}(15) = 5 + 0 + 4 + 9    = 18 \equiv \mathbf{1} \pmod{17} \\
    \hat{v}_3 &= \mathbf{v}(9)  = 5 + 0 + 13 + 15  = 33 \equiv \mathbf{16} \pmod{17}
\end{align*}
El vector transformado es $\hat{\mathbf{v}} = (0, 3, 1, 16)$.\vspace{0.3cm}

\textbf{Paso 2: Multiplicación Punto a Punto}

Calculamos el producto de Hadamard $\hat{\mathbf{y}} = \hat{\mathbf{w}} \circ \hat{\mathbf{v}}$ en el dominio de la frecuencia:
\begin{align*}
    \hat{y}_0 &= \hat{w}_0 \cdot \hat{v}_0 = 15 \cdot 0  = 0  \equiv \mathbf{0} \pmod{17} \\
    \hat{y}_1 &= \hat{w}_1 \cdot \hat{v}_1 = 13 \cdot 3  = 39 \equiv \mathbf{5} \pmod{17} \\
    \hat{y}_2 &= \hat{w}_2 \cdot \hat{v}_2 = 11 \cdot 1  = 11 \equiv \mathbf{11} \pmod{17} \\
    \hat{y}_3 &= \hat{w}_3 \cdot \hat{v}_3 = 16 \cdot 16 = 256 \equiv \mathbf{1} \pmod{17}
\end{align*}

El resultado en el dominio NTT es $\hat{\mathbf{y}} = (0, 5, 11, 1)$. Mediante solo 4 multiplicaciones escalares, hemos encapsulado la información completa del producto polinómico.
\end{ejemplo}

\subsection{Algoritmo $\text{NTT}^{-1}$ al Código Python}

La recuperación de los coeficientes polinómicos originales requiere revertir el proceso algebraico. El estándar FIPS 204 \cite{nist_fips204} especifica este procedimiento en su Algoritmo 42, el cual ha sido transcrito a la implementación en Python mostrada en el Código \ref{lst:invNTT_python}.

\begin{algorithm}[H]
\caption{NTT$^{-1}(\hat{w})$ - Algoritmo 42 del estándar FIPS 204}
\label{alg:invNTT}
\textbf{Input:} Vector en dominio NTT $\hat{w} = (\hat{w}[0], \dots, \hat{w}[255]) \in T_q$. \\
\textbf{Output:} Polinomio $w(X) = \sum_{j=0}^{255} w_j X^j \in R_q$.
\begin{algorithmic}[1]
\FOR{$j=0$ \TO $255$}
    \STATE $w_j \gets \hat{w}[j]$
\ENDFOR
\STATE $m \gets 256$
\STATE $len \gets 1$
\WHILE{$len < 256$}
    \STATE $start \gets 0$
    \WHILE{$start < 256$}
        \STATE $m \gets m - 1$
        \STATE $z \gets -ZETAS[m]$ \COMMENT{$z \gets -\zeta^{\text{BitRev8}(m)} \bmod q$}
        \FOR{$j=start$ \TO $start + len - 1$}
            \STATE $t \gets w_j$
            \STATE $w_j \gets (t + w_{j+len}) \bmod q$
            \STATE $w_{j+len} \gets (t - w_{j+len}) \bmod q$
            \STATE $w_{j+len} \gets (z \cdot w_{j+len}) \bmod q$
        \ENDFOR
        \STATE $start \gets start + 2 \cdot len$
    \ENDWHILE
    \STATE $len \gets 2 \cdot len$
\ENDWHILE
\STATE $f \gets 8347681$ \COMMENT{$f = 256^{-1} \bmod q$}
\FOR{$j=0$ \TO $255$}
    \STATE $w_j \gets (f \cdot w_j) \bmod q$
\ENDFOR
\RETURN $w$
\end{algorithmic}
\end{algorithm}

\begin{minipage}{\linewidth}
\begin{lstlisting}[language=Python, caption={Implementación de la NTT$^{-1}$ reflejando el Algoritmo 42.}, label={lst:invNTT_python}]
def ntt_inv(w_hat: list) -> list:
    w = w_hat.copy()
    m = 256
    length = 1
    
    # 6: Bucle de etapas inversas (Gentleman-Sande)
    while length < 256:
        start = 0
        while start < 256:
            m = m - 1
            z = (-ZETAS[m]) % Q # Twiddle factor inverso
            
            # 11: Operacion de mariposa inversa
            for j in range(start, start + length):
                t = w[j]
                w[j] = (t + w[j + length]) % Q
                w[j + length] = (t - w[j + length]) % Q
                w[j + length] = (z * w[j + length]) % Q
                
            start = start + 2 * length
        length = length * 2
        
    # 21: Multiplicacion por el inverso modular de n (n^-1 mod q)
    f = 8347681 
    for j in range(256):
        w[j] = (w[j] * f) % Q
        
    return w
\end{lstlisting}
\end{minipage}

A diferencia de la NTT directa, la $\text{NTT}^{-1}$ utiliza el algoritmo de \textit{Gentleman-Sande}. Se recorre el árbol de dependencias en sentido inverso, es decir, comenzando con \texttt{length = 1} y multiplicando por 2 y, se añade un paso final de normalización: multiplicar todos los coeficientes resultantes por el inverso modular de $n$ bajo módulo $q$.

\begin{ejemplo}[Recuperación del Polinomio mediante $\text{NTT}^{-1}$]
Para finalizar la demostración matemática con el vector en el dominio de la frecuencia $\hat{\mathbf{y}} = (0, 5, 11, 1)$, se aplica la transformación inversa $\text{NTT}^{-1}$. Para evaluar el polinomio inverso, se emplean las raíces de evaluación inversas ($r_j^{-1} \pmod{17}$) y se multiplica la suma resultante por el inverso escalar de la dimensión ($n^{-1} \pmod q$).

Dado $n=4$, su inverso modular es $4^{-1} \equiv 13 \pmod{17}$, puesto que $4 \times 13 = 52 \equiv 1 \pmod{17}$. Las raíces inversas correspondientes al conjunto $\{2, 8, 15, 9\}$ son $\{9, 15, 8, 2\}$.

La ecuación para recuperar cada coeficiente $y_k$ en el dominio estándar se define como:
$$y_k = n^{-1} \sum_{j=0}^{n-1} \hat{y}_j \cdot (r_j^{-1})^k \pmod q$$

El cálculo de los cuatro coeficientes se desglosa a continuación:
\begin{align*}
    y_0 &= 13 \cdot [0(9^0) + 5(15^0) + 11(8^0) + 1(2^0)] \\
        &= 13 \cdot (0 + 5 + 11 + 1) = 221 \equiv \mathbf{0} \pmod{17} \\[1ex]
    y_1 &= 13 \cdot [0(9^1) + 5(15^1) + 11(8^1) + 1(2^1)] \\
        &= 13 \cdot (0 + 75 + 88 + 2) = 2145 \equiv \mathbf{3} \pmod{17} \\[1ex]
    y_2 &= 13 \cdot [0(9^2) + 5(15^2) + 11(8^2) + 1(2^2)] \\
        &= 13 \cdot (0 + 20 + 143 + 4) = 2171 \equiv \mathbf{12} \pmod{17} \\[1ex]
    y_3 &= 13 \cdot [0(9^3) + 5(15^3) + 11(8^3) + 1(2^3)] \\
        &= 13 \cdot (0 + 45 + 22 + 8) = 975 \equiv \mathbf{6} \pmod{17}
\end{align*}

Recomponiendo los coeficientes calculados, se obtiene la ecuación polinómica resultante de la multiplicación en el dominio estándar:
\begin{equation}
    \mathbf{y}(x) = 0 + 3x + 12x^2 + 6x^3 \in \mathbb{Z}_{17}[x] / \langle x^4 + 1 \rangle
\end{equation}

De esta forma, la implementación orientada a código certifica el cumplimiento algorítmico del estándar FIPS 204 \cite{nist_fips204}.
\end{ejemplo}


\section{Análisis del Core}
\label{sec:analisis_core}

El módulo \texttt{core/} contiene los métodos principales del protocolo. Para certificar la implementación se ha mantenido una correspondencia con la especificación del NIST. A continuación, se analizan los puntos críticos del esquema, que son la generación, firma y verificación.

\subsection{Expansión de la Matriz Pública}

En la fase de generación de claves (\texttt{keygen}), el mayor desafío se encuentra en la \textbf{expansión de la matriz pública $\mathbf{\hat{A}}$}. Esta matriz tiene dimensiones $k \times l$ y sus elementos son polinomios en el dominio NTT ($T_q$). Generar y almacenar esta matriz estáticamente en memoria o transmitirla completa por la red resultaría inasumible en entornos reales.

Por ello, el estándar FIPS 204 soluciona este problema de diseño mediante la generación al vuelo. En lugar de la matriz, la clave pública solo almacena una semilla de 32 bytes ($\rho$). El Algoritmo \ref{alg:expand_a} muestra cómo el NIST define la reconstrucción de la matriz a partir de esta semilla.

\begin{algorithm}[H]
\caption{\texttt{ExpandA}($\rho$) - Algoritmo 32 del estándar FIPS 204}
\label{alg:expand_a}
\textbf{Input:} Una semilla $\rho \in \mathbb{B}^{32}$. \\
\textbf{Output:} Matriz $\mathbf{\hat{A}} \in (T_q)^{k \times l}$.
\begin{algorithmic}[1]
\FOR{$r=0$ \TO $k-1$}
    \FOR{$s=0$ \TO $l-1$}
        \STATE $\rho' \gets \rho \parallel \text{IntegerToBytes}(s, 1) \parallel \text{IntegerToBytes}(r, 1)$
        \STATE $\mathbf{\hat{A}}[r, s] \gets \text{RejNTTPoly}(\rho')$ \COMMENT{La semilla $\rho'$ depende de $s$ y $r$}
    \ENDFOR
\ENDFOR
\RETURN $\mathbf{\hat{A}}$
\end{algorithmic}
\end{algorithm}

La implementación desarrollada en Python (Código \ref{lst:expand_a_python}) refleja una fiel representación al pseudocódigo anterior.
\vspace{0.3cm}

\begin{minipage}{\linewidth}
\begin{lstlisting}[language=Python, caption={Implementación de \texttt{ExpandA} asegurando la separación de dominio.}, label={lst:expand_a_python}]
def expand_a(rho: bytes, k: int, l: int) -> list:
    # Inicializamos matriz vacia de k filas por l columnas
    A_hat = [[None for _ in range(l)] for _ in range(k)]
    
    for r in range(k):
        for s in range(l):
            # Garantiza separacion de dominio para cada celda
            rho_prime = (rho + integer_to_bytes(s, 1) + 
                         integer_to_bytes(r, 1))
            
            # Generamos polinomio en T_q usando muestreo de rechazo
            A_hat[r][s] = rej_ntt_poly(rho_prime)
            
    return A_hat
\end{lstlisting}
\end{minipage}

Esta implementación aplica la \textbf{separación de dominio}. En la línea 9, la semilla $\rho$ se concatena con las coordenadas de la celda ($s$ y $r$) codificadas como \textit{bytes}. Esto garantiza que la función XOF interna inicia un estado único para cada iteración. De esta forma, se evita la colisión de polinomios entre diferentes celdas de la matriz, cumpliendo estrictamente con la distribución uniforme.


\subsection{Muestreo por Rechazo y Fiat-Shamir}
\label{subsec:sign_reject}

El proceso de firma implementado corresponde al paradigma \textit{Fiat-Shamir con Abortos}\cite{lyubashevsky2009fiat}. A diferencia de los esquemas clásicos como RSA o ECDSA, la generación de la firma en ML-DSA no es un cálculo algebraico directo, sino un proceso iterativo. 

El mayor riesgo de seguridad en los esquemas basados en retículos es la filtración de la clave privada $\mathbf{s}_1, \mathbf{s}_2$ a través de la firma pública $\mathbf{z}$. Para evitar esto, el algoritmo genera una máscara aleatoria $\mathbf{y}$, y el resultado se somete a estrictas comprobaciones. Si los coeficientes exceden un límite seguro de magnitud, el algoritmo destruye la firma parcial, incrementa un contador $\kappa$ y vuelve a empezar.

El Algoritmo \ref{alg:sign_reject} describe parcialmente el diseño lógico de este bucle de rechazo del FIPS 204.

\begin{algorithm}[H]
\caption{Fragmento Algoritmo 7: Bucle de Muestreo por Rechazo}
\label{alg:sign_reject}
\begin{algorithmic}[1]
\WHILE{$(\mathbf{z}, \mathbf{h}) = \bot$}
    \STATE $\mathbf{y} \gets \text{ExpandMask}(\rho'', \kappa)$
    \STATE \dots \COMMENT{Cálculo del compromiso $\mathbf{w}_1$, desafío $c$ y productos cruzados}
    \STATE $\mathbf{z} \gets \mathbf{y} + c \mathbf{s}_1$
    \STATE $\mathbf{r}_0 \gets \text{LowBits}(\mathbf{w} - c \mathbf{s}_2)$
    
    \IF{$\|\mathbf{z}\|_\infty \ge \gamma_1 - \beta$ \OR $\|\mathbf{r}_0\|_\infty \ge \gamma_2 - \beta$}
        \STATE $(\mathbf{z}, \mathbf{h}) \gets \bot$ \COMMENT{Abortar: Riesgo de filtración de clave}
    \ELSE
        \STATE \dots \COMMENT{Cálculo del vector de ayuda $\mathbf{h}$ y $c\mathbf{t}_0$}
        \IF{$\|c\mathbf{t}_0\|_\infty \ge \gamma_2$ \OR $\text{peso}(\mathbf{h}) > \omega$}
            \STATE $(\mathbf{z}, \mathbf{h}) \gets \bot$ \COMMENT{Abortar: Fallo de compresión de la firma}
        \ENDIF
    \ENDIF
    \STATE $\kappa \gets \kappa + l$ \COMMENT{Avanzar estado determinista}
\ENDWHILE
\end{algorithmic}
\end{algorithm}

La implementación en Python se expone en el Código \ref{lst:sign_python}. Destaca el uso de la Norma Infinita ($\|\cdot\|_\infty$) y el conteo del peso de Hamming como barreras de seguridad antes de emitir los \textit{bytes} finales.

\begin{minipage}{\linewidth}
\begin{lstlisting}[language=Python, caption={Bucle iterativo estocástico de la firma (Fiat-Shamir con abortos).}, label={lst:sign_python}]
# Comienzo del bucle de muestreo por rechazo
kappa = 0
while True:
    # Generacion determinista de la mascara candidata (y)
    y = expand_mask(rho_pp, kappa, l, gamma1)

    # Calculo del compromiso w y el desafio c 

    # La "respuesta" criptografica
    z = [[(y[i][j] + cs1[i][j]) % Q for j in range(256)] for i in range(l)]

    # r0 <- LowBits(w - <<cs2>>)
    w_minus_cs2 = [[(w[i][j] - cs2[i][j]) % Q for j in range(256)] 
                   for i in range(k)]
    r0 = lowbits(w_minus_cs2, gamma2)

    # Norma Infinita
    # Evita que 'z' sea demasiado grande y filtre geometria de la clave 's1'
    if _inf_norm(z) >= (gamma1 - beta) or _inf_norm(r0) >= (gamma2 - beta):
        kappa += l
        continue 
    
    # Peso de Hamming y Norma
    # Garantiza que el vector de ayuda 'h' cumpla con el limite de peso
    peso_h = sum(sum(1 for x in p if x != 0) for p in h)
    if _inf_norm(ct0) >= gamma2 or peso_h > omega:
        kappa += l
        continue
        
    # Firma matematicamente valida, proceder al empaquetado final
    z_cent = [[(x % Q if x % Q <= Q//2 else (x % Q) - Q) for x in p] 
              for p in z]
    return sig_encode(c_tilde, z_cent, h, gamma1, omega, k)
\end{lstlisting}
\end{minipage}

 Al abortar (\texttt{continue}), no se solicita al sistema operativo nueva entropía. En su lugar, se incrementa el contador \texttt{k} en un valor $l$ y se inyecta nuevamente en la función \texttt{expand\_mask}. Esto garantiza que, si se firma el mismo documento con la misma semilla $\rho''$, los intentos fallidos serán exactamente los mismos en cualquier ordenador del mundo y produciendo una firma idéntica.



 \subsection{Reconstrucción del Compromiso}
\label{subsec:verify_reconstruction}

En la criptografía asimétrica tradicional, como RSA, la verificación se fundamenta en "desencriptar" la firma utilizando la clave pública para recuperar el \textit{hash} original del mensaje. Sin embargo, en el estándar ML-DSA, cambia por completo: la verificación consiste en un proceso de \textbf{reconstrucción}.

El verificador utiliza la respuesta del firmante $\mathbf{z}$, el vector de ayuda $\mathbf{h}$ y la parte $\mathbf{t}_1$ de la clave pública, para recalcular una aproximación del compromiso original $\mathbf{w}'$. Posteriormente, si al reempaquetar los bits altos de esta aproximación el \textit{hash} resultante coincide con el de la firma, se garantiza la integridad matemática.

El Algoritmo \ref{alg:verify_core} ilustra el núcleo de este mecanismo de reconstrucción.

\begin{algorithm}[H]
\caption{Fragmento Algoritmo 8: Reconstrucción y Verificación}
\label{alg:verify_core}
\begin{algorithmic}[9]
\STATE $\mathbf{w}' \gets \text{NTT}^{-1}(\mathbf{\hat{A}} \circ \text{NTT}(\mathbf{z}) - \text{NTT}(c) \circ \text{NTT}(\mathbf{t}_1 \cdot 2^d))$
\STATE $\mathbf{w}_1' \gets \text{UseHint}(\mathbf{h}, \mathbf{w}')$ \COMMENT{Reconstrucción del compromiso del firmante}
\STATE $c' \gets \text{H}(\mu \parallel \text{w1Encode}(\mathbf{w}_1'), \lambda/4)$ \COMMENT{Cálculo del hash de verificación}
\RETURN $(\|\mathbf{z}\|_\infty < \gamma_1 - \beta) \text{ \AND } (c == c')$
\end{algorithmic}
\end{algorithm}

La implementación de este bloque (Código \ref{lst:verify_python}) demuestra la eficiencia de la NTT durante el proceso de verificación. Calcular la expresión $\mathbf{A}\mathbf{z} - c\mathbf{t}_1 \cdot 2^d$ exigiría pesadas convoluciones de complejidad $\mathcal{O}(n^2)$. 

Para optimizar este proceso, el motor criptográfico transforma previamente todos los operandos al dominio establecido por el FIPS 204. De este modo, $\mathbf{\hat{A}} \circ \mathbf{\hat{z}} - \mathbf{\hat{c}} \circ \mathbf{\hat{t}}_1 \cdot 2^d$, donde la matriz matemática se reduce a multiplicaciones escalares coeficiente a coeficiente de complejidad $\mathcal{O}(n)$, aplicando únicamente la transformada inversa al resultado final.\vspace{0.3cm}

\begin{minipage}{\linewidth}
\begin{lstlisting}[language=Python, caption={Cálculo de la aproximación y verificación de la firma.}, label={lst:verify_python}]
# 9: w' <- Az - c*t1*2^d (Aproximacion del compromiso original)
# Calculo acelerado en el dominio NTT y posterior transformada inversa
Az_hat = matrix_vector_ntt(A_hat, z_hat)
ct1_hat = scalar_vector_ntt(c_hat, t1_2d_hat)

w_approx = []
for i in range(k):
    # Resta coeficiente a coeficiente aplicando el modulo Q
    diff_hat = [(Az_hat[i][j] - ct1_hat[i][j]) % Q for j in range(256)]
    w_approx.append(ntt_inv(diff_hat))

# 10: w1' <- UseHint(h, w_approx)
# Recuperacion de los bits altos exactos corrigiendo errores con 'h'
w1_prime = use_hint(h, w_approx, gamma2)

# 12: c'_tilde <- H(mu || w1Encode(w1'))
c_tilde_prime = H.digest(mu + w1_encode(w1_prime, gamma2), c_tilde_len)

# 13: Verificacion bidimensional: 
# 1. Que 'z' sea un vector corto (Norma Infinita)
# 2. Que la reconstruccion coincida con el compromiso original
z_is_short = _inf_norm(z) < (gamma1 - beta)
hashes_match = (c_tilde == c_tilde_prime)

return z_is_short and hashes_match
\end{lstlisting}
\end{minipage}

Es importante destacar la validación de la línea 22 (\texttt{z\_is\_short}). Aunque los \textit{hashes} coincidan perfectamente (\texttt{hashes\_match == True}), el verificador debe rechazar la firma si la norma infinita de $\mathbf{z}$ supera el umbral $\gamma_1 - \beta$. Esta comprobación mitiga los ataques de "falsificación de firmas" donde un atacante maligno inyecta coeficientes grandes para forzar colisiones en el algoritmo \texttt{UseHint}.


\section{Pruebas Unitarias}
\label{sec:validacion_unitaria}

La implementación de un estándar criptográfico de la magnitud del FIPS 204 exige un rigor absoluto. Un fallo en la gestión de un solo bit durante los procesos de compresión o un desbordamiento en el módulo $q$ comprometería la seguridad de todo el esquema. Por ello, la validación del motor se ha enfocado mediante una metodología de Desarrollo Guiado por Pruebas TDD.

\subsection{Estructura de la Suite de Pruebas TDD}

Se ha diseñado un conjunto de validación modular estructurada en espejo a la arquitectura del código fuente. El objetivo fundamental ha sido certificar individualmente la correctitud matemática de los más de 50 algoritmos auxiliares definidos en el estándar antes de su integración en las funciones core de firma y verificación.

El banco de pruebas desarrollado se compone de los siguientes módulos críticos:
\begin{itemize}
    \item \texttt{test\_ntt}: Validación de las transformadas directas e inversas y comprobación de las propiedades algebraicas en el anillo $R_q$.
    \item \texttt{test\_encoding}: Certificación de los algoritmos de empaquetado y desempaquetado de \textit{bytes} (\texttt{pkEncode}, \texttt{sigDecode}, \texttt{w1Encode}), verificando el cumplimiento del formato \textit{little-endian}.
    \item \texttt{test\_sampling}: Pruebas estadísticas sobre las funciones de generación aleatoria, como el muestreo de rechazo y \texttt{SampleInBall}.
    \item \texttt{test\_decomposition}: Verificación de las funciones de extracción de bits (\texttt{HighBits}, \texttt{LowBits}, \texttt{MakeHint}).
    \item \texttt{test\_core}: Validación de la orquestación final de los métodos \texttt{KeyGen}, \texttt{Sign} y \texttt{Verify}.
\end{itemize}

\subsection{Ejecución y Resultados}

Para ilustrar la validación, el Código \ref{lst:test_ntt} muestra un fragmento del test \texttt{test\_ntt.py}. En este test no solo se comprueba el funcionamiento básico, sino que se estresa el algoritmo con datos aleatorios.

\begin{minipage}{\linewidth}
\begin{lstlisting}[language=Python, caption={Pruebas unitarias del NTT.}, label={lst:test_ntt}]
class TestNTTTransformations(unittest.TestCase):

    def test_ntt_roundtrip_zeros(self):
        """Caso base: El polinomio cero debe seguir siendo cero en el dominio NTT."""
        w = [0] * 256
        w_hat = ntt(w)
        self.assertEqual(w_hat, [0] * 256)
        
        w_recuperado = ntt_inv(w_hat)
        self.assertEqual(w, w_recuperado)

    def test_ntt_roundtrip_random_stress(self):
        """Prueba de estres: 100 polinomios aleatorios de ida y vuelta."""
        for _ in range(100):
            w_original = [random.randint(0, Q - 1) for _ in range(256)]
            
            # NTT (ida) y NTT Inversa (vuelta)
            w_hat = ntt(w_original)
            w_recuperado = ntt_inv(w_hat)
            
            self.assertEqual(w_original, w_recuperado, 
                             "Error de desbordamiento en la NTT")
            
            # Verificacion de pertenencia estricta al anillo [0, Q-1]
            self.assertTrue(all(0 <= c < Q for c in w_hat))

    def test_ntt_linearity(self):
        """Propiedad homomorfica: NTT(a + b) == NTT(a) + NTT(b)."""
        a = [random.randint(0, Q - 1) for _ in range(256)]
        b = [random.randint(0, Q - 1) for _ in range(256)]
        
        # 1. Suma en dominio normal y posterior NTT
        a_plus_b = [(a[i] + b[i]) % Q for i in range(256)]
        ntt_suma = ntt(a_plus_b)
        
        # 2. NTT individual y posterior suma en dominio de frecuencia
        ntt_a = ntt(a)
        ntt_b = ntt(b)
        suma_ntts = add_ntt(ntt_a, ntt_b)
        
        self.assertEqual(ntt_suma, suma_ntts, 
                         "La NTT debe ser una operacion estrictamente lineal")
\end{lstlisting}
\end{minipage}

La superación de estos tests garantiza tres pilares fundamentales para el proyecto:
\begin{enumerate}
    \item \textbf{Reversibilidad perfecta:} Se confirma que la aplicación de la transformada inversa recupera los datos originales sin degradación.
    \item \textbf{Integridad de los límites:} Se certifica que todas las operaciones aritméticas respetan el módulo $q$ sin producir errores de desbordamiento.
    \item \textbf{Consistencia algorítmica:} Al validar cada uno de los 50 sub-algoritmos, se asegura que la integración final del estándar FIPS 204 es robusta y fiel a la especificación oficial.
\end{enumerate}

Con la validación técnica del motor concluida, se dispone de una base criptográfica sólida para su despliegue en el entorno distribuido que se detalla en el siguiente capítulo.



\chapter{Diseño y Desarrollo del Sistema Distribuido}
\label{ch:desarrollo_sistema}

Este capítulo detalla el proceso de diseño, arquitectura e implementación de la infraestructura que da soporte al motor criptográfico ML-DSA. El objetivo es transformar un algoritmo matemático complejo en una plataforma funcional que permita la gestión de identidades y la firma de documentos PDF en un entorno web real.

\section{Visión General y Estrategia de Diseño}
\label{sec:vision_general_arquitectura}

\subsection{Naturaleza y Finalidad de la Herramienta}
\label{subsec:naturaleza_finalidad}

La transición hacia la criptografía post-cuántica PQC no es únicamente un desafío matemático, sino un profundo reto de ingeniería de \textit{software}. El estándar FIPS 204 define el funcionamiento de ML-DSA \cite{nist_fips204}; sin embargo, su integración en ecosistemas web modernos puede presentar problemas como: mayor tamaño de claves, latencia computacional y la necesidad de mantener la integridad en formatos de archivo como el PDF.

El sistema desarrollado es una \textbf{arquitectura de referencia y prueba de concepto PoC} orientada a resolver la integración práctica de ML-DSA en un entorno distribuido. 

La finalidad principal de la plataforma es demostrar la viabilidad técnica de crear criptografía resistente a la computación cuántica. La herramienta aborda tres objetivos de ingeniería fundamentales:

\begin{enumerate}
    \item \textbf{Abstracción Criptográfica:} Transforma un motor basado en retículos en un servicio accesible mediante una interfaz RESTful, aislando la complejidad algebraica del usuario final.
    \item \textbf{Interoperabilidad Documental:} Implementa un método para inyectar firmas post-cuánticas, cuya carga útil es superior a las firmas convencionales, en documentos PDF bajo el estándar ISO 32000-1 \cite{iso_pdf} sin invalidar su estructura.
    \item \textbf{Gestión de Identidad y Persistencia:} Establece un flujo seguro donde las claves y registros se vinculan a sesiones web persistentes, permitiendo operaciones remotas sin exponer el material criptográfico en canales no seguros.
\end{enumerate}

Por lo tanto, este desarrollo se presenta como un modelo que representa una implementación real para iniciar la transición hacia la \textit{Agilidad Criptográfica}, anticipándose a la amenaza que supone la futura computación cuántica a gran escala.


\subsection{Ecosistema Tecnológico}
\label{subsec:ecosistema_tecnologico}

Para crear la arquitectura descrita, se ha seleccionado un conjunto de herramientas tecnológicas fundamentado en la seguridad, el procesamiento asíncrono y la defensa en profundidad. La elección de estas herramientas responde directamente a los requisitos por este tipo de criptografía tan avanzada.

\begin{itemize}
    \item \textbf{Capa de Presentación (React y TypeScript):} El cliente web se ha desarrollado utilizando la librería React \cite{react_meta}. La decisión crítica en esta capa ha sido la adopción estricta de TypeScript. En una aplicación que manipula flujos de \textit{bytes} en crudo y estructuras binarias, el tipado estático previene vulnerabilidades de corrupción de datos en tiempo de ejecución. Adicionalmente, el ecosistema de React facilita una futura migración hacia arquitecturas de cifrado local mediante la \textit{Web Crypto API}.
    
    \item \textbf{Capa de Lógica y Cómputo (FastAPI):} Dado que la manipulación algebraica y las implementaciones de referencia de ML-DSA poseen un fuerte soporte en ecosistemas Python y C, el servidor RESTful se ha construido sobre el \textit{framework} FastAPI \cite{fastapi_ramirez}. La principal ventaja arquitectónica de FastAPI es su soporte nativo para el estándar ASGI \cite{asgi_spec}. Dado que la generación de firmas post-cuánticas requiere un uso intensivo de la CPU, un servidor síncrono tradicional bloquearía el hilo de ejecución. El modelo asíncrono permite a la API mantener una alta disponibilidad y gestionar peticiones concurrentes sin perder experiencia del usuario.
    
    \item \textbf{Capa de Persistencia e Identidad (Supabase / PostgreSQL):} Implementar un gestor de identidad desde cero en una aplicación criptográfica supone un riesgo crítico de seguridad. Para mitigar esta superficie de ataque, la persistencia se delega en Supabase \cite{supabase_docs}. Este servicio proporciona un núcleo PostgreSQL robusto acoplado a un motor de autenticación centralizado. Su mayor aportación al sistema es la implementación de políticas RLS (\textit{Row Level Security}), que garantizan a nivel de motor de base de datos que un usuario únicamente pueda consultar y manipular sus propios registros y firmas, estableciendo un aislamiento total de la información.
\end{itemize}




\subsection{JSON Web Tokens JWT}
\label{subsec:jwt_identidad}

La arquitectura distribuida propuesta, donde el cliente \textit{React}, el servidor de lógica \textit{FastAPI} y el servicio de persistencia \textit{Supabase} operan en nodos independientes, requiere un mecanismo de autenticación que sea ligero y seguro. Para resolver este desafío, se ha implementado el estándar \textbf{JSON Web Token (RFC 7519)} \cite{ietf_jwt}.

\subsubsection{Concepto y Función}

En este ecosistema, el JWT actúa como un pasaporte digital. Una vez que el usuario se autentica en Supabase, recibe un token firmado que acredita su identidad. A partir de ese momento, el cliente no necesita enviar sus credenciales en cada petición; simplemente presenta este token ante la API. La API, al poseer la clave secreta compartida, puede validar la autenticidad del token de forma local y asíncrona.

\subsubsection{Anatomía Técnica de un JWT}

Un token se compone de tres bloques codificados en \textit{Base64URL} y separados por puntos. Esta estructura permite que el token sea transportado de forma eficiente en cabeceras HTTP:

\begin{enumerate}
    \item \textbf{Header (Cabecera):} Contiene los metadatos de control, especificando el tipo de token (\texttt{JWT}) y el algoritmo de firma utilizado, que en este caso es RS256, como especifica el estándar RFC 7518 \cite{rfc7518}.
    \item \textbf{Payload (Cuerpo/Carga útil):} Es el núcleo de la información. Contiene los denominados \textit{claims}. En este proyecto, se utilizan \textit{claims} registrados como \texttt{sub} (ID único del usuario), \texttt{email} y \texttt{exp} (fecha de expiración), permitiendo que la API identifique al firmante sin realizar consultas adicionales.
    \item \textbf{Signature (Firma):} Es el componente que garantiza la integridad. Se genera calculando un \textit{hash} de la cabecera y el cuerpo utilizando una clave secreta. Cualquier modificación mínima en los datos del token invalidaría la firma, permitiendo que FastAPI rechace la petición de forma inmediata.
\end{enumerate}

\subsubsection{El Flujo en la Arquitectura JWT}

El desafío principal resuelto mediante JWT es la \textbf{validación descentralizada}. El flujo de trabajo se detalla a continuación (ver Figura \ref{fig:flujo_jwt}):

\begin{figure}[H]
    \centering
    \begin{tikzpicture}[
        >=Stealth,
        node distance=3cm and 3cm,
        caja/.style={rectangle, draw=black, thick, fill=white, minimum width=3.5cm, minimum height=1.5cm, align=center, font=\sffamily},
        proceso/.style={rectangle, draw=black, thick, fill=gray!10, rounded corners, minimum width=4cm, minimum height=1.2cm, align=center, text width=3.8cm, font=\footnotesize\sffamily},
        alerta/.style={rectangle, draw=red!80, fill=red!5, dashed, text width=3.5cm, align=center, font=\scriptsize\sffamily, rounded corners}
    ]

    \node[caja] (cliente) {\textbf{Cliente Web} \\ (React)};
    \node[caja, right=of cliente] (supabase) {\textbf{Emisor de Identidad} \\ (Supabase Auth)};
    \node[caja, below=3cm of cliente] (api) {\textbf{Backend} \\ (FastAPI)};

    \draw[->, thick] ([yshift=3mm]cliente.east) -- node[above, font=\scriptsize] {1. Petición Login} ([yshift=3mm]supabase.west);
    \draw[<-, thick] ([yshift=-3mm]cliente.east) -- node[below, font=\scriptsize] {2. Emite JWT Firmado} ([yshift=-3mm]supabase.west);

    \draw[->, thick] ([xshift=3mm]cliente.south) -- node[right, font=\scriptsize, align=left, text width=3cm] {3. Petición HTTP \\ \texttt{Auth: Bearer \textless token\textgreater}} ([xshift=3mm]api.north);

    \node[proceso, right=of api] (validacion) {\textbf{4. Validación} \\ Decodifica y verifica firma con \texttt{JWT\_SECRET}};
    \draw[->, thick] (api.east) -- (validacion.west);

    \draw[<-, thick] ([xshift=-3mm]cliente.south) -- node[left, font=\scriptsize, align=right, text width=2.5cm] {5. Respuesta \\ (ej. PDF Firmado)} ([xshift=-3mm]api.north);

    \node[alerta, above=0.6cm of validacion] (nota) {¡FastAPI no consulta a Supabase en cada petición!};
    \draw[->, dashed, draw=red!80, thick] (nota.south) -- (validacion.north);

    \end{tikzpicture}
    \caption{Secuencia del flujo de autenticación \textit{stateless} mediante JWT.}
    \label{fig:flujo_jwt}
\end{figure}

\begin{itemize}
    \item \textbf{Emisión:} Supabase Auth actúa como el emisor de confianza (\textit{Issuer}), generando el token tras validar el \textit{login} del usuario.
    \item \textbf{Tránsito:} El cliente almacena el token y lo inyecta en la cabecera \texttt{Authorization: Bearer <token>} de cada petición hacia la API.
    \item \textbf{Verificación Local:} FastAPI recibe la petición y, mediante un \textit{middleware} de seguridad, decodifica el token y verifica la firma contra el \texttt{JWT\_SECRET} de Supabase.
\end{itemize}

Esta metodología elimina el cuello de botella que supondría obligar a la API a consultar la base de datos para verificar la sesión en cada operación criptográfica, lo cual es importante dada la carga computacional que ya impone el motor ML-DSA.

\subsubsection{Ventajas del Modelo}

La adopción de JWT frente a las sesiones tradicionales basadas en \textit{cookies} aporta tres beneficios críticos al sistema:
\begin{itemize}
    \item \textbf{Escalabilidad Horizontal:} Al no existir un estado de sesión en la memoria del servidor, se podrían desplegar múltiples instancias de la API sin necesidad de sincronizar bases de datos de sesiones.
    \item \textbf{Seguridad Cross-Domain:} Facilita la comunicación entre el dominio del frontend y el dominio del backend.
    \item \textbf{Agilidad Criptográfica:} El estándar JWT permite actualizar los algoritmos de firma de forma transparente, facilitando la transición hacia firmas de identidad post-cuánticas en futuras iteraciones del protocolo.
\end{itemize}


\subsection{Infraestructura de Datos y Persistencia}
\label{subsec:infraestructura_datos}

Diseñar la capa de persistencia en un criptosistema distribuido requiere un análisis previo: primero se debe definir qué información es vital para el sistema, qué reglas la gobiernan y, finalmente, cómo se estructura físicamente. Para esta prueba, el modelo relacional en PostgreSQL (Supabase) se ha construido siguiendo este proceso.

\subsubsection{Fase 1: Identificación de Requisitos de Información}

El paso inicial en el diseño de la base de datos es definir las dimensiones de datos necesarias para satisfacer los casos de uso del sistema. Se han identificado cuatro pilares fundamentales, todas ellas sobre la existencia de un usuario autenticado:

\begin{enumerate}
    \item \textbf{Dimensión de Identidad:} El sistema necesita almacenar información biográfica y profesional básica del usuario (nombre, especialidad, organización) para dar contexto a las autorías de las firmas en la interfaz web.
    \item \textbf{Dimensión Criptográfica:} Es el núcleo del sistema. Se requiere un almacenamiento persistente que asocie a un usuario con su par de claves ML-DSA (pública y privada), permitiendo su recuperación en futuras sesiones.
    \item \textbf{Dimensión de Telemetría y Cuotas:} Para evitar un uso excesivo del servidor, el sistema debe monitorizar cuántas operaciones le restan a cada usuario según su nivel de autorización.
    \item \textbf{Dimensión de Auditoría:} El sistema necesita un historial que registre el "quién, qué y cuándo" de cada operación de generación de claves o firma de documentos.
\end{enumerate}

\subsubsection{Fase 2: Reglas y Restricciones del Dominio}

Una vez establecidos los requisitos, se han aplicado reglas de diseño para garantizar la solidez del sistema frente a casos límite y vulnerabilidades:

\begin{itemize}
    \item \textbf{Lenguaje y Normalización:} Para evitar la ambigüedad, se ha forzado un lenguaje único. La entidad central es un identificador criptográfico \texttt{user\_id} gestionado centralizadamente por el proveedor de identidad \textit{Supabase Auth}.
    \item \textbf{Inmutabilidad del Registro:} Dado el requisito de auditoría, la tabla que registre los eventos se diseña bajo la regla de "solo inserción". No existirán borrados físicos ni lógicos sobre el historial. Si un usuario elimina su cuenta, se retendrá los \textit{logs} de forma anónima, pero aplicará un borrado en cascada sobre su bóveda de claves para inutilizar la identidad criptográfica.
    \item \textbf{Aislamiento de Secretos:} Los datos biométricos o de perfil no pueden convivir en la misma tabla que el material criptográfico. Esta separación permite aplicar políticas de Seguridad a Nivel de Fila RLS extremadamente restrictivas sobre las claves, sin bloquear el acceso al perfil público.
\end{itemize}

\subsubsection{Fase 3: Entidades, Cardinalidad y Modelo Físico}

A la hora de crear el diagrama, se han definido cuatro entidades fuertemente tipadas. La cardinalidad se ha diseñado evitando la complejidad del Muchos a Muchos (N:M):

\begin{itemize}
    \item \textbf{Relaciones 1:1:} Las entidades \texttt{profiles}, \texttt{crypto\_identities} y \texttt{user\_usage} mantienen una relación uno a uno con el servicio de autenticación.
    \item \textbf{Relaciones 1:N:} La entidad \texttt{activity\_logs} representa la única relación uno a muchos.
\end{itemize}

El siguiente diagrama ilustra la arquitectura final de datos, dando una importancia clave a la autenticación que ofrece Supabase, así como los atributos y claves (PK/FK) que actúan como parte fundamental del sistema.

\begin{figure}[H]
    \centering
    \begin{tikzpicture}[
        node distance=1.5cm and 2.5cm,
        every node/.style={font=\sffamily\scriptsize},
        table/.style={rectangle, draw=black, thick, fill=white, inner sep=5pt},
        authnode/.style={rectangle, rounded corners, draw=black, thick, fill=gray!15, align=center, inner sep=8pt, font=\bfseries\small}
    ]

    \node[authnode] (auth) {Supabase \\ Auth Service};

    \node[table, above left=1cm and 2cm of auth] (profiles) {
        \begin{tabular}{@{}l@{}}
            \textbf{profiles} \\
            \hline
            \underline{id} (UUID, PK, FK) \\
            first\_name (Text) \\
            last\_name (Text) \\
            full\_name (Text) \\
            specialty (Text) \\
            organization (Text) \\
            avatar\_url (Text) \\
            created\_at (Timestamptz)
        \end{tabular}
    };

    \node[table, below left=1cm and 2cm of auth] (logs) {
        \begin{tabular}{@{}l@{}}
            \textbf{activity\_logs} \\
            \hline
            \underline{id} (UUID, PK) \\
            user\_id (UUID, FK) \\
            action\_type (Text) \\
            file\_name (Text) \\
            created\_at (Timestamptz)
        \end{tabular}
    };

    \node[table, above right=1cm and 2cm of auth] (crypto) {
        \begin{tabular}{@{}l@{}}
            \textbf{crypto\_identities} \\
            \hline
            \underline{user\_id} (UUID, PK, FK) \\
            public\_key (Text) \\
            private\_key (Text) \\
            email (Text) \\
            created\_at (Timestamptz)
        \end{tabular}
    };

    \node[table, below right=1cm and 2cm of auth] (usage) {
        \begin{tabular}{@{}l@{}}
            \textbf{user\_usage} \\
            \hline
            \underline{user\_id} (UUID, PK, FK) \\
            plan\_rank (Text) \\
            ops\_remaining (Int4) \\
            updated\_at (Timestamptz)
        \end{tabular}
    };

    \draw[thick] ([yshift=3mm]auth.west) -| node[pos=0.1, above] {1} node[pos=0.9, left] {1} (profiles.south);
    \draw[thick] ([yshift=-3mm]auth.west) -| node[pos=0.1, below] {1} node[pos=0.9, left] {N} (logs.north);
    
    \draw[thick] ([yshift=3mm]auth.east) -| node[pos=0.1, above] {1} node[pos=0.9, right] {1} (crypto.south);
    \draw[thick] ([yshift=-3mm]auth.east) -| node[pos=0.1, below] {1} node[pos=0.9, right] {1} (usage.north);

    \end{tikzpicture}
    \caption{Diagrama de la base de datos.}
    \label{fig:diagrama_relacional}
\end{figure}


\section{Fase I: Prototipo Funcional y API REST}
\label{sec:fase1_api}

El núcleo del sistema reside en su capacidad para ejecutar los algoritmos definidos en el estándar FIPS 204 y exponerlos como servicios web. Esta fase detalla la construcción de la API RESTful utilizando FastAPI, actuando como puente de comunicación estricto entre el motor criptográfico en el servidor y el cliente final.

\subsection{Configuración del Servidor Asíncrono}
\label{subsec:configuracion_servidor}

El despliegue de un motor criptográfico basado en retículos en un entorno web presenta un desafío crítico: las operaciones de generación de claves y firma, que involucran Transformadas NTT y múltiples ciclos de muestreo de rechazo, tareas altamente intensivas en CPU. Si estas operaciones se ejecutaran de forma síncrona en el hilo principal del servidor, bloquearían cualquier otra petición entrante, provocando una denegación de servicio (DoS) técnica en escenarios de concurrencia.

Para mitigar este cuello de botella y garantizar la escalabilidad, la infraestructura base del servidor se ha diseñado aprovechando el estándar \textit{Asynchronous Server Gateway Interface} ASGI de FastAPI, estructurado en tres pilares de configuración:

\subsubsection*{1. Gestión de la Concurrencia (Event Loop y ThreadPool)}

En el desarrollo de servidores web, es necesario distinguir entre dos tipos de carga de trabajo. Por un lado, las operaciones de red, como recibir una petición HTTP o consultar una base de datos, conocidas como \textit{I/O-bound}. FastAPI maneja estas tareas de forma extremadamente eficiente mediante su \textit{Event Loop} basado en la librería \texttt{asyncio} de Python \cite{python_asyncio}: un hilo principal que atiende peticiones de forma asíncrona sin quedarse bloqueado.

Sin embargo, generar firmas con ML-DSA es una tarea puramente matemática que satura el procesador. Si permitiéramos que el hilo principal realizara estos cálculos, el servidor entero se congelaría hasta terminar la firma, dejando a cualquier otro usuario en un estado de espera excesivamente largo.

Para solucionar este problema de arquitectura, la configuración del servidor delega las llamadas al motor ML-DSA a un \textit{ThreadPoolExecutor} \cite{python_concurrent}. Este mecanismo nativo envía las operaciones matemáticas pesadas a un grupo de subprocesos secundarios. Mientras los núcleos del procesador resuelven la aritmética polinómica en segundo plano, el hilo principal de FastAPI queda libre para seguir escuchando y respondiendo a nuevas peticiones concurrentes.

\subsubsection*{2. Políticas de Seguridad y Enrutamiento CORS}

Dado que la arquitectura distribuida tiene el cliente web y la API en dominios independientes, el navegador bloqueará por defecto cualquier comunicación cruzada. Para solventarlo, se ha implementado un \textit{middleware} CORS. El servidor aplica políticas de \textit{Zero Trust}, autorizando estrictamente peticiones procedentes del dominio del \textit{frontend}, limitando los métodos HTTP a \texttt{GET} y \texttt{POST}, y exigiendo la presencia explícita de cabeceras de autorización.

\subsubsection*{3. Inyección de Dependencias para la Identidad}

Para materializar el flujo definido en la Sección \ref{subsec:jwt_identidad}, se aprovecha el sistema nativo de Inyección de Dependencias de FastAPI. Se ha programado un interceptor que protege las rutas privadas:
    
\begin{itemize}
    \item Intercepta la petición HTTP antes de que alcance el controlador lógico de la ruta.
    \item Extrae el \textit{token} JWT y valida matemáticamente su firma utilizando la clave secreta del entorno (\texttt{JWT\_SECRET}).
    \item Si el \textit{token} es inválido o ha expirado, aborta la operación con un código \texttt{401 Unauthorized}. Si es válido, inyecta el \texttt{user\_id} verificado directamente en el contexto de la petición para uso del motor criptográfico.
\end{itemize}


\subsection{Implementación de los Endpoints Core}
\label{subsec:endpoints_core}

El control de la lógica de la plataforma se gestionan a través de tres \textit{endpoints} fundamentales en la API de FastAPI. Cada uno de ellos ha sido diseñado respetando los principios de mínimo privilegio y separación de responsabilidades.

\subsubsection{1. Generación de Identidades: \texttt{POST /api/generate}}

Este \textit{endpoint} actúa como el puente entre la autenticación web y el motor ML-DSA. Su función no es solo generar el par de claves post-cuánticas, sino vincularlas de forma segura a la sesión del usuario.

El flujo de ejecución implementa una técnica crítica de \textbf{Defensa en Profundidad}:
\begin{enumerate}
    \item El \textit{middleware} de FastAPI valida el \textit{token} JWT entrante y extrae el \texttt{user\_id}.
    \item Se invoca la función \texttt{keygen()} del motor criptográfico, generando la clave pública \texttt{pk} y secreta \texttt{sk}.
    \item \textbf{Delegación de Identidad y Contexto RLS:} En los diseños tradicionales, el servidor suele acceder a la base de datos utilizando una clave de administrador con privilegios absolutos, lo que anula cualquier política de seguridad a nivel de tabla. Para evitar este riesgo, la API implementa un modelo de delegación: el servidor crea una conexión temporal a Supabase pasándole exactamente el mismo \textit{token} JWT que envió el cliente. De este modo, la base de datos asume la identidad del usuario real de forma transparente, evaluando y aplicando rigurosamente sus reglas de Seguridad a Nivel de Fila RLS antes de permitir la inserción de las claves.
\end{enumerate}

\subsubsection{2. Arquitectura de Firma: \texttt{POST /api/sign}}

Para optimizar la privacidad y el rendimiento del sistema, no se utiliza un modelo de firma centralizado donde el servidor procesa el documento completo. En su lugar, se implementa un flujo de \textbf{Firma por Delegación de Hash}, siguiendo los estándares de seguridad modernos:

\begin{enumerate}
    \item \textbf{Pre-procesamiento en Cliente:} El \textit{frontend} estampa visualmente la firma en el PDF y, calcula de forma local el resumen criptográfico \textit{hash} SHA-256 del archivo modificado.
    \item \textbf{Tránsito Eficiente:} El cliente envía al \textit{backend} exclusivamente el \textit{hash} resultante. Esto garantiza que el contenido del documento nunca abandone el entorno local del usuario, cumpliendo con principios de privacidad por diseño.
    \item \textbf{Firma en Backend:} El servidor recibe el \textit{hash}, recupera la clave privada del usuario mediante el flujo de identidad delegado Sección \ref{subsec:endpoints_core}.1 y ejecuta el algoritmo \texttt{sign} de ML-DSA sobre dicha huella digital.
    \item \textbf{Sujeción de la Firma:} El servidor devuelve la firma post-cuántica al cliente, quien se encarga de inyectar este sello criptográfico en la estructura del PDF original para su descarga final.
\end{enumerate}

Esta arquitectura garantiza que el servidor actúe únicamente como un módulo hardware de seguridad lógico, comprendiendo la matemática del estándar FIPS 204 sin comprometer la confidencialidad de los documentos del usuario.

\subsubsection{3. Verificación Pública y Mitigación DoS: \texttt{POST /api/verify}}

A diferencia de los procesos anteriores, la verificación de una firma asimétrica debe ser universalmente accesible y no requerir autenticación. Este \textit{endpoint} recibe tres parámetros: el documento original, el archivo de firma (\texttt{.sig}) y la clave pública del autor.

Dado que la función \texttt{hash\_verify} requiere reconstruir la matriz polinómica y aplicar transformadas matemáticas, este \textit{endpoint} es altamente susceptible a ataques de Denegación de Servicio DoS mediante fuerza bruta computacional. Para mitigarlo, se ha implementado un patrón de \textbf{Rate Limiting} (\texttt{@limiter.limit("10/minute")}), restringiendo el número de validaciones que una misma dirección IP puede solicitar, protegiendo así los recursos asíncronos del servidor.



\subsection{Validación de Endpoints}
\label{subsec:mvp_conexion}

Tras el desarrollo de la lógica en el \textit{backend}, se implementó una interfaz de usuario provisional con el objetivo de validar el comportamiento de los \textit{endpoints} en un entorno real de cliente-servidor. Esta fase fue crítica para asegurar que el contrato de la API, entradas y salidas de datos, era correcto antes de proceder al diseño de la interfaz definitiva.

La validación se realizó siguiendo un orden lógico de dependencia funcional:

\begin{enumerate}
    \item \textbf{Pruebas de Persistencia \texttt{/generate}:} Se utilizó el cliente provisional para disparar la generación de claves. Se verificó visualmente que, tras la llamada, los registros aparecieran correctamente en las tablas de Supabase, confirmando que la delegación del \textit{token} JWT y las políticas RLS funcionaban según lo previsto.
    \item \textbf{Validación del Flujo de Firma (\texttt{/sign}):} Una vez confirmada la existencia de claves, se probó el envío de \textit{hashes} de documentos. El objetivo aquí fue validar que el servidor recuperara la clave privada correcta y devolviera un objeto de firma compatible con el estándar ML-DSA.
    \item \textbf{Cierre del Ciclo (\texttt{/verify}):} Finalmente, se integraron ambas partes para comprobar que una firma generada por el sistema podía ser validada con éxito mediante el \textit{endpoint} público, asegurando la integridad total del proceso criptográfico.
\end{enumerate}

Este enfoque permitió detectar y corregir fallos, evitando problemas de gestión en el servidor y estableciendo una base técnica estable para la Fase II del proyecto. Sin esta fase previa de comprobación, la creación de la interfaz de usuario habría sido más compleja, ya que podrían haberse propagado errores ocultos.


\section{Arquitectura del Cliente Web e Interfaz de Usuario}
\label{sec:arquitectura_cliente_web}

\subsection{Enrutamiento, Control de Acceso y Esqueleto de la Aplicación}
\label{subsec:enrutamiento_acceso}

La arquitectura del cliente web se ha diseñado bajo la modalidad de aplicación de página única (SPA), haciendo uso de la librería React Router. La navegación se gestiona íntegramente en el navegador mediante el componente central \texttt{<BrowserRouter>}. Esto elimina el coste de latencia por recarga de la página, permitiendo transiciones instantáneas.

El aspecto más crítico de esta capa arquitectónica es proteger las rutas privadas, como el \textit{Portal de Firmas} o el panel \textit{Mi Identidad}. Para aislar estas páginas, se ha desarrollado un guardián de rutas denominado \textit{Route Guard}, que se encuentra en el componente \texttt{<ProtectedRoute>}. Este componente evalúa el estado global de autenticación antes de permitir el montaje en memoria de la vista solicitada.

Si un usuario legítimo recarga la página forzosamente en su navegador, se debe validar su sesión y evitar que esta se cierre de forma inesperada. Para ello, el estado de la sesión se mantiene temporalmente en espera mientras el sistema consulta asíncronamente la validez del \textit{token} JWT en segundo plano. Si la validación del contexto falla, el guardián anula la petición de acceso y expulsa al usuario. Esta protección evita que un usuario malicioso pueda acceder a direcciones privadas, previniendo así posibles vectores de ataque como Open Redirect.

El ataque de \textit{Open Redirect} es una vulnerabilidad de seguridad web tipificada oficialmente como CWE-601 \cite{cwe601}. Este fallo de diseño ocurre cuando una aplicación acepta una entrada controlada por el usuario y la procesa para ejecutar una redirección HTTP hacia una URL arbitraria, sin someter dicha entrada a un proceso riguroso de validación o sanitización previa \cite{owasp_open_redirect}.

\subsection{Gestión de Estado Global y Autenticación}
\label{subsec:gestion_estado_autenticacion}

En una aplicación de página única SPA, la información de la sesión del usuario es un recurso que debe estar accesible desde multitud de componentes independientes, por ejemplo, la barra de navegación para mostrar el avatar, o el guardián de rutas para permitir el acceso. Si pasamos esta información manualmente de componente en componente, desde el nivel superior hasta el más profundo de la interfaz, caemos en un antipatrón de diseño conocido como \textit{prop drilling} o perforación de propiedades. Este fenómeno obliga a los componentes intermedios a recibir y transportar datos que realmente no necesitan, ensuciando el código y volviendo la arquitectura frágil ante futuros cambios \cite{react_prop_drilling}.

Para evitar esta cadena de transmisión ineficiente, se ha implementado un modelo de estado global centralizado utilizando la \textit{Context API} nativa de React combinada con el SDK de Supabase \cite{supabase_auth_reference}. Esta solución crea un canal de comunicación directo: cualquier componente de la interfaz puede solicitar y consumir los datos de la sesión sin alterar ni involucrar a sus componentes padre.

Este sistema se fundamenta en un patrón de diseño Proveedor-Consumidor. Se define un componente de nivel superior denominado \texttt{AuthProvider}. Este proveedor encapsula y expone cuatro variables de estado fundamentales para el ciclo de vida de la aplicación:

\begin{itemize}
    \item \texttt{user}: Objeto que contiene los metadatos de identidad del usuario extraídos del \textit{token}.
    \item \texttt{session}: El objeto de sesión activo, el cual encapsula los \textit{Access Tokens} y \textit{Refresh Tokens} en formato JWT \cite{rfc7519}.
    \item \texttt{isAuthenticated}: Un indicador booleano derivado que determina la presencia de una sesión válida.
    \item \texttt{isLoading}: Un semáforo de sincronización asíncrona que bloquea el renderizado de la interfaz mientras se recupera o verifica el estado de la sesión local contra el servidor en el arranque de la aplicación.
\end{itemize}

El principal desafío de gestionar la sesión en el navegador es mantenerla sincronizada con el servidor real. Por ejemplo, es necesario detectar instantáneamente si el \textit{token} JWT ha caducado o si el usuario ha cerrado sesión. Para resolver esto, el proveedor inicializa un oyente de eventos (\textit{Event Listener}) utilizando el método \texttt{onAuthStateChange} del SDK de Supabase, el cual avisa a la aplicación web de cualquier cambio en la cuenta.

La mecánica de esta sincronización sigue tres pasos lógicos:

\begin{enumerate}
    \item \textbf{Carga Inicial:} Al abrir la página web, el componente \texttt{AuthProvider} busca en la memoria del navegador (conocida como \textit{Local Storage}) si ya existe una sesión guardada de una visita anterior. Mientras dura esta rápida comprobación, el semáforo \texttt{isLoading} se mantiene en \texttt{true} para evitar que el sistema se equivoque y expulse al usuario a la pantalla de \textit{login} antes de tiempo.
    
    \item \textbf{Escucha Continua:} Tras la carga inicial, el sistema se queda escuchando en segundo plano. Si ocurre algún cambio en la sesión (el usuario inicia sesión, hace \textit{logout}, o el \textit{token} caduca por tiempo), Supabase dispara un aviso interno. Al recibirlo, la aplicación actualiza instantáneamente sus variables de memoria (\texttt{user} y \texttt{session}).
    
    \item \textbf{Actualización Automática:} En cuanto el \texttt{AuthProvider} actualiza los datos de la sesión, React detecta el cambio. En lugar de recargar toda la página web, React avisa de forma automática y exclusiva a los componentes que están usando el \textit{hook} \texttt{useAuth}. Así, si el usuario cierra sesión, el guardián de rutas lo expulsa y el menú superior cambia el avatar por el botón de \textit{login} de forma transparente.
\end{enumerate}

Gracias a esta arquitectura, la aplicación no tiene que programar manualmente la comprobación continua de la sesión o la renovación de los \textit{tokens}; ese trabajo pesado se delega en el SDK de Supabase. La interfaz web simplemente reacciona a los eventos, garantizando que lo que ve el usuario coincida en todo momento con sus permisos reales.

\subsection{Entorno Público y Herramientas \textit{Zero-Trust}}
\label{subsec:entorno_publico}

El entorno público de la aplicación actúa como la capa de presentación y validación del sistema. Está compuesto por una serie de vistas accesibles sin necesidad de autenticación, diseñadas para informar al visitante y, fundamentalmente, para permitir la auditoría de los documentos. Este entorno incluye rutas estáticas como la página de Inicio o la denominada \textit{Landing Page}, la documentación técnica sobre algoritmos poscuánticos y los apartados de resolución de dudas FAQ. 

Sin embargo, la pieza central de esta zona pública es el \textbf{Centro de Verificación} (\texttt{/dashboard/verify}). La decisión arquitectónica de mantener esta herramienta abierta y sin barreras de registro responde directamente al principio de seguridad \textit{Zero-Trust} Confianza Cero.

En un sistema de firma digital funcional, el acto de firmar es privado y exclusivo del propietario de la clave, pero el acto de verificar debe ser universal. Si un usuario de la plataforma firma un contrato y se lo envía a un tercero, por ejemplo, a un cliente o a un auditor, este tercero debe tener una vía rápida, transparente y gratuita para validar la autenticidad del documento. Exigir la creación de una cuenta de usuario únicamente para comprobar una firma introduciría una fricción inaceptable que arruinaría la usabilidad del sistema.

Cuando un visitante accede al Centro de Verificación, la interfaz le permite cargar un documento PDF. El sistema extrae el sobre criptográfico inyectado en el archivo y utiliza la clave pública correspondiente para validar matemáticamente la firma mediante ML-DSA. Al no requerir un \textit{token} de sesión JWT para ejecutar esta comprobación, el cliente web demuestra que la validez del documento reside en la integridad matemática de los datos, y no en los privilegios de acceso a la base de datos de la plataforma. Esta capacidad de firmar y de auditar es lo que otorga validez real a la herramienta en un entorno de producción.

\begin{figure}[H]
    \centering
    \includegraphics[width=0.75\textwidth]{images/verificacion.png}
    \caption{Interfaz pública de verificación de firmas ML-DSA.}
    \label{fig:verificacion}
\end{figure}


\subsection{El Dashboard Privado y el Núcleo Operativo}
\label{subsec:dashboard_privado}

Al superar la barrera de autenticación, el sistema muestra dos nuevas rutas. El usuario autenticado mantiene acceso total a las herramientas de la capa pública, como el Centro de Verificación o la documentación técnica, pero desbloquea la gestión de su identidad y la ejecución de la firma.\vspace{0.3cm}

\textbf{1. Panel "Mi Identidad": Centro de Auditoría Criptográfica}\vspace{0.3cm}

A nivel visual, este apartado (\texttt{/dashboard/identity}) adopta la usabilidad de un perfil de usuario moderno, priorizando la claridad informativa. Sin embargo, a nivel funcional, opera como un panel de auditoría personal donde el usuario puede consultar su clave pública y su actividad dentro del sistema.

La interfaz está configurada como una vista de solo lectura  que expone tres ejes de información fundamentales:
\begin{itemize}
    \item \textbf{Metadatos de Registro:} Refleja la información estática y de contacto vinculada a la cuenta durante la creación de la identidad.
    \item \textbf{Transparencia de Claves:} Muestra el material criptográfico público asociado al usuario, es decir su clave pública, permitiéndole auditar la credencial frente a terceros.
    \item \textbf{Historial de Operaciones:} Despliega un registro histórico de la actividad del usuario, permitiendo al sistema de trazabilidad sobre los eventos y firmas realizados.
\end{itemize}

Aunque la arquitectura de componentes de esta vista está diseñada para ser extensible y permitir la futura edición de perfiles o la rotación manual de claves, la implementación actual no permite el cambio de los datos. Esta decisión de diseño prioriza la estabilidad del sistema en su fase inicial.\vspace{0.3cm}

\textbf{2. Portal de Firmas: El Motor de Operaciones}\vspace{0.3cm}

Este es el núcleo de la aplicación (\texttt{/dashboard/sign}). Una vez que la identidad criptográfica está configurada, el usuario utiliza este portal para interactuar con los documentos. La interfaz guía al usuario y tiene el siguiente flujo:

\begin{enumerate}
    \item \textbf{Carga del Documento:} El usuario selecciona o arrastra un archivo PDF a la zona de carga temporal en el navegador.
    \item \textbf{Procesamiento Previo:} En lugar de subir el archivo completo al servidor (lo que consumiría un ancho de banda innecesario y plantearía vulnerabilidades de privacidad), el cliente web extrae la información y prepara el resumen \textit{hash} del documento de forma local.
    \item \textbf{Delegación de la Firma:} El cliente web se comunica de forma segura con el \textit{endpoint} \texttt{/sign} del \textit{backend}, solicitando que el documento sea sellado utilizando la clave privada que el usuario configuró previamente en su panel de identidad.
    \item \textbf{Retorno Visual:} La interfaz gestiona la espera asíncrona y, al recibir la firma poscuántica del servidor, confirma visualmente al usuario que la operación ha sido un éxito, ofreciéndole el documento final listo para descargar.
\end{enumerate}

Este diseño de flujo garantiza que el usuario no se abrume con configuraciones de seguridad complejas mientras realiza una tarea administrativa rutinaria, mejorando la experiencia de usuario (UX) y manteniendo el rigor técnico del sistema.

\begin{figure}[H]
    \centering
    \includegraphics[width=0.75\textwidth]{images/vista_firma.png}
    \caption{Panel de operación privada para la firma de documentos con ML-DSA.}
    \label{fig:vista_firma}
\end{figure}

\subsection{Interacción Asíncrona y \textit{Feedback} Visual (UI/UX)}
\label{subsec:interaccion_asincrona}

En una aplicación orientada a la seguridad, el diseño de la interfaz no es solo una cuestión estética, sino una herramienta fundamental para evitar errores y generar confianza. Las operaciones criptográficas poscuánticas, como generar claves o firmar un documento con ML-DSA, no son instantáneas. Requieren un tiempo de cómputo para resolver los cálculos algebraicos y para establecer la comunicación por red con el servidor.

Si durante esos segundos de espera la pantalla no hace nada, el usuario puede pensar que la aplicación se ha quedado bloqueada. La reacción natural suele ser pulsar el botón de "Firmar" varias veces seguidas o recargar la página de forma brusca, lo cual puede saturar el servidor o enviar órdenes duplicadas que corrompan el proceso de firma.

Para mitigar este problema, la interfaz web utiliza las librerías Tailwind CSS y Framer Motion para guiar al usuario y proteger el sistema, aplicando tres medidas de diseño básicas:

\begin{itemize}
    \item \textbf{Prevención de acciones duplicadas:} En el instante en que el usuario pulsa un botón crítico (como "Firmar Documento"), la aplicación desactiva temporalmente todos los botones interactivos de la pantalla. Visualmente se vuelven semitransparentes y el cursor cambia, bloqueando físicamente que el usuario pueda enviar la misma petición dos veces por impaciencia.
    
    \item \textbf{Indicadores visuales de carga:} En lugar de congelar la pantalla, se utilizan transiciones suaves para mostrar elementos de carga. Esto comunica al usuario que el sistema no ha fallado, sino que se encuentra procesando activamente la operación matemática solicitada en segundo plano.
    
    \item \textbf{Avisos claros en caso de error:} Si el servidor detecta un fallo criptográfico, como por ejemplo, un documento alterado o una firma no válida, la aplicación captura ese error y lo muestra en pantalla mediante pequeñas notificaciones emergentes. En lugar de romper la página, se informa al usuario de manera clara para que pueda tomar medidas.
\end{itemize}

Gracias a estas medidas, conseguimos que un proceso matemáticamente complejo y con tiempos de espera inevitables se perciba como una herramienta fluida y robusta, evitando que la confusión del usuario provoque fallos técnicos en el sistema.


\subsection{Motor Cliente de Procesamiento Documental}
\label{subsec:motor_cliente_procesamiento}

El procesamiento de los documentos constituye el nexo de unión entre la representación visual del documento PDF y el proceso criptográfico de la firma digital, que representa la matemática pura del proceso. Para que la herramienta sea verdaderamente privada y cumpla con el principio de \textit{Zero-Trust}, el cliente web debe extraer la esencia matemática del archivo directamente en el navegador del usuario, antes de que ninguna información viaje por la red.

Este proceso se divide en tres fases secuenciales: preparación visual, extracción de la huella matemática y la inyección de la firma mediante la técnica \textit{EOF Trick}.\vspace{0.3cm}

\textbf{1. Preparación Visual del Documento}\vspace{0.3cm}

Un documento firmado digitalmente requiere una marca visible para que cualquier auditor reconozca a simple vista que ha sido procesado. Para lograr esto, la aplicación utiliza la herramienta de manipulación \texttt{pdf-lib} en el propio navegador. 

Antes de realizar ninguna operación criptográfica, el cliente web dibuja un cajetín visual en la última página del PDF, inyectando textos informativos y un código QR generado al vuelo. Esta modificación es de carácter estético; no aporta seguridad matemática, pero es indispensable para la confianza del usuario final.\vspace{0.3cm}

\textbf{2. Extracción de la Huella y Privacidad por Diseño}\vspace{0.3cm}

En lugar de enviar el documento PDF completo al servidor (lo que consumiría un gran ancho de banda y plantearía un grave riesgo de privacidad si el documento contiene datos sensibles), la aplicación utiliza la API nativa de criptografía del navegador \texttt{Web Crypto API}.

El navegador lee el archivo PDF y calcula localmente su huella digital única utilizando el algoritmo de resumen SHA-256. Lo único que viaja a través de la red hacia el servidor \textit{backend} es esta pequeña cadena alfanumérica. El servidor recibe este \textit{hash}, lo introduce en el algoritmo poscuántico ML-DSA junto con la clave privada del usuario, y devuelve la firma matemática. El archivo PDF original nunca abandona el ordenador del usuario, garantizando una confidencialidad absoluta.\vspace{0.3cm}

\textbf{3. Inyección del Sobre Criptográfico \textit{EOF Trick}}\vspace{0.3cm}

El último paso del flujo resuelve el mayor reto del sistema: ¿cómo unimos la firma poscuántica devuelta por el servidor al archivo PDF sin corromper el formato del documento?

En lugar de crear páginas adicionales o alterar las complejas estructuras internas del PDF, el sistema implementa una técnica matemática directa y altamente eficiente conocida como \textit{EOF Trick} (Inyección al final del archivo, por las siglas de \textit{End Of File}). La mecánica es la siguiente:

\begin{itemize}
    \item El navegador recibe la firma y la empaqueta junto con el identificador del firmante en una estructura de texto, delimitada con las etiquetas personalizadas \texttt{<QPROOF>} y \texttt{</QPROOF>}.
    \item Esta estructura de texto se transforma en una secuencia de \textit{bytes}.
    \item Finalmente, el cliente web concatena estos datos matemáticos justo al final de la secuencia binaria del archivo PDF local.
\end{itemize}

La principal ventaja de este enfoque radica en su compatibilidad universal. Los visores de PDF tradicionales (como Adobe Acrobat) leen el archivo de principio a fin, y cuando encuentran la marca estándar de "fin de archivo", ignoran lo que hay detrás. Por tanto, el usuario puede abrir el documento con normalidad. Sin embargo, cuando este mismo documento se audita en nuestro Centro de Verificación, el algoritmo inspecciona directamente el final del archivo binario, extrae la etiqueta \texttt{<QPROOF>} y valida matemáticamente la autenticidad sin depender de servidores externos.

\section{Seguridad del Cliente Web}
\label{sec:bastionado_seguridad}

En un ecosistema donde la firma digital autoriza a un tercero para que gestione y firme documentos en nombre del titular original, el cliente web es la principal superficie de ataque. Por muy robustos que sean los algoritmos poscuánticos en el \textit{backend}, si la interfaz de usuario es vulnerable, el sistema en su conjunto queda comprometido. Para garantizar la integridad del entorno, se han implementado defensas perimetrales en la aplicación React, prestando especial atención a los vectores de ataque sobre la autenticación \cite{owasp_auth}.\vspace{0.3cm}

\textbf{1. Prevención de Inyecciones XSS}

El ataque de \textit{Cross-Site Scripting} XSS ocurre cuando un actor malicioso logra inyectar código ejecutable en la página web. Para mitigar este ataque, la arquitectura se apoya en la naturaleza declarativa de React. Todo el contenido dinámico, como los nombres de usuario o datos extraídos de los documentos, es analizado automáticamente antes de ser inyectado en el DOM. React convierte cualquier cadena de texto en contenido seguro, neutralizando etiquetas HTML o \textit{scripts} maliciosos \cite{owasp_xss}.\vspace{0.3cm}

\textbf{2. Flujo de Autenticación}\vspace{0.3cm}

El portal de acceso e identidad (\texttt{Auth.tsx}) es el perímetro más crítico. Se ha diseñado bajo un enfoque de defensa, aplicando las siguientes lógicas algorítmicas para repeler ataques automatizados:

\begin{itemize}
    \item \textbf{Limitación de Tasa (\textit{Rate Limiting}) y Anti-Fuerza Bruta:} Se ha programado un escudo temporal en el cliente que monitoriza los intentos de acceso. Si se detectan 5 intentos fallidos en menos de un minuto, el algoritmo bloquea físicamente la interacción con el formulario. Esto neutraliza de raíz los ataques de diccionario o la saturación por \textit{bots}.
    
    \item \textbf{Anti-Enumeración de Usuarios:} Un error clásico de diseño es devolver mensajes específicos como ``El usuario no existe'' o ``Contraseña incorrecta''. Esto permite a un atacante mapear qué correos están registrados en la plataforma. Para evitar esta fuga de información, el sistema responde siempre con un error genérico, como por ejemplo, ``Credenciales inválidas''. Ante los ojos de un atacante, la respuesta es indistinguible en ambos casos.
    
    \item \textbf{Análisis de las entradas:} Antes de transmitir cualquier petición de registro al servidor, una capa de validación de esquemas estricta recorta espacios en blanco innecesarios, formateando posibles caracteres maliciosos, valida la estructura del correo y restringe la longitud de los campos. Esto previene ataques de desbordamiento en la base de datos \cite{owasp_input_validation}.
    
    \item \textbf{Delegación Criptográfica y Doble Activación:} La aplicación web nunca manipula ni transmite las contraseñas en texto plano hacia el servidor poscuántico. El proceso se delega al SDK de Supabase, garantizando que el material sensible sea procesado mediante algoritmos de \textit{hashing} robustos (como bcrypt/Argon2) \cite{owasp_password_storage}. Además, las cuentas nacen bloqueadas por defecto hasta que el usuario resuelve un flujo \textit{Zero-Trust} haciendo clic en el enlace de confirmación enviado a su correo electrónico real.
\end{itemize}

\textbf{3. Políticas de Origen Cruzado CORS}\vspace{0.3cm}

Dado que el cliente web y el servidor operan en dominios diferentes, se utilizan Políticas de Origen Cruzado CORS. El servidor de la API se ha configurado con una lista blanca (\textit{whitelist}) estricta \cite{mdn_cors}. Únicamente se aceptan peticiones HTTP que provengan exactamente de la URL donde está alojado el cliente web oficial. Cualquier intento de conexión desde una web de terceros es rechazado a nivel de red, blindando el \textit{backend} contra accesos no autorizados.


\chapter{Conclusiones y Trabajo Futuro}
\label{ch:conclusiones}

\section{Conclusiones Generales}
\label{sec:conclusiones_generales}

El desarrollo de este Trabajo de Fin de Grado ha supuesto un desafío arquitectónico que ha obligado a unir dos disciplinas que habitualmente operan por separado: el rigor algebraico de la criptografía teórica y la usabilidad de la ingeniería de \textit{software} web. 

La amenaza de la computación cuántica sobre los esquemas asimétricos clásicos, como RSA \cite{rivest1978} o ECC \cite{miller1985}, ya no es una hipótesis lejana, sino un riesgo sistémico que ha motivado la estandarización urgente de algoritmos como ML-DSA \cite{nist_fips204}. La principal conclusión de este proyecto es que la elevada complejidad matemática de la criptografía poscuántica no tiene por qué traducirse en una mala usabilidad para el usuario final.

El resultado es un sistema distribuido plenamente funcional que emplea operaciones computacionalmente pesadas en una interfaz de usuario fluida, garantizando la integridad de documentos PDF con firmas poscuánticas sin corromper la estructura estándar del archivo.

\section{Grado de Consecución de los Objetivos}
\label{sec:grado_consecucion}

A partir de la planificación y los hitos técnicos definidos en el primer capítulo de esta memoria, se establece la siguiente evaluación crítica sobre el cumplimiento de los objetivos planteados:

\begin{itemize}
    \item \textbf{Asimilación de fundamentos algebraicos:} Se ha superado la elevada barrera de entrada matemática que exige la criptografía poscuántica. Se ha comprendido y documentado la viabilidad de los problemas de retículos (MLWE y MSIS), sentando la base teórica necesaria para justificar la seguridad del sistema frente a ataques con algoritmos cuánticos como el de Shor \cite{shor1994}.
    
    \item \textbf{Análisis crítico del estándar FIPS 204:} Se ha desglosado la complejidad de la especificación técnica publicada por el NIST \cite{nist_fips204}. Esto ha permitido comprender no solo el paradigma de firma de \textit{Fiat-Shamir con abortos} \cite{lyubashevsky2009fiat}, sino también los parámetros de seguridad que equilibran el tamaño de las claves con la velocidad de ejecución.
    
    \item \textbf{Desarrollo del motor criptográfico:} Se ha construido con éxito un núcleo algorítmico en Python que ejecuta la generación de claves, la firma y la verificación matemática de manera fiel a la especificación oficial. La estabilidad matemática de las transformadas NTT y la gestión de la entropía han sido validadas mediante técnicas de \textit{Test-Driven Development} TDD.
    
    \item \textbf{Diseño e integración del sistema distribuido:} El objetivo de llevar la matemática a un entorno real se ha llevado a la práctica mediante la construcción de una arquitectura cliente-servidor. Se ha desplegado una API REST asíncrona conectada a una aplicación de página única SPA protegida, resolviendo los problemas de delegación de firmas y manipulación binaria de archivos mediante la técnica del \textit{EOF Trick}.
    
    \item \textbf{Validación de rendimiento en hardware clásico:} De manera empírica, el funcionamiento de la aplicación web demuestra que el esquema ML-DSA es operativo sobre procesadores clásicos de arquitectura Von Neumann. Los tiempos de respuesta observados durante la serialización, comunicación por red y firma matemática se mantienen en rangos aceptables para una aplicación web moderna, descartando la necesidad de \textit{hardware} especializado para la adopción temprana de la tecnología poscuántica.
\end{itemize}

\section{Líneas de Trabajo Futuro}
\label{sec:trabajo_futuro}

Aunque la arquitectura actual cumple con los requisitos operativos propuestos y demuestra la viabilidad técnica de ML-DSA en entornos web, el ciclo de vida del \textit{software} permite identificar vías de mejora técnica para escalar el sistema hacia un entorno de producción masiva:

\begin{itemize}
    \item \textbf{Despliegue de Infraestructura y Contenedorización:} El motor criptográfico en Python requiere de una infraestructura dedicada para su alojamiento en producción. Se plantea la contenedorización del \textit{backend} mediante Docker y su despliegue en un Servidor Privado Virtual VPS. Esto mitigará los problemas de latencia, bloqueos de red y limitaciones de tiempo de ejecución \textit{timeouts} propios de las arquitecturas sin servidor \textit{serverless} gratuitas, permitiendo un rendimiento sostenido bajo concurrencia.
    
    \item \textbf{Integración con Infraestructura de Clave Pública PKI:} Actualmente, el sistema confía en las claves públicas expuestas en el panel de identidad del usuario. El siguiente paso evolutivo es integrar el estándar X.509 adaptado para ML-DSA \cite{ietf_mldsa_pki}, permitiendo la emisión de certificados digitales autorizados por una Autoridad de Certificación.

    \item \textbf{Transición hacia Tokens de Autenticación Poscuánticos:} El sistema actual utiliza JWT con el algoritmo clásico RS256 para la gestión de sesiones. Esto genera una asimetría de seguridad, constituyendo el eslabón criptográfico más débil ante capacidades cuánticas. Una prioridad futura es la migración de la capa de autenticación hacia tokens estructurados con algoritmos PQC. Aunque existen implementaciones experimentales y pruebas de concepto en la comunidad de código abierto, la estandarización oficial para integrar algoritmos como ML-DSA  se encuentra aún en fase de borrador en el IETF \cite{ietf_pqc_jose}. Su futura adopción garantizará una resistencia cuántica en toda la arquitectura del sistema.
    
\end{itemize}

\section{Impacto Tecnológico y Social}
\label{sec:impacto}


La llegada del denominado \textit{Q-Day} (el hito tecnológico en el que un ordenador cuántico alcance la capacidad de ejecutar el algoritmo de Shor de manera eficiente) amenaza con colapsar la confianza digital, invalidando la autenticidad de contratos, historiales médicos y transacciones financieras a nivel global. Lejos de ser una hipótesis a largo plazo, la trayectoria de investigación reciente muestra una reducción exponencial en los requisitos de hardware necesarios para romper la criptografía actual. Mientras que en 2012 se estimaba que quebrar RSA-2048 requeriría cerca de mil millones de qubits físicos, esta cifra descendió a 20 millones en 2019, a menos de un millón en 2025 \cite{google2025quantum}, y a menos de 100.000 en 2026 \cite{tqi2026qday}. Esta reducción constante, impulsada por mejoras algorítmicas, evidencia que la amenaza se acelera. A esto se suma el riesgo actual del \textit{Harvest Now, Decrypt Later} (HNDL), mediante el cual actores maliciosos ya están interceptando y almacenando tráfico cifrado con la intención de descifrarlo en cuanto el hardware cuántico sea viable \cite{ibm2025readiness}.

Ante este escenario, el impacto principal de este proyecto radica en su carácter anticipatorio y práctico. Al proporcionar una herramienta usable y accesible que implementa desde hoy el estándar NIST poscuántico (FIPS 204), se democratiza el acceso a una seguridad documental a prueba de futuro. El software desarrollado demuestra que las organizaciones no necesitan esperar a la llegada de la era cuántica, ni realizar inversiones masivas en hardware especializado, para empezar a proteger su tejido de confianza digital frente a la criptografía del mañana.

\chapter{Anexo: Interfaces y Flujo del Usuario}
\label{anexo:interfaces}

Este anexo documenta de forma secuencial el flujo de interacción de la plataforma, mostrando la usabilidad del sistema criptográfico y la robustez del entorno frente a diferentes estados de autenticación.

\section{Entorno Público y Divulgación Tecnológica}

\subsection{Landing Page y Especificaciones Técnicas}
La página de inicio actúa como punto de entrada e incluye una sección donde se detallan las especificaciones matemáticas del estándar FIPS 204.
\begin{figure}[H]
    \centering
    \includegraphics[width=1\textwidth]{images/landingpage.png}
    \caption{Página principal: Acceso directo a documentación técnica del algoritmo.}
\end{figure}

\begin{figure}[H]
    \centering
    \includegraphics[width=1\textwidth]{images/tecnologia.png}
    \caption{Página Técnica: Detalle de las especificaciones matemáticas del estándar FIPS 204}.
\end{figure}

\subsection{Centro de Verificación Pública}
Este módulo implementa el paradigma \textit{Zero-Trust}, permitiendo la validación de firmas sin necesidad de registro previo.
\begin{figure}[H]
    \centering
    \includegraphics[width=1\textwidth]{images/verifiacionpanle.png}
    \caption{Módulo de validación de firmas ML-DSA.}
\end{figure}

\section{Gestión de Identidad y Sesión}

\subsection{Registro y Autenticación}
El sistema emplea un protocolo de registro que vincula la identidad del usuario con la generación determinista de claves públicas y privadas según los parámetros de seguridad establecidos.

\begin{figure}[H]
    \centering
    \includegraphics[width=0.85\textwidth]{images/login.png}
    \caption{Interfaz de acceso seguro mediante credenciales (login).}
    \label{fig:login}
\end{figure}

\begin{figure}[H]
    \centering
    \includegraphics[width=0.5\textwidth]{images/resgitro.png}
    \caption{Formulario de registro y generación de par de claves para nuevos usuarios.}
    \label{fig:registro}
\end{figure}

\section{Operaciones Privadas}

\subsection{Panel de Firma y Generación de Evidencias}
Una vez autenticado, el usuario accede al \textit{dashboard} donde se ejecuta el motor de firma.
\begin{figure}[H]
    \centering
    \includegraphics[width=1\textwidth]{images/firmaotra.png}
    \caption{Dashboard privado: Ejecución de la operación de firma criptográfica.}
\end{figure}

\begin{figure}[H]
    \centering
    \includegraphics[width=1\textwidth]{images/evidencia.png}
    \caption{Interior del documento PDF.}
\end{figure}

\subsection{Verificación de Evidencias}
Finalmente, el sistema permite la descarga y posterior verificación del documento firmado, confirmando su integridad mediante la reconstrucción del compromiso.

\begin{figure}[H]
    \centering
    \includegraphics[width=1\textwidth]{images/VERIFICACION2.png}
    \caption{Interfaz general de validación.}
\end{figure}

En caso de que el documento haya sido alterado o la firma no coincida con el compromiso original, el sistema ejecuta el muestreo por rechazo y notifica la invalidez de la operación.

\begin{figure}[H]
    \centering
    \includegraphics[width=0.5\textwidth]{images/ERROR_VERIFICACION.png}
    \caption{Validación fallida.}
\end{figure}

\newpage
\addcontentsline{toc}{chapter}{Bibliografía}
\bibliographystyle{ieeetr}
\bibliography{bibliografia}

\end{document}